
%%FILE blueprint/src/chapter/AdeleMiniproject.tex
\chapter{Miniproject: Adeles}\label{Adele_miniproject}

\section{Status}

This is an active miniproject.

\section{The goal}

There are several goals to this miniproject.

\begin{enumerate}
  \item Define the adeles $\A_K$ of a number field~$K$ and
    give them the structure of a $K$-algebra (status: now in mathlib thanks to
    Salvatore Mercuri);
  \item Prove that $\A_K$ is a locally compact topological ring (status:
      \href{https://github.com/smmercuri/adele-ring_locally-compact}{
      also proved by Mercuri} but not yet in mathlib);
  \item Base change: show that if $L/K$ is a finite extension of number fields then the
    natural map $L\otimes_K\A_K\to\A_L$ is an isomorphism, both algebraic and
    topological; (status: not
    formalized yet, but there is a plan -- see the project dashboard);
  \item Prove that $K \subseteq \A_K$ is a discrete subgroup and the quotient
    is compact (status: not formalized yet, but there is a plan -- see the project
    dashboard);
  \item Get this stuff into mathlib (status: (1) done, (2)--(4) not done).
\end{enumerate}

We briefly go through the basic definitions. Let $K$ be a number field.
Let $\Zhat=\projlim_{N\geq1}(\Z/N\Z)$ be the profinite completion of $\Z$,
equipped with the projective limit topology.

A cheap definition of the finite adeles $\A_K^\infty$ of $K$ is $K\otimes_{\Z}\Zhat$,
equipped with the $\Zhat$-module topology.
A cheap definition of the infinite adeles
$K_\infty$ of $K$ is $K\otimes_{\Q}\R$ with the $\R$-module topology (this is a
finite-dimensional $\R$-vector space so this is just the usual topology on $\R^n$).
A cheap definition of the adeles of $K$ is $\A_K^\infty\times K_\infty$ with
the product topology. This is a commutative topological ring.

However in the literature (and in mathlib) we see different definitions.
The finite adeles of $K$ are usually defined in the books
as the so-called restricted product $\prod'_{\mathfrak{p}}K_{\mathfrak{p}}$ over the completions
$K_{\mathfrak{p}}$ of $K$ at all maximal ideals $\mathfrak{p}\subseteq\mathcal{O}_K$ of the
integers of $K$. Here the restricted product is the subset of $\prod_{\mathfrak{p}}K_{\mathfrak{p}}$
consisting of elements which are in the integers $\mathcal{O}_{K,\mathfrak{p}}$ of
$K_{\mathfrak{p}}$ for all but finitely many $\mathfrak{p}$. This is the definition given in
mathlib. Mathlib also has the proof that they're a topological ring;
furthermore the construction of the finite adeles in mathlib works for any
Dedekind domain (this was pointed out to me by Mar\'ia In\'es
de Frutos Fern\'andez; the adeles
are an arithmetic object, but the finite adeles are an algebraic object).

Similarly the infinite adeles of a number field~$K$
are usually defined as $\prod_v K_v$,
the product running over the archimedean completions of~$K$, and this is
the mathlib definition.

The adeles of a number field $K$ are the product of the finite and infinite
adeles, and mathlib knows that they're a $K$-algebra and a topological ring.

\section{Local compactness}

As mentioned above, Salvatore Mercuri was the first to give a complete formalisation of the proof
that the adele ring is locally compact as a topological space. His work is in
\href{https://github.com/smmercuri/adele-ring_locally-compact}{his own repo} and proved the
result using the ``ad hoc'' topology on the adeles which we initially had. Since then,
adeles have been refactored to have the direct limit topology and mathlib has
\href{https://leanprover-community.github.io/mathlib4_docs/Mathlib/Topology/Algebra/RestrictedProduct.html\#RestrictedProduct.locallyCompactSpace_of_addGroup}
{\tt RestrictedProduct.locallyCompactSpace\_of\_addGroup}, the result
that a restricted product of topological additive groups $K_v$ over compact open
subgroups $A_v$ is locally compact.

What we need then is this (note that this is not true for a general Dedekind domain):
\begin{theorem}
  \label{NumberField.instCompactSpaceAdicCompletionIntegers}
  \lean{NumberField.instCompactSpaceAdicCompletionIntegers}
  \discussion{451}
  \leanok
  If $K$ is a number field and $v$ is a nonzero prime ideal of the integers of $K$,
  then the integers of $K_v$ is a compact open subgroup.
\end{theorem}
\begin{proof} Openness should follow from the fact that the integers are
  $\{x : v(x)<v(1/\pi)\}$ where $\pi$ is a uniformizer. Compactness needs
  finiteness of the residue field $\mathcal{O}_K/v$.
\end{proof}

Once we have this, the above result from mathlib gives us
\begin{theorem}
  \lean{NumberField.AdeleRing.locallyCompactSpace}
  \label{NumberField.AdeleRing.locallyCompactSpace}
  \uses{NumberField.instCompactSpaceAdicCompletionIntegers}
  \discussion{253}
  \leanok
  The adeles of a number field are locally compact.
\end{theorem}
\begin{proof}
  The adeles of a number field are a product of the finite adeles and the infinite adeles
  so it suffices to prove that the finite and infinite adeles are locally compact.
  The infinite adeles are just isomorphic to $\R^n$ as a topological space, so they're certainly
  locally compact. As for the finite adeles,
  the mathlib theorem {\tt RestrictedProduct.locallyCompactSpace\_of\_addGroup}
  says that a restricted product of locally compact additive groups with respect to open compact
  subgroups is locally compact, so this reduces us the previous result.
\end{proof}

\section{Base change}

The ``theorem'' we want is that if $L/K$ is a finite extension of number fields,
then $\A_L=L\otimes_K\A_K$. This isn't a theorem though, this is actually a \emph{definition}
(the map between the two objects) and a theorem about
the definition (that it's an isomorphism). In fact the full claim is that it is both a homeomorphism
and an $L$-algebra isomorphism. Before we can prove the theorem, we need to make the
definition.

Recall that the adeles $\A_K$ of a number field is a product $\A_K^\infty\times K_\infty$
of the finite adeles and the infinite adeles. So our ``theorem'' follows immediately from
the ``theorem''s that $\A_L^\infty=L\otimes_K\A_K^\infty$ and $L_\infty=L\otimes_KK_\infty$
(both of these equalities mean an algebraic and topological isomorphism).
We may thus treat the finite and infinite results separately.

\subsection{Base change for nonarchimedean completions.}

As pointed out above, the theory of finite adeles works fine for Dedekind domains.
So for the time being let~$A$ be a Dedekind domain. Recall that the \emph{height one spectrum}
of $A$ is the nonzero prime ideals of~$A$. Note that because we stick to the literature,
rather than to common sense, fields are Dedekind domains in mathlib, and the
height one spectrum of a field is empty. The reason I don't like allowing fields
to be Dedekind domains is that geometrically the definition of Dedekind
domain in the literature is ``smooth affine curve, or a point''. But many theorems in algebraic
geometry begin ``let $C$ be a smooth curve'', rather than ``let $C$ be a smooth curve or a point''.

Let $K$ be the field of fractions of $A$. If $v$ is in the height one spectrum of $A$,
then we can put the $v$-adic topology on $A$ and on $K$, and consider the completions
$A_v$ and $K_v$. The finite adele ring $\mathbb{A}_{A,K}^\infty$ is defined to be
the restricted product of the $K_v$ with respect to the $A_v$, as $v$ runs over
the height one spectrum of $A$. It is topologised by making $\prod_v A_v$ open with
the product topology (here $A_v$ has the $v$-adic topology).

Now let~$L/K$ be a finite separable extension, and let $B$ be the integral closure of~$A$ in~$L$.
We want to relate the finite adeles of $K$ and of $L$. We work place by place, starting by fixing
one place $w$ of $B$ and analysing the relation of $L_w$ and $B_w$ to the completions $K_v$
and $A_v$ where $v$ is the place of $A$ dividing $w$.

So let $w$ be a nonzero prime ideal of $B$. Say $w$ lies over $v$, a prime ideal of $A$.
Then we can put the $w$-adic topology on $L$ and the $v$-adic topology on~$K$. Furthermore
we can equip $K$ with an additive $v$-adic valuation, that is,
a function also called $v$ from $K$ to $\Z\cup\{\infty\}$ normalised so that if $\pi$ is a uniformiser
for $v$ then $v(\pi)=1$. Similarly we consider $w$ as a function from $L$ to $\Z\cup\{\infty\}$.
The next lemma explains how these valuations are related.

\begin{lemma}
  \label{IsDedekindDomain.HeightOneSpectrum.valuation_comap}
  \lean{IsDedekindDomain.HeightOneSpectrum.valuation_comap}
  \leanok
  If $i:K\to L$ denotes the inclusion then for $k\in K$ we have
  $e\times w(i(k))=v(k)$, where $e$ is the ramification index of $w/v$
  (recall that valuations here are written additively, unlike in mathlib).
\end{lemma}
\begin{proof}
  \leanok
  Standard (and formalized).
\end{proof}

\begin{definition}
  \lean{IsDedekindDomain.HeightOneSpectrum.Extension.adicCompletionSemialgHom}
  \label{IsDedekindDomain.HeightOneSpectrum.Extension.adicCompletionSemialgHom}
  \uses{IsDedekindDomain.HeightOneSpectrum.valuation_comap}
  \leanok
  There's a natural ring map $K_v\to L_w$ extending the map $K\to L$.
  It is defined by completing
  the inclusion $K\to L$ at the finite places $v$ and $w$ (which can be done
  because the previous lemma shows that the map is uniformly continuous for the $v$-adic
  and $w$-adic topologies).
\end{definition}

\begin{lemma}
  \lean{IsDedekindDomain.HeightOneSpectrum.Extension.valued_adicCompletionSemialgHom}
  \label{IsDedekindDomain.HeightOneSpectrum.Extension.valued_adicCompletionSemialgHom}
  \uses{IsDedekindDomain.HeightOneSpectrum.valuation_comap,
  IsDedekindDomain.HeightOneSpectrum.Extension.adicCompletionSemialgHom}
  \leanok
  If $i_v:K_v\to L_w$ denotes the map of the previous definition
  then for $x\in K_v$ we have
  $e\times w(i(k))=v(k)$, where $e$ is the ramification index of $w/v$.
\end{lemma}
\begin{proof}
  \leanok
  Follows by continuity from lemma~\ref{IsDedekindDomain.HeightOneSpectrum.valuation_comap}.
\end{proof}

\begin{lemma}
  \lean{IsDedekindDomain.HeightOneSpectrum.Extension.adicCompletionSemialgHom_image_adicCompletionIntegers}
  \label{IsDedekindDomain.HeightOneSpectrum.Extension.adicCompletionSemialgHom_image_adicCompletionIntegers}
  \uses{IsDedekindDomain.HeightOneSpectrum.Extension.valued_adicCompletionSemialgHom}
  \leanok
  The map $i_v:K_v\to L_w$ sends the integer ring $A_v$ into $B_w$.
\end{lemma}
\begin{proof}
  \leanok
  The integer ring is defined by $v\geq0$ (or $v\leq 1$ in mathlib, which uses multiplicative
  valuations) so the result follows from \ref{IsDedekindDomain.HeightOneSpectrum.Extension.valued_adicCompletionSemialgHom}.
\end{proof}
\begin{theorem}
  \lean{IsDedekindDomain.HeightOneSpectrum.adicCompletion.instIsModuleTopology}
  \label{IsDedekindDomain.HeightOneSpectrum.adicCompletion.instIsModuleTopology}
  \uses{IsDedekindDomain.HeightOneSpectrum.Extension.adicCompletionSemialgHom}
  \leanok
  \discussion{326}
  Giving $L_w$ the $K_v$-algebra structure coming from the natural map $K_v\to L_w$,
  the $w$-adic topology on $L_w$ is the $K_v$-module topology.
\end{theorem}
\begin{proof}
  Any basis for $L$ as a $K$-vector space spans $L_w$ as a $K_v$-module, so $L_w$ is
  finite-dimensional over $K_v$ and the module topology is the same as the product
  topology. So we need to establish that the product topology on $L_w=K_v^n$ is
  the $w$-adic topology. But the $w$-adic topology is induced by the $w$-adic norm,
  which makes $L_w$ into a normed $K_v$-vector space, and (after picking a basis)
  the product norm on $L_w=K_v^n$ also makes $L_w$ into a normed $K_v$-vector space.
  So the result follows from the standard fact (see for example the lemma on p52
  of Cassels-Froelich, formalized as {\tt ContinuousLinearEquiv.ofFinrankEq} in mathlib)
  that any two norms on a finite-dimensional vector space over
  a complete field are equivalent (and thus induce the same topology).
\end{proof}

Because of the commutative diagram
\begin{center}
\begin{tikzcd}
K_v \arrow{r} & L_w  \\
K \arrow{u} \arrow{r} & L \arrow{u}
\end{tikzcd}
\end{center}

we can view $L_w$ as an $L\otimes_KK_v$-algebra.

Now instead of fixing $w$ upstairs, we fix $v$ downstairs and consider all $w$ lying
over it at once. So say $v$ is in the height one spectrum of $A$.

\begin{lemma}
  \lean{IsDedekindDomain.HeightOneSpectrum.Extension.finite}
  \label{IsDedekindDomain.HeightOneSpectrum.Extension.finite}
  \leanok
  There are only finitely many primes $w$ of $B$ lying above $v$.
\end{lemma}
\begin{proof}
  \leanok
  This is a standard fact about Dedekind domains. The key input is
  mathlib's theorem {\tt primesOver\_finite}.
\end{proof}

We write $w|v$ to denote the fact that $w$ is a prime of $B$ above $v$ of $A$.

\begin{definition}
  \lean{IsDedekindDomain.HeightOneSpectrum.adicCompletion.semialgHomPi}
  \label{IsDedekindDomain.HeightOneSpectrum.adicCompletion.semialgHomPi}
  \uses{IsDedekindDomain.HeightOneSpectrum.Extension.adicCompletionSemialgHom}
  \leanok
  The product of the maps $K_v\to L_w$ for $w|v$ is a natural ring map $K_v\to\prod_{w|v}L_w$
  lying over $K\to L$.
\end{definition}

Because $K_v\to\prod_{w|v}L_w$ lies over $K\to L$, there's an induced $L$-algebra
map $L\otimes_KK_v\to\prod_{w|v}L_w$. We are now able to state one of the key results
in this section. The proof is probably the hardest proof
in this section to formalize.

\begin{theorem}
  \lean{IsDedekindDomain.HeightOneSpectrum.adicCompletion.baseChangeAlgEquiv}
  \label{IsDedekindDomain.HeightOneSpectrum.adicCompletion.baseChangeAlgEquiv}
  \uses{IsDedekindDomain.HeightOneSpectrum.adicCompletion.semialgHomPi,
  IsDedekindDomain.HeightOneSpectrum.Extension.finite}
  \leanok
  The induced $L$-algebra homomorphism $L\otimes_KK_v\to\prod_{w|v}L_w$ is an
  isomorphism of rings.
\end{theorem}
\begin{proof}

  My current proposal to formalize this is as follows. The map is surjective
  because the image is dense and closed; this has been formalized already.
  It is also a $K_v$-algebra homomorphism if we give $L_w$ the obvious $K_v$-algebra
  structure. Thus we can conclude the result if we can prove that both spaces are
  finite-dimensional and have the same dimension. The $K_v$-dimension of $L\otimes_KK_v$
  is equal to the $K$-dimension of $L$, which is $\sum_{w|v}e_wf_w$ using the standard
  notation that $e_w$ is the ramification index of $w$ and $f_w$ the residue degree
  (this result is in mathlib). So it suffices to prove that $[L_w:K_v]=e_wf_w$.
  We already have that $e_w$ (defined globally) is equal to the local ramification
  index (defined as the factor by which the valuations differ on $K$). So what is left
  is to prove that (i) the residue field extension induced by $L_w/K_v$ has degree is equal to the
  globally-defined $f_w$, (ii) an extension of local fields has degree $ef$. Now (i) sounds
  straightforward given what we have (the map from $A$ to $\mathcal{O}_v$ has kernel $v$ and
  dense image) and (ii) is true for any complete discretely-valued field; I am not suggesting
  we formalize the following proof, but at least it represents a rigorous justification:
  A field complete with respect to a discrete valuation is \emph{stable} in the sense
  of the book by Bosch-G\"{u}ntzer-Remmert (Prop 3.6.2.1), so every finite extension of such a field
  is cartesian (def 3.6.1.1) and thus $ef=n$ (Prop 3.6.2.4, (iii) implies (ii)). Note
  that if you weaken the hypotheses too much then there are counterexamples; it's possible
  to have $ef<n$ and BGR goes into details.

\end{proof}

\begin{theorem}
  \label{IsDedekindDomain.HeightOneSpectrum.adicCompletion.instIsModuleTopologyPi}
  \lean{IsDedekindDomain.HeightOneSpectrum.adicCompletion.instIsModuleTopologyPi}
  \uses{IsDedekindDomain.HeightOneSpectrum.adicCompletion.instIsModuleTopology}
  \leanok
  For $v$ fixed, the product topology on $\prod_{w|v}L_w$ is the $K_v$-module
  topology.
\end{theorem}
\begin{proof}
  \leanok
  This is a finite product of $K_v$-modules each of which has the $K_v$-module topology
  by~\ref{IsDedekindDomain.HeightOneSpectrum.adicCompletion.instIsModuleTopology},
  and the product topology is the module topology for a finite product of modules each of which
  has the module topology (this is in mathlib).
\end{proof}

\begin{theorem}
  \label{IsDedekindDomain.HeightOneSpectrum.adicCompletion.baseChangeContinuousAlgEquiv}
  \lean{IsDedekindDomain.HeightOneSpectrum.adicCompletion.baseChangeContinuousAlgEquiv}
  \uses{IsDedekindDomain.HeightOneSpectrum.adicCompletion.instIsModuleTopologyPi}
  If we give $L\otimes_KK_v$ the $K_v$-module topology then the $L$-algebra isomorphism
  $L\otimes_K K_v\cong\prod_{w|v}L_w$ is also a homeomorphism.
\end{theorem}
\begin{proof} Indeed, is a $K_v$-algebra isomorphism between two modules each of which
  have the module topology, and any module map is automorphically continuous for the
  module topologies.
\end{proof}

We now start thinking about what's going on at the integral level. We write $A_v$
for the integers of $K_v$ and $B_w$ for the integers of $L_w$.

\begin{theorem}
  \lean{IsDedekindDomain.HeightOneSpectrum.range_baseChange_comp_tensorAdicCompletionTo_eq_pi}
  \label{IsDedekindDomain.HeightOneSpectrum.range_baseChange_comp_tensorAdicCompletionTo_eq_pi}
  \uses{IsDedekindDomain.HeightOneSpectrum.adicCompletion.baseChangeAlgEquiv}
  \leanok
  The isomorphism $L\otimes_KK_v\to\prod_{w|v}L_w$ induces an isomorphism
  $B\otimes_AA_v\to \prod_{w|v}B_w$
  for all $v$ in the height one spectrum of $A$.
\end{theorem}
\begin{proof}
  \leanok
  Certainly the image of the integral elements are integral. The argument in the other
  direction is more delicate. My original plan was to follow Cassels--Froehlich,
  Cassels' article ``Global fields'', section 12 lemma, p61, which proves it for
  all but finitely many primes, but \href{https://github.com/ImperialCollegeLondon/FLT/pull/400}
  {a PR by Matthew Jasper} gives another approach which works for all primes.
  Jasper's argument is to show that the closure of $A$ in $K_v$ is $A_v$
  for a valuation on a Dedekind domain, and then that the closure of $A$ in $\prod_{v\in S}K_v$
  is $\prod_{v\in S}A_v$ for $S$ a finite set of valuations (using the Chinese
  remainder theorem). Applying this to $B$ we get that the closure of $B$ in $\prod_{w|v}L_w$
  is $\prod_{w|v}B_w$. He then shows that this closure is the image of
  $B\otimes_A\mathcal{O}_v$ (by showing that this image is closed because it's open),
  giving surjectivity; injectivity follows from the statement
  that $L\otimes_KK_v=\prod_{w|v}L_w$.
\end{proof}

A summary of what we have so far: for all finite places $v$ of $A$
we have shown that the natural map $L\otimes_KK_v\to\prod_wL_w$
is an isomorphism of $L$-algebras, and that if $L\otimes_KK_v$ has
the $K_v$-module topology and each $L_w$ has the valuation topology
then this map is also a homeomorphism. Furthermore we have shown
that there is an induced algebraic isomorphism $B\otimes_AA_v\equiv\prod_w B_w$
on the subrings of the left and right hand sides.

Recall that the finite adeles $\A_{A,K}^\infty$ is defined in mathlib to be
the restricted product of the $K_v$ with respect to the $A_v$, equipped with a certain
restricted product topology (which is not the subspace topology of the product
topology, indeed $\prod_v A_v$ is open in this topology). We have seen in
definition~\ref{IsDedekindDomain.HeightOneSpectrum.Extension.adicCompletionSemialgHom} that
there's a map $K_v\to L_w$ if $w|v$, extending $K\to L$, and we have seen in
theorem~\ref{IsDedekindDomain.HeightOneSpectrum.Extension.adicCompletionSemialgHom_image_adicCompletionIntegers}
that this sends $A_v$ to $B_w$. We conclude

\begin{definition}
  \label{IsDedekindDomain.FiniteAdeleRing.mapSemialgHom}
  \lean{IsDedekindDomain.FiniteAdeleRing.mapSemialgHom}
  \uses{IsDedekindDomain.HeightOneSpectrum.Extension.adicCompletionSemialgHom,
  IsDedekindDomain.HeightOneSpectrum.Extension.adicCompletionSemialgHom_image_adicCompletionIntegers}
  \leanok
  There's a natural ring homomorphism
  $\A_{A,K}^\infty\to\A_{B,L}^\infty$ lying over $K\to L$.
\end{definition}

Hence there's a natural $L$-algebra homomorphism $L\otimes_K\A_{A,K}^\infty\to\A_{B,L}^\infty$.

Our next goal in this section is the following two results. First the algebraic claim:

\begin{theorem}
  \label{IsDedekindDomain.FiniteAdeleRing.baseChangeAlgEquiv}
  \lean{IsDedekindDomain.FiniteAdeleRing.baseChangeAlgEquiv}
  \uses{IsDedekindDomain.HeightOneSpectrum.range_baseChange_comp_tensorAdicCompletionTo_eq_pi,
    RestrictedProduct.relabelIso,
    IsDedekindDomain.FiniteAdeleRing.baseChangeIntegralAlgEquiv,
    IsDedekindDomain.AKLB.tensorProduct_module_algEquiv}
  \leanok
  \discussion{243}
  This natural map $L\otimes_K\A_{A,K}^\infty\to\A_{B,L}^\infty$ is an isomorphism.
\end{theorem}

Now $L\otimes_K\A_{A,K}^\infty$ is an $\A_{A,K}^\infty$-module and hence can be given
the $\A_{A,K}^\infty$-module topology. We also claim

\begin{theorem}
  \label{IsDedekindDomain.FiniteAdeleRing.baseChangeContinuousAlgEquiv}
  \lean{IsDedekindDomain.FiniteAdeleRing.baseChangeContinuousAlgEquiv}
  \uses{IsDedekindDomain.HeightOneSpectrum.range_baseChange_comp_tensorAdicCompletionTo_eq_pi}
  \leanok
The induced $L$-algebra morphism
  $L\otimes_K\A_{A,K}^\infty\to\A_{B,L}^\infty$ is a topological isomorphism.
\end{theorem}

Informally, the proofs are simple: componentwise we know
that $L\otimes_KK_v$ is isomorphic both algebraically and
topologically to $\prod_{w|v}L_w$, and that this isomorphism
sends the open set $B\otimes_AA_v$ homeomorphically onto
$\prod_{w|v}B_w$, so now it's ``just a case of putting everything
together''. Formally, we really need to spell this out, as there is a lot
going on. We do this in the
next subsection.

\subsection{Base change for nonarchimedean completions.}

As usual we are in the AKLB set-up, so in particular $K$ is the field
of fractions of the Dedekind domain $A$, $L/K$ is a finite
separable extension, and $B$ is the integral closure of $A$
in $L$. The goal in this subsection is to spell out the following argument: Assume that
$L\otimes_KK_v\cong\prod_{w|v}L_w$
algebraically and topologically for all $v$, with $B\otimes_AA_v$ identified with $\prod_{w|v}B_w$.
Then $L\otimes_K\A_K^\infty\cong\A_L^\infty$, algebraically and topologically. Here
the tensor products $L\otimes_K R$ (for $R$ a $K$-algebra with a topology) are all being
given the $R$-module topology, which if we choose a basis for $L/K$ is just the product
topology.

We start with the following observation. If $M$ is a $K$-module then there's a canonical
map $B\otimes_AM\to L\otimes_KM$ sending $b\otimes m$ to $b\otimes m$ (this follows from the
universal property of the tensor product). Our first goal
is to show that this map is an isomorphism. Let us establish some lemmas first.

\begin{lemma}
  \label{IsDedekindDomain.dvd_norm}
  If $0\not=b\in B$ then there exists $0\not=a\in A$ such that $b$ divides
  the image of $a$ in $B$.
\end{lemma}
\begin{remark} Is this already in mathlib?
\end{remark}
\begin{proof} Let $a=N_{L/K}(b)$, the norm. This is known to take nonzero elements of $L$
to nonzero elements of $K$ (because the norm is the determinant of an invertible linear map)
and integral elements to integral elements. Furthermore $a/b\in L$ is the the product of the
conjugates of $b$ in some normal closure of $L$, and hence it is integral, and thus in $B$.
\end{proof}

\begin{corollary}
  \label{IsDedekindDomain.AKLB.surjective_tensorProduct_map}
  \uses{IsDedekindDomain.dvd_norm}
  The $A$-bilinear map $B\times K\to L$ sending $(b,k)$ to $bk$ is surjective.
\end{corollary}
\begin{proof} Given $\lambda\in L$ write it as $n/d$ with $0\not=d\in B$. Choose $0\not=a\in A$
  and $b\in B$ with $db=a$ and then note $\lambda=nb/a=nb\times a^{-1}$.
\end{proof}

\begin{corollary}
  \label{IsDedekindDomain.FiniteAdeleRing.tensorProduct_algEquiv}
  \uses{IsDedekindDomain.dvd_norm}
  The natural map $B\otimes_AK\to L$ is a $B$-algebra isomorphism.
\end{corollary}
\begin{proof}

We write down an inverse. Regard $B\otimes_AK$ as a $B$-algebra via the action on the left.
Note that at this point it's not even clear that $B\otimes_AK$ is a field. We have the
structure map $B\to B\otimes_AK$ sending $b$ to $b\otimes1$, which is $B$-linear. I claim
that every nonzero element of $B$ gets sent to an invertible element of $B\otimes_AK$.
Indeed, if $b\not=0$ and (using the previous lemma) we choose $0\not=a\in A$ such that
$bb'=a$, then $(b\otimes1)(b'\otimes\frac1a)=1$. Thus by the universal property of
localisation, the $B$-linear map $B\to B\otimes_AK$ extends to a ring homomorphism from the field
of fractions of $B$ to $B\otimes_AK$, which we claim is our desired inverse.
Checking that both composites are the identity should be straightforward. Starting
with $B\otimes_AK$ we only have to check on elements of the form $b\otimes k$;
starting with $L$ we only have to check on elements of $B$. Hopefully both are
straightforward.
\end{proof}

\begin{corollary}
  \label{IsDedekindDomain.AKLB.tensorProduct_module_algEquiv}
  %%% \uses{IsDedekindDomain.FiniteAdeleRing.baseChangeAlgEquiv} -- uncomment to break leanblueprint
  \lean{IsDedekindDomain.linearEquivTensorProductModule}
  \leanok
  If $M$ is any $K$-module then the canonical map $B\otimes_A M\to L\otimes_K M$
  is an isomorphism.
\end{corollary}
\begin{proof}
  \leanok
  We can factor this map as $B\otimes_AM\cong B\otimes_A(K\otimes_KM)\cong
  (B\otimes_A K)\cong_KM\to L\otimes_KM$ and we just showed that the latter map was an
  isomorphism.
\end{proof}

We now need to explain how tensor products sometimes commute with restricted products.
Something we will need along the way is

\begin{theorem} $B$ is a finitely-presented $A$-module.
  \label{IsDedekindDomain.AKLB.finitePresentation}
\end{theorem}
\begin{proof} $A$ is Noetherian as it is a Dedekind domain, so it suffices to prove that $B$ is
  finitely-generated as an $A$-module. But this is in mathlib already (a proof is
  around line 101 of {\tt BaseChange.lean} in FLT at the time of writing).
\end{proof}

The reason we care about this is the following.

\begin{theorem}
  \label{pi_tensorProduct_of_finitePresentation}
  If $R$ is a commutative ring, if $M$ is a finitely presented $R$-module
  and if $N_i$ are a collection of $R$-modules, then the canonical map
  $M\otimes_R\prod_i N_i\to\prod_i(M\otimes_R N_i)$ is an isomorphism.
\end{theorem}
\begin{proof} If $M$ is finite and free then Maddy Crim has already formalized this
  in FLT. For the general case present $M$ as $R^a\to R^b\to M\to 0$ and use that tensor
  products and arbitrary products preserve surjections.
\end{proof}

\begin{corollary}
  \label{IsDedekindDomain.pi_tensorProduct}
  \uses{IsDedekindDomain.AKLB.finitePresentation,pi_tensorProduct_of_finitePresentation}
  If $S$ is a finite set of nonzero primes of $A$ then the natural map
  $B\otimes((\prod_{v\in S}K_v)\times(\prod_{v\notin S}A_v))\to
  (\prod_{v\in S}(B\otimes_AK_v))\times(\prod_{v\notin S}(B\otimes_AA_v))$
  is an isomorphism.
\end{corollary}
\begin{proof} Follows from the previous two theorems.
\end{proof}

Recall that $\A_K^\infty$ is the finite adeles of $K$,
defined as the restricted product of the $K_v$ with respect to the $A_v$,
where $v$ runs through the nonzero primes of $A$. Let $R$ denote the restricted
product of the $B\otimes_A K_v$ with respect to the $B\otimes_A A_v$.

\begin{corollary}
  \label{IsDedekindDomain.FiniteAdeleRing.IntegraltensorProductAlgEquiv_aux1}
  \uses{IsDedekindDomain.pi_tensorProduct}
  The natural map $B\otimes_A\A_K^\infty\to R$ is a $B$-algebra isomorphism.
\end{corollary}
\begin{proof} This follows from the previous corollary and the fact that
  tensor products commute with filtered colimits.
\end{proof}

Recall from earlier in this section that if $v$ is a finite place of $A$ then the natural map from
$B\otimes_A K_v$ to $L\otimes_KK_v$ is an isomorphism, and recall from the previous section
that the natural map from $L\otimes_KK_v$ to $\prod_{w|v}L_w$ was also an isomorphism.
This isomorphism sends $B\otimes_A A_v$ to $\prod_{w|v}B_w$ (I thank Matthew Jasper for
pointing out to me that this statement was true at all primes, not just at unramified primes).
Finally, the set of $w$ of $B$ dividing a fixed place $v$ of $A$ is finite.
Let's now formalize the abstract statement which we need. Rather than following the
notation for restricted product in the literature and writing $\A_K^\infty=\prod'_vK_v$
with the $\mathcal{O}_v$ implicit, we will write $\prod'_v(K_v,\mathcal{O}_v)$ in the below.

\begin{definition}
  \label{RestrictedProduct.relabelIso}
  Let $V$ and $W$ be index sets, and let $f:W\to V$ be a map with finite fibres.
  Let $X_v$ be sets, with subsets $C_v$, let $Y_w$ be sets with subsets $D_w$,
  and say for all $v\in V$ we're given a bijection $X_v\to\prod_{w|f(w)=v}Y_w$,
  identifying $C_v$ with $\prod_{w:f(w)=v}D_w$. Then there's an induced bijection between
  the restricted products $\prod'_v(X_v,C_v)$ and $\prod'_w(Y_w,D_w)$.
\end{definition}

\begin{corollary}
  \label{IsDedekindDomain.FiniteAdeleRing.IntegraltensorProductAlgEquiv_aux2}
  \uses{RestrictedProduct.relabelIso}
  The ring $R$ introduced above (the restricted
  product of the $B\otimes_A K_v$ with respect to the $B\otimes_A A_v$)
  is isomorphic to $\mathbb{A}_L$.
\end{corollary}
\begin{proof} Let $V$ be the finite places of $K$ and $W$ the finite places of $L$,
  let $X_v$ be $B\otimes_A K_v$, let $C_v$ be $B\otimes_A A_v$, let $Y_w$ be $L_w$,
  let $D_w$ be $B_w$ and the result follows from the previous definition, given
  theorem~\ref{IsDedekindDomain.HeightOneSpectrum.range_baseChange_comp_tensorAdicCompletionTo_eq_pi}.
\end{proof}
From this, we can deduce the theorem we claimed earlier:

\begin{theorem}
  \label{IsDedekindDomain.FiniteAdeleRing.baseChangeIntegralAlgEquiv}
  \uses{RestrictedProduct.relabelIso,
  IsDedekindDomain.FiniteAdeleRing.IntegraltensorProductAlgEquiv_aux1,
  IsDedekindDomain.FiniteAdeleRing.IntegraltensorProductAlgEquiv_aux2}
  The natural map $B\otimes_A\A_K^\infty\to\A_L^\infty$ is a $B$-algebra
  isomorphism.
\end{theorem}
\begin{proof}
  This map factors through the auxiliary ring~$R$ so the result follows
  from the previous two constructions.
\end{proof}

Because this map factors through the isomorphism $B\otimes_A\A_K^\infty\to L\otimes_K\A_K^\infty$
we can finally deduce that the natural map $L\otimes_K\A_K^\infty\to\A_L^\infty$ is an algebraic
  isomorphism.
\begin{proof}
  \proves{IsDedekindDomain.FiniteAdeleRing.baseChangeAlgEquiv}
  Follows immediately from theorem~\ref{IsDedekindDomain.FiniteAdeleRing.baseChangeIntegralAlgEquiv}
  and theorem~\ref{IsDedekindDomain.AKLB.tensorProduct_module_algEquiv}.
\end{proof}

We still need to talk about topologies though, so let's finish by doing this. Let's start with some trivialities, expressed as definitions rather than theorems
because they're constructions.

\begin{definition} If $X_v$ and $Y_v$ are families of topological spaces indexed by $v\in V$,
  if $f_v:X_v\to Y_v$ is a continuous map sending the subset $C_v\subseteq X_v$ into
  $D_v\subseteq Y_v$ then there's an induced continuous map $\prod'_v(X_v,C_v)\to\prod'_v(Y_v,D_v)$.
\end{definition}

\begin{definition} If all the $f_v$ are homeomorphisms identifying $C_v$ with $D_v$ then
  the induced map on restricted products is also a homeomorphism (proof: apply the previous
  construction to $f_v$ and their inverses)
\end{definition}

We now allow a slight change of index set. Unfortunately I don't think that we can deduce
the results just stated above from this one, in Lean, because the product of $Y_w$ over a
set of size 1 is not definitionally equal to $Y_w$.

Recall definition~\ref{RestrictedProduct.relabelIso}, giving us a bijection between two restricted
products.

\begin{theorem} In the same setup as definition~\ref{RestrictedProduct.relabelIso}
  ($V,W$ index sets, $f:W\to V$,
  $C_v\subseteq X_v$ and $D_w\subseteq Y_w$, bijections $b_v:X_v\to\prod_{w:f(w)=v}Y_w$
  identifying $C_v$ with $\prod_{w:f(w)=v}D_w$), if all the $X_v$ and $Y_w$ are furthermore
  topological spaces, all the $C_v$ and $D_w$ are open, and all the $b_v$ are homeomorphisms,
  then the induced
  map $\prod'_v(X_v,C_v)\to\prod'_w(Y_w,D_w)$ is also a homeomorphism.
\end{theorem}
\begin{proof} I have only thought about the cofinite filter case, where this
  should follow easily from the definition of the topology.
\end{proof}

\begin{corollary} $\mathbb{\A_L^\infty}$ is homeomorphic to $\prod_v(B\otimes_AK_v,B\otimes_AA_v)$.
\end{corollary}
\begin{proof} Follows from the previous theorem with $X_v=B\otimes_AK_v$ $D_w=L_w$ etc.
\end{proof}

Recall that if $$R$$ is a commutative ring, and two $$R$$-modules both have the $$R$$-module
topology, then any $$R$$-linear morphism between them is automatically continuous. We know
that $\A_L^\infty$ is $\A_K^\infty$-linearly isomorphic to $L\otimes_K\A_K^\infty$ and our claim is that
if $L\otimes_K\A_K^\infty$ is given the $\A_K^\infty$-module topology then this isomorphism is also a
homeomorphism; to prove this, we thus just need to check that $\A_L^\infty$ has the $\A_K^\infty$-module
topology. Equivalently, by the previous result, we need to check that
the restricted product topology on the $\A_K^\infty$-algebra $\prod'_v(B\otimes_AK_v,B\otimes_AA_v)$
is the $\A_K^\infty$-module topology.

We now need to make restricted products of modules into modules over restricted product of rings.
The API, which should be straightforward so we don't give details here, is: if $R_v$ are rings
with subrings $S_v$ and if $M_v$ are $R_v$-modules with $S_v$-stable submodules $N_v$, then
$\prod'_v(M_v,N_v)$ is naturally a module over $\prod'_v(R_v,S_v)$, and that $R_v$-morphisms
$M_v\to M_v'$ sending $N_v$ to $N_v'$ induce $\prod'_v(R_s,S_v)$-linear maps
$\prod'_v(M_v,N_v)\to\prod'_v(M'_v,N'_v)$. From this one deduces that isomorphisms on the factors
induce isomorphisms on the restricted products.


Now $A_v$ is a PID, so $B\otimes_AA_v$ is free (as it is finitely-generated and torsion-free).
This means that there is an isomorphism $B\otimes_AA_v\cong(A_v)^n$, which extends to an isomorphism
$B\otimes_AK_v\cong K_v^n$. These isomorphisms are also homeomorphisms. If we fix such isomorphisms
for all $v$ then we get an induced $\A_K^\infty$-module isomorphism + homeomorphism
$\prod'_v(B\otimes_AK_v,B\otimes_AA_v)=\prod'_v(K_n^n,A_v^n)$. So it suffices to prove
that the $\prod'_v(K_v,A_v)$-module $\prod'_v(K_v^n,A_v^n)$ has the $\prod'_v(K_v,A_v)$-module
topology. This follows from the fact that the product topology on two modules with the
module topology is the module topology (a fact in mathlib) and the following result.

\begin{lemma}
  If $X_v$ and $Y_v$ are topological spaces with open subspaces $C_v$ and $D_v$, then
  the obvious bijection $\prod'_v(X_v \times Y_v,C_v\times D_v) \cong
  \left(\prod'_v(X_v,C_v)\right)\times\left(\prod'_v(Y_v,D_v)\right)$ is a homeomorphism,
  where the restricted products have the restricted product topology and the binary
  product has the product topology.
\end{lemma}
\begin{proof} This should hopefully be straightforward using {\tt RestrictedProduct.continuous\_dom\_prod}
\end{proof}

As a corollary one can prove by induction on $n$ that the restricted product of $n$th powers
is homeomorphic to the $n$th power of the restricted product and this is the result
we require.

\subsection{Base change for infinite adeles}

Recall that if $K$ is a number field then the infinite adeles of $K$ are defined
to be the product $\prod_{v\mid\infty} K_v$ of all the completions of $K$ at the
infinite places.

The result we need here is that if $L/K$ is a finite extension of number fields,
then the map $K\to L$ extends to a continuous $K$-algebra map $K_\infty\to L_\infty$,
and thus to a continuous $L$-algebra isomorphism $L\otimes_KK_\infty\to L_\infty$.
Perhaps a cheap proof would be to deduce it from the fact that $K_\infty=K\otimes_{\Q}\R$.

The overall strategy is to first establish, for each infinite place $v$ of $K$, homeomorphisms
between for the completion $K_v$ and the product $\prod_{w\mid v}L_w$ of completions of $L$ at
all infinite places $w$ of $L$ lying above $v$.
We then use these homeomorphisms to construct base change for the infinite adele ring.

\subsubsection{Weak approximation at infinite places}

First, we require a preliminary result that $K$ is dense inside any product of
completions $\prod_{v\in S} K_v$ of $K$ at infinite places.

\begin{theorem}
  \label{NumberField.InfinitePlace.Completion.denseRange_algebraMap_subtype_pi}
  \lean{NumberField.InfinitePlace.Completion.denseRange_algebraMap_subtype_pi}
  \leanok
  Let $S$ be a set of infinite places of $K$. The image of $K$ under the embedding
  $K\hookrightarrow (K_v)_{v\in S}; k \mapsto (k)_v$ is dense in the product topology.
\end{theorem}
\begin{proof}
  Let $(K, v)$ denote $K$ equipped with the topology induced by the infinite place $v$.
  It suffices to show that the image of $K$ under the embedding
  $K\hookrightarrow \prod_{v\mid \infty}(K, v)$ is dense in the product topology.
  By a standard analytic argument, for each $v$ it is possible to select a sequence $(x_n^{(v)})_n$
  with the property that $x_n^{(v)} \to 1$ in $v$'s topology, while $x_n^{(v)} \to 0$ in any other
  infinite place's topology.
  Let $(z_v)_v \in \prod_{v\mid \infty}(K, v)$.
  For each $v$, the sequence $y_n := \sum_{v\mid\infty} x_n^{(v)}z_v$ in $K$ converges to
  $z_v$ in $v$'s topology.
  So the embedded image of $y_n$ in $\prod_{v\mid \infty}(K, v)$ converges to $(z_v)_v$ in the
  product topology.
\end{proof}

\subsubsection{Dimensionality of $\prod_{w\mid v}L_w$ as a $K_v$-vector space}

This subsection contains a result that the $K_v$-dimension of
$L \otimes_K K_v$ is equal to the $K_v$-dimension of $\prod_{w\mid v}L_w$.

\begin{theorem}
  \label{NumberField.InfinitePlace.Completion.finrank_pi_eq_finrank_tensorProduct}
  \lean{NumberField.InfinitePlace.Completion.finrank_pi_eq_finrank_tensorProduct}
  For a fixed infinite place $v$ of $K$, we have
  $\text{dim}_{K_v} \prod_{w\mid v} L_w = \text{dim}_{K_v} L\otimes_K K_v$.
\end{theorem}

\subsubsection{Base change at infinite places}

\begin{definition}
  \label{NumberField.InfinitePlace.Completion.piExtension}
  \lean{NumberField.InfinitePlace.Completion.piExtension}
  \leanok
  Let $v$ be an infinite place of $K$. There is a continuous $K$-algebra homomorphism
  $K_v \to \prod_{w\mid v}L_w$, whose restriction to $K$ corresponds to the global embedding
  of $K$ into $(L_w)_w$.
\end{definition}

The map in~\ref{NumberField.InfinitePlace.Completion.piExtension} can be lifted to an
$L$-algebra homomorphism defined on $L\otimes_K K_v$.

\begin{definition}
  \label{NumberField.InfinitePlace.Completion.baseChange}
  \lean{NumberField.InfinitePlace.Completion.baseChange}
  \leanok
  \uses{NumberField.InfinitePlace.Completion.piExtension}
  Let $v$ be an infinite place of $K$. There is a natural $L$-algebra homomorphism
  $L\otimes_K K_v \to \prod_{w\mid v}L_w$, whose restriction to $1\otimes_K K_v$ corresponds to
  the map in~\ref{NumberField.InfinitePlace.Completion.piExtension}.
\end{definition}

\begin{theorem}
  \label{NumberField.InfinitePlace.Completion.baseChange_surjective}
  \lean{NumberField.InfinitePlace.Completion.baseChange_surjective}
  \leanok
  \uses{NumberField.InfinitePlace.Completion.denseRange_algebraMap_subtype_pi}
  For a fixed infinite place $v$ of $K$, the map $L\otimes_K K_v \to\prod_{w\mid v}L_w$ is
  surjective.
\end{theorem}
\begin{proof}
  Let $(x_i)_i$ be a $K_v$-basis of $\prod_{w\mid v}L_w$. By the density of $L$ in
  $\prod_{w\mid v}L_w$
  (Theorem~\ref{NumberField.InfinitePlace.Completion.denseRange_algebraMap_subtype_pi}), we can
  find $\alpha_i \in L$ arbitrarily close to $x_i\prod_{w\mid v}L_w$ with respect to the sup norm
  when embedded globally in $\prod_{w\mid v}L_w$.
  In particular, it is possible to choose such $\alpha_i$ so that the matrix representing
  the vector $((\alpha_i)_{w \mid v})_i$ in the basis $(x_i)_i$ has non-zero determinant.
  Since $(\alpha_i)_{w \mid v}$ is the image of $1\otimes \alpha_i$ under base change, we have
  that $(1 \otimes \alpha_i)_i$ also forms a basis of $L\otimes_K K_v$, and base change
  is surjective.
\end{proof}

\begin{theorem}
  \label{NumberField.InfinitePlace.Completion.baseChange_injective}
  \lean{NumberField.InfinitePlace.Completion.baseChange_injective}
  \leanok
  \uses{NumberField.InfinitePlace.Completion.finrank_pi_eq_finrank_tensorProduct}
  For a fixed infinite place $v$ of $K$, the map $L\otimes_K K_v \to\prod_{w\mid v}L_w$ is
  injective.
\end{theorem}
\begin{proof}
  The $L$-algebra map $L\otimes_K K_v \to\prod_{w\mid v}L_w$ can equivalently be thought of
  as $K_v$-linear, which is injective, because it is surjective by
  Theorem~\ref{NumberField.InfinitePlace.Completion.baseChange_surjective}, and both side have the same
  $K_v$-dimension by
  Theorem~\ref{NumberField.InfinitePlace.Completion.finrank_pi_eq_finrank_tensorProduct}.
\end{proof}

We have established that the map of
Definition~\ref{NumberField.InfinitePlace.Completion.baseChange} gives an $L$-algebra isomorphism
between $L\otimes_K K_v$ and $\prod_{w\mid v}L_w$.
The left-hand side is given the $K_v$-module topology, while we show that the right-hand side also
has the $K_v$-module topology.
\begin{theorem}
  \label{NumberField.InfinitePlace.Completion.instIsModuleTopologyValEqComapAlgebraMap_fLT}
  \lean{NumberField.InfinitePlace.Completion.instIsModuleTopologyValEqComapAlgebraMap_fLT}
  \leanok
  If $w \mid v$ is an infinite place of $L$ lying above the infinite place $v$ of $K$, then
  $L_w$ has the $K_v$-module topology.
\end{theorem}
\begin{proof}
  Because $L_w$ is a finite-dimensional normed $K_v$ vector space, there exists a $K_v$-linear
  linear homeomorphism $L_w \cong K_v^n$, from which $L_w$ has the $K_v$-module topology.
\end{proof}

\begin{theorem}
  \label{NumberField.InfinitePlace.Completion.baseChangeEquiv}
  \lean{NumberField.InfinitePlace.Completion.baseChangeEquiv}
  \leanok
  \uses{NumberField.InfinitePlace.Completion.baseChange,
    NumberField.InfinitePlace.Completion.baseChange_surjective,
    NumberField.InfinitePlace.Completion.baseChange_injective,
    NumberField.InfinitePlace.Completion.instIsModuleTopologyValEqComapAlgebraMap_fLT}
  Let $v$ be an infinite place of $K$. There is a natural $L$-algebra homeomorphism
  $L\otimes_K K_v \cong_L \prod_{w\mid v}L_w$, whose restriction to $1\otimes_K K_v$ corresponds to
  the map in~\ref{NumberField.InfinitePlace.Completion.piExtension}.
\end{theorem}
\begin{proof}
  The map in~\ref{NumberField.InfinitePlace.Completion.baseChange} is an $L$-algebra
  isomorphism by Theorem~\ref{NumberField.InfinitePlace.Completion.baseChange_surjective}
  and Theorem~\ref{NumberField.InfinitePlace.Completion.baseChange_injective}.
  Every $K_v$-algebra isomorphism between two $K_v$-module topological spaces is a homeomorphism.
  Since the $L$-algebra isomorphism of
  Definition~\ref{NumberField.InfinitePlace.Completion.baseChange} can equivalently be viewed as
  a $K_v$-algebra isomorphism, it is also a homeomorphism.
\end{proof}

\begin{theorem}
  \label{NumberField.InfinitePlace.Completion.piEquiv}
  \lean{NumberField.InfinitePlace.Completion.piEquiv}
  \leanok
  \uses{NumberField.InfinitePlace.Completion.baseChangeEquiv}
  Let $v$ be an infinite place of $K$. There is a natural $K_v$-linear homeomorphism
  $K_v^{[L:K]} \cong_{K_v} \prod_{w\mid v}L_w$.
\end{theorem}
\begin{proof}
  Compose the $K_v$-linear isomorphism $K_v^{[L:K]} \cong \prod_{w\mid v}L_w$ with the $K_v$-linear
  version of base change~\ref{NumberField.InfinitePlace.Completion.baseChangeEquiv} to get the
  isomorphism in the statement.
  Since both sides have the $K_v$-module topology, then this is also a homeomorphism.
\end{proof}

\subsubsection{Base change for the infinite adele ring}

First, we induce a $K_{\infty}$-algebra on $L_{\infty}$ from the action of each $K_v$ on
$\prod_{w\mid v}L_w$.
Specifically, this means that for $x \in K_{\infty}$ and $y \in L_{\infty}$, we have
$(x \cdot y)_w = x_{v_w} \cdot y_w$, where $v_w$ is the restriction of $w$ to $K$.
We show that $L_{\infty}$ has the $K_{\infty}$-module topology.

\begin{theorem}
  \label{NumberField.InfiniteAdeleRing.piEquiv}
  \lean{NumberField.InfiniteAdeleRing.piEquiv}
  \leanok
  \uses{NumberField.InfinitePlace.Completion.piEquiv}
  There is a natural $K_{\infty}$-linear homeomorphism
  $K_{\infty}^{[L:K]} \cong_{K_{\infty}} L_{\infty}$.
\end{theorem}
\begin{proof}
  Using the isomorphisms $K_v^{[L:K]} \cong_{K_v} \prod_{w\mid v}L_w$ from
  Theorem~\ref{NumberField.InfinitePlace.Completion.piEquiv}, we clearly have a bijection
  $K_{\infty}^{[L:K]} \cong \prod_v\prod_{w \mid v} L_w \cong \prod_w L_w$.
  The $K_v$-linearity of each component isomorphism extends to $K_{\infty}$-linearity if the
  action of $\prod_v K_v$ on $\prod_w L_w$ is constant on the fibers of the restriction map on
  infinite places.
  In other words, if, for all $x \in K_{\infty}$ and $y \in L_{\infty}$, we have
  $(x \cdot y)_w = x_{v_w} \cdot y_w$, which is true by definition.
\end{proof}

\begin{theorem}
  \label{NumberField.InfiniteAdeleRing.instIsModuleTopology_fLT}
  \lean{NumberField.InfiniteAdeleRing.instIsModuleTopology_fLT}
  \leanok
  \uses{NumberField.InfiniteAdeleRing.piEquiv}
  $L_{\infty}$ has the $K_{\infty}$-module topology.
\end{theorem}
\begin{proof}
  Since $L_{\infty}$ is homeomorphic to a finite product of $K_{\infty}$ as a $K_{\infty}$-vector
  space, it has the $K_{\infty}$-module topology.
\end{proof}

\begin{theorem}
  \label{NumberField.InfiniteAdeleRing.baseChangeAlgEquiv}
  \lean{NumberField.InfiniteAdeleRing.baseChangeAlgEquiv}
  \leanok
  \uses{NumberField.InfinitePlace.Completion.baseChangeEquiv}
  There is a natural $L$-algebra isomorphism
  $L \otimes_K K_{\infty} \cong_L L_{\infty}$.
\end{theorem}
\begin{proof}
  This follows from the following chain of isomorphisms:
  \[
    L \otimes_K K_{\infty} \cong_L \prod_v (L \otimes_K K_v) \cong_L
      \prod_v \prod_{w\mid v}L_w \cong_L L_{\infty} .
  \]
  The first isomorphism is the standard $L$-algebra isomorphism
  $L \otimes_K \prod_v K_v \cong_L \prod_v (L \otimes_K K_v)$.
  The second isomorphism is given by the component $L$-algebra isomorphisms
  $L \otimes_K K_v \cong_L \prod_{w\mid v}L_w$ from
  Theorem~\ref{NumberField.InfinitePlace.Completion.baseChangeEquiv}.
\end{proof}

It remains to show that the map in~\ref{NumberField.InfiniteAdeleRing.baseChangeAlgEquiv} is a
homeomorphism.
Since both sides have the $K_{\infty}$-module topology, and since the $L$-algebra isomorphism
of~\ref{NumberField.InfiniteAdeleRing.baseChangeAlgEquiv} can equivalently be viewed as a
$K_{\infty}$-linear isomorphism, it is also a homeomorphism.

\begin{theorem}
  \label{NumberField.InfiniteAdeleRing.baseChangeEquiv}
  \lean{NumberField.InfiniteAdeleRing.baseChangeEquiv,
    NumberField.InfiniteAdeleRing.instIsModuleTopology_fLT}
  \leanok
  \uses{NumberField.InfiniteAdeleRing.baseChangeAlgEquiv}
  If $K\to L$ is a ring homomorphism between two number fields then there is a natural isomorphism
  (both topological and algebraic) $L\otimes_KK_\infty\cong L_\infty$.
\end{theorem}
\begin{proof}
  Since both sides of the $L$-algebra isomorphism
  in~\ref{NumberField.InfiniteAdeleRing.baseChangeAlgEquiv} have the $K_{\infty}$-module topology,
  and since the isomorphism can equivalently be viewed as a $K_{\infty}$-linear isomorphism,
  it is also a homeomorphism.
\end{proof}


\subsection{Base change for adeles}

From the previous results we deduce immediately that if $L/K$ is a finite extension
of number fields then there's a natural (topological and algebraic) isomorphism
$L\otimes_K\A_K\to \A_L$.

\begin{theorem}
  \label{NumberField.AdeleRing.baseChangeEquiv}
  \lean{NumberField.AdeleRing.baseChangeEquiv}
  \leanok
  \uses{NumberField.InfiniteAdeleRing.baseChangeEquiv,
  IsDedekindDomain.FiniteAdeleRing.baseChangeAlgEquiv}
  If $K\to L$ is a ring homomorphism between two number fields then there is a natural isomorphism
  (both topological and algebraic) $L\otimes_K\A_K\cong\A_L$.
\end{theorem}
\begin{proof}
  \leanok
  Follows from the previous results.
\end{proof}

Something else we shall need:

\begin{theorem}
  \label{NumberField.AdeleRing.baseChange_moduleTopology}
  \lean{NumberField.AdeleRing.instIsModuleTopology}
  \leanok
  \uses{NumberField.AdeleRing.baseChangeEquiv}
  If $K\to L$ is a ring homomorphism between two number fields then the topology on $\A_L$
  is the $\A_K$-module topology, where the module structure comes from the
  natural map $\A_K\to\A_L$.
\end{theorem}
\begin{proof}
  \leanok
  Indeed $\A_L\cong L\otimes_K\A_K$ is a homeomorphism, and
  the right hand side has the $\A_K$-module topology.
\end{proof}

\section{Discreteness and compactness}

We need that if $K$ is a number field then
$K\subseteq\mathbb{A}_K$ is discrete, and the quotient (with the
quotient topology) is compact. Here is a proposed proof.

\begin{theorem}
  \lean{Rat.AdeleRing.zero_discrete}
  \label{Rat.AdeleRing.zero_discrete}
  \leanok
  There's an open subset of $\A_{\Q}$ whose intersection with $\Q$ is $\{0\}$.
\end{theorem}
\begin{proof}
  Use $\prod_p{\Z_p}\times(-1,1)$. Any rational $q$ in this set is a $p$-adic
  integer for all primes $p$ and hence (writing it in lowest terms as $q=n/d$)
  satisfies $p\nmid d$, meaning that $d=\pm1$ and thus $q\in\Z$. The fact
  that $q\in(-1,1)$ implies $q=0$.
\end{proof}

\begin{theorem}
  \lean{NumberField.AdeleRing.zero_discrete}
  \label{NumberField.AdeleRing.zero_discrete}
  \uses{Rat.AdeleRing.zero_discrete,NumberField.AdeleRing.baseChangeEquiv}
  \leanok
  There's an open subset of $\A_{K}$ whose intersection with $K$ is $\{0\}$.
\end{theorem}
\begin{proof}
  By a previous result, we have $\A_K=K\otimes_{\Q}\A_{\Q}$.
  Choose a basis of $K/\Q$; then $K$ can be identified with $\Q^n\subseteq(\A_{\Q})^n$
  and the result follows from the previous theorem.
\end{proof}

\begin{theorem}
  \lean{NumberField.AdeleRing.discrete}
  \label{NumberField.AdeleRing.discrete}
  \uses{NumberField.AdeleRing.zero_discrete}
  \leanok
  The additive subgroup $K$ of $\A_K$ is discrete.
\end{theorem}
\begin{proof}
  If $x\in K$ and $U$ is the open subset in the previous lemma, then
  it's easily checked that $K\cap U=\{0\}$ implies $K\cap (U+x)=\{x\}$,
  and $U+x$ is open.
\end{proof}

For compactness we follow the same approach.

\begin{theorem}
  \lean{Rat.AdeleRing.cocompact}
  \label{Rat.AdeleRing.cocompact}
  \leanok
  The quotient $\A_{\Q}/\Q$ is compact.
\end{theorem}
\begin{proof}
  The space $\prod_p\Z_p\times[0,1]\subseteq\A_{\Q}$ is a product of compact spaces
  and is hence compact. I claim that it surjects onto $\A_{\Q}/\Q$. Indeed,
  if $a\in\A_{\Q}$ then for the finitely many prime numbers $p\in S$ such that $a_p\not\in\Z_p$
  we have $a_p\in\frac{r_p}{p^{n_p}}+\Z_p$ with $r_p/p^{n_p}\in\Q$, and
  if $q=\sum_{p\in S}\frac{r_p}{p^{n_p}}\in\Q$ then $a-q\in \prod_p\Z_p\times\R$.
  Now just subtract $\lfloor a_{\infty}-q\rfloor$ to move into $\prod_p\Z_p\times[0,1)$
  and we are done.
\end{proof}

\begin{theorem}
  \lean{NumberField.AdeleRing.cocompact}
  \label{NumberField.AdeleRing.cocompact}
  \uses{Rat.AdeleRing.cocompact,NumberField.AdeleRing.baseChangeEquiv}
  \leanok
  The quotient $\A_K/K$ is compact.
\end{theorem}
\begin{proof}
  We proceed as in the discreteness proof above, by reducing to $\Q$. As before, choosing
  a $\Q$-basis of $K$ gives us $\A_K/K\cong(\A_{\Q}/\Q)^n$ so the result follows from
  the previous theorem.
\end{proof}


%%FILE blueprint/src/chapter/FrobeniusProject.tex
\chapter{Miniproject: Frobenius elements}\label{Frobenius_project}

\section{Status}

This miniproject has been a success: the main results
are sorry-free and merged into mathlib
(see \href{https://github.com/leanprover-community/mathlib4/pull/19926}{this PR}).
As a result there will be no more work
on this miniproject in the FLT repo. Below is a fairly
detailed sketch of the argument used.

\section{Introduction and goal}

When this project started, I had thought that the existence of Frobenius elements was
specific to the theory of local and global fields, and a slightly more general result
held for Dedekind domains. Then Joel Riou pointed out on the Lean Zulip
an extremely general result from from Bourbaki's Commutative Algebra
(Chapter V, Section 2, number 2, theorem 2, part (ii)). This beautiful
result is surely what we want to see in mathlib. Before we state Bourbaki's
theorem, let us set the scene.

\section{Statement of the theorem}

The set-up throughout this project:
$G$ is a finite group acting (via ring isomorphisms) on a commutative ring $B$,
and $A$ is the subring of $G$-invariants.

Say $Q$ is a prime ideal of $B$, whose pullback to $A$ is the prime ideal $P$.
Note that $G$ naturally acts on the ideals of $B$. Let's define the
\emph{decomposition group} $D_Q$ of $Q$ to be the subgroup of $G$ which
stabilizes $Q$ (just to be clear: $g\in D_Q$ means
$\{g\cdot q\, :\, q \in Q\}=Q$, not $\forall q\in Q, g\cdot q=q$).

Let $L$ be the field of fractions of the integral domain $B/Q$, and let $K$ be the
field of fractions of the subring $A/P$. Then $L$ is naturally a $K$-algebra.
In this generality $L/K$ may not even be finite or Galois, but we can still talk about
$\Aut(L/K)$.

In the next definition we write down a group homomorphism $\phi$ from $D_Q$ to $\Aut(L/K)$.

\begin{definition}
  \label{IsFractionRing.stabilizerHom}
  \lean{IsFractionRing.stabilizerHom}
  Choose $g\in D_Q$. Then the action of $g$ on $B$ gives us an induced
  $A/P$-algebra automorphism of $B/Q$ which extends to a $K$-algebra automorphism $\phi(g)$ of $L$.
  This construction $g\mapsto \phi(g)$ defines a group homomorphism from $D_Q$
  to $\Aut(L/K)$ (all the proofs implicit in the definition here are straightforward).
  \mathlibok
\end{definition}

The theorem we want in this mini-project is
\begin{theorem}
  \label{IsFractionRing.stabilizerHom_surjective}
  \lean{IsFractionRing.stabilizerHom_surjective}
  \uses{IsFractionRing.stabilizerHom}
  \mathlibok
  The map $g\mapsto \phi_g$ from $D_Q$ to $\Aut(L/K)$ defined above is surjective.
\end{theorem}

The goal of this mini-project is to get this theorem formalised and ideally into mathlib.

In particular, $\Aut(L/K)$ is finite
as a corollary. What is so striking about this theorem to me is that the only finiteness hypothesis
is on the group $G$ which acts; there are no finiteness or Noetherian
hypotheses on the rings at all.

As a trivial consequence we get Frobenius elements for finite Galois extensions in both
the local and global field setting, as $\Aut(L/K)$ is just a Galois group of finite fields
in these cases, so by surjectivity we can lift a Frobenius element.

Even though $G$ is finite, it is possible in characteristic $p>0$ for the extension $L/K$ to be
infinite (and mostly inseparable). The theorem implies that $\Aut(L/K)$ is always finite;
what is actually happening is that $L/K$ is algebraic and normal, and its maximal separable
subextension is finite of degree at most $|G|$. However, we can prove surjectivity directly
without reference to this maximal separable subextension.

\subsection{Examples}

These do not need to be formalised.

The basic example is when $B$ is the ring of integers of a finite Galois extension of $\Q$
(or equivalently the field obtained by adjoining all of the roots of
a monic polynomial with rational coefficients).
Then the Galois group $G$ acts on $B$, and the invariants $A$ are just $B\cap\Q=\Z$. If $Q$
is a nonzero prime ideal of $B$ then a standard result is that $B/Q$ is a finite field. Let's
call this field $k$. We deduce that $P=Q\cap \Z$ is a nonzero ideal (as $\Z/P$ injects into $B/Q$)
and hence must be
the principal ideal generated by a prime number $p$. This number $p$ is ``the prime below $Q$'',
and also the characteristic of $k$.

If $D$ is the subgroup of $G$ stabilizing $Q$ as a set, then $D$ acts on $k$, so we get
a map from $D$ to $\Gal(k/\F_p)$. The theorem states that this map is surjective.
Its kernel is the inertia subgroup of $Q$, which for all but finitely many primes of $B$
is the trivial group. So in these cases we get an \emph{isomorphism} from $D$ to $\Gal(k/\F_p)$
meaning that $D$ is cyclic, and furthermore has two canonical generators, one called $\Frob_Q$
(by Artin) and the other one unfortunately also called $\Frob_Q$ (by Deligne), which are inverses
to each other. Some people think that both Frobenii are canonical, some people (for example
the Shimura variety crowd) think that Deligne's choice is more canonical, and others
(for example the Heegner point crowd) think that Artin's choice is more canonical.
It was this example which convinced me that the word ``canonical'' should be banished
from mathematics if not accompanied by a clear explanation of what the author means
by the word. I know several examples of papers in the literature where $\Frob_Q$ is
used with no explanation as to which one it is. And sometimes it doesn't matter.
For example, with either choice of normalization the following is true.
If $Q$ and $Q'$ are two primes above $p$ then there's some $g\in G$ such that
$gQ=Q'$ and one can deduce from this that $\Frob_Q$ and $\Frob_{Q'}$ are conjugate. In particular
if $G$ is abelian then $\Frob_Q$ and $\Frob_{Q'}$ are equal, so we can call them both $\Frob_p$.

An example which demonstrates that things can get a bit stranger in characteristic $p$ is the
following (we restrict to $p=2$ but a generalisation of this pathology exists for all $p>0$).
We let $B=\mathbb{F}_2[X,Y]$ and $G=\{1,\tau\}$ with the involution $\tau$
fixing $Y$ and switching $X$, $X+Y$, i.e., $(\tau f)(X,Y)=f(X+Y, Y)$.
Hence $\tau$ fixes the sum and product of these two polynomials, and thus
$A\supseteq \mathbb{F}_2[X^2+XY,Y]$. One can check (and this needs checking) that this
is in fact an equality, and I believe that $B$ is free of
rank 2 over $A$ with basis $\{1,X\}$.
Now set $Q=(Y)$ so $P=YA$ and the quotient $(B/Q)/(A/P)$ is
$\mathbb{F}_2[X]/\mathbb{F}_2[X^2]$, and the corresponding extension
of the fields of fractions of these integral domains is a radical extension
$\mathbb{F}_2(X)/\mathbb{F}_2(X^2)$ of degree 2.
In other words, it is finite and normal, but not separable.

It appears that this construction behaves well with respect to tensor
products over $\mathbb{F}_2$ and filtered colimits (I have not thought about
how to make this rigorous), and assuming this is true, it will explain Bourbaki's
counterexample to the hope that $L/K$ is always finite. We now describe the counterexample
(exercise 9 of number 2)

The example is from the exercises in Bourbaki (exercise 9 of section 2 above, found at the end
of Chapter V). We let $B=\mathbb{F}_2[X_0,Y_0,X_1,Y_1,X_2,\ldots]$
be a polynomial ring in infinitely many variables $X_i$ and $Y_i$ indexed by the
naturals (or indeed by any infinite set), and $G$ is cyclic of order 2 with
the generator acting on $B$ via $X_n\mapsto X_n+Y_n$ and $Y_n\mapsto Y_n$.
If $Q$ is the ideal generated by the $Y_n$ then now $L/K$ is a radical extension of
infinite degree; the generator swaps $X_n$ and $X_n+Y_n$, so it fixes
their product $a_n\in A$, which becomes $X_n^2$ modulo $Q$, so all
of the $X_{i}$ will be algebraically independent in $L/K$ and $X_{i}^2\in K$.

\section{The extension \texorpdfstring{$B/A$}{B/A}.}

The precise set-up we'll work in is the following. We fix $G$ a finite group acting
on $B$ a commutative ring, and we have another commutative ring $A$ such
that $B$ is an $A$-algebra and the image of $A$ in $B$ is precisely the $G$-invariant
elements of $B$. We don't ever need the map $A\to B$ to be injective so we don't assume this.

We start with a construction which is fundamental to everything,
and which explains why we need $G$ to be finite.

\begin{definition}
  \label{MulSemiringAction.charpoly}
  \lean{MulSemiringAction.charpoly}
  \mathlibok
  If $b\in B$ then define the \emph{characteristic polynomial}
  $F_b(X) \in B[X]$ of $b$ to be $\prod_{g\in G}(X-g\cdot b)$.
\end{definition}

Clearly $F_b$ is a monic polynomial.

\begin{lemma}
  \label{MulSemiringAction.monic_charpoly}
  \lean{MulSemiringAction.monic_charpoly}
  \uses{MulSemiringAction.charpoly}
  \mathlibok
  $F_b$ is monic.
\end{lemma}
\begin{proof}
  \mathlibok
  Obvious.
\end{proof}

It's also clear that $F_b$ has degree $|G|$ and has $b$ as a root.
Also $F_b$ is $G$-invariant, because acting by some $\gamma\in G$
just permutes the order of the factors. Hence we can descend $F_b$
to a monic polynomial $M_b(X)$ of degree $|G|$ in $A[X]$. We will
also refer to $M_b$ as the characteristic polynomial of $b$, even though
it's not even well-defined if the map $A\to B$ isn't injective.

\begin{lemma}
  \label{Algebra.IsInvariant.charpoly_mem_lifts}
  \lean{Algebra.IsInvariant.charpoly_mem_lifts}
  \uses{MulSemiringAction.charpoly}
  \mathlibok
  $F_b$ is the lift of some monic polynomial $M_b$ in $A[X]$.
\end{lemma}
\begin{proof}
  \mathlibok
  The coefficients of $F_b$ are $G$-invariant, and thus lie in the image of $A$.
\end{proof}

\begin{theorem}
  \label{Algebra.IsInvariant.isIntegral}
  \lean{Algebra.IsInvariant.isIntegral}
  \mathlibok
  $B/A$ is integral.
\end{theorem}
\begin{proof}
  \uses{MulSemiringAction.monic_charpoly, Algebra.IsInvariant.charpoly_mem_lifts}
  \mathlibok
  Use $M_b$.
\end{proof}

\section{The extension \texorpdfstring{$(B/Q)/(A/P)$}{(B/Q)/(A/P)}.}

Note that $P$ and $Q$ are primes, so the quotients $A/P$ and $B/Q$ are integral domains.

The following technical lemma constructs an element of $B$ with nice characteristic polynomial
modulo $Q$.

\begin{lemma}
  \label{fixed_of_fixed1_aux1}
  \mathlibok
  There exist elements $a,b \in B$, with $a \notin Q$ and $a$ in the image of $A$ such that
  for all $g\in G$,
  \begin{itemize}
    \item If $g \cdot Q = Q$, then $g \cdot b \equiv a \pmod{Q}$.
    \item If $g \cdot Q \neq Q$, then $g \cdot b \equiv 0 \pmod{Q}$.
  \end{itemize}
\end{lemma}
\begin{proof}
  \uses{Algebra.IsInvariant.charpoly_mem_lifts}
  The ideals $g \cdot Q \neq Q$ are not contained in $Q$.
  Since $Q$ is a prime ideal, this implies that the intersection of all $g \cdot Q \neq Q$ is
  still not contained in $Q$.
  Then we can find an element $c \notin Q$ with $c \in g \cdot Q$ for all $g \cdot Q \neq Q$.
  Let $F_c$ be the characteristic polynomial of $c$, and write $F_c(X) \equiv X^j R(X) \pmod{Q}$.
  Let $a$ be the coefficient of $X^j$ in $F_c(X)$, and choose $R(X)$ so that $R(0) = a$.
  Let $b = R(0) - R(c)$.
  Note that $F_c(c) = 0$ and $c \not\equiv 0 \pmod{Q}$, so $R(c) \equiv 0 \pmod{Q}$.
  Then $b \equiv a \pmod{Q}$, so $g \cdot b \equiv a \pmod{Q}$ whenever $g \cdot Q = Q$.
  But if $g \cdot Q \neq Q$, then $c \equiv 0 \pmod {g \cdot Q}$.
  Then $b \equiv 0 \pmod {g \cdot Q}$, so $g \cdot b \equiv 0 \pmod{Q}$
  whenever $g \cdot Q \neq Q$.
  \mathlibok
\end{proof}

A slight modification allows us to take an element of $B$ fixed by $D_Q$ as input.

\begin{lemma}
  \label{fixed_of_fixed1_aux2}
  \mathlibok
  Let $b_0 \in B$.
  Suppose that the image of $b_0$ in the quotient $B/Q$ is fixed by the stabilizer subgroup
  $D_Q$.
  Then there exist elements $a,b \in B$, with $a\notin Q$ and $a$ in the image of $A$ such that
  for all $g\in G$,
  \begin{itemize}
    \item If $g \cdot Q = Q$, then $g \cdot b \equiv ab_0 \pmod{Q}$.
    \item If $g \cdot Q \neq Q$, then $g \cdot b \equiv 0 \pmod{Q}$.
  \end{itemize}
\end{lemma}
\begin{proof}
  \uses{fixed_of_fixed1_aux1}
  Multiply the $b$ from~\ref{fixed_of_fixed1_aux1} by $b_0$.
  \mathlibok
\end{proof}

\section{The extension \texorpdfstring{$L/K$}{L/K}.}

Let $L^{D_Q}$ denote the fixed field of the action $D_Q$ on $L$.
Our strategy for proving surjectivity of $D_Q \to \Aut(L/K)$ will be to write this map as the
composition $D_Q \to \Aut(L/L^{D_Q}) \to \Aut(L/K)$.

The surjectivity of the first map is a general fact of Galois theory.

\begin{theorem}
  \label{FixedPoints.toAlgAut_surjective}
  \lean{FixedPoints.toAlgAut_surjective}
  Let $H$ be a finite group acting on a field $F$ by field automorphisms.
  Then the map $H \to \Aut(F/F^H)$ is surjective.
  \mathlibok
\end{theorem}
\begin{proof}
  This is a general fact of Galois theory and was already in mathlib.
  \mathlibok
\end{proof}

For surjectivity of the second map, we need to know that every element of $L^{D_Q}$ is fixed by
$\Aut(L/K)$.
We first show this for elements of $B/Q$ fixed by $D_Q$.

\begin{proposition}
  \label{fixed_of_fixed1}
  \uses{IsFractionRing.stabilizerHom}
  \mathlibok
  Let $b_0 \in B/Q$.
  Suppose that $b_0$ is fixed by the stabilizer subgroup $D_Q$.
  Then $b_0$ is fixed by $\Aut(L/K)$.
\end{proposition}
\begin{proof}
  \uses{fixed_of_fixed1_aux1, fixed_of_fixed1_aux2}
  Let $a,b\in B$ be elements from~\ref{fixed_of_fixed1_aux2}.
  Let $M_b \in A[X]$ be the characteristic polynomial of $b$.
  We can map $M_b$ to $L[X]$ in two different ways: via $B[X]$ and via $K[X]$.
  Going via $B[X]$ tells us that the image of $M_b(X)$ in $L[X]$ is exactly
  $$(X - ab_0)^{\lvert D_Q\rvert}X^{\lvert G\rvert-\lvert D_Q\rvert}.$$
  But going via $K[X]$ tells us that this image lies in $K[X]$, so we must have $ab_0\in K$.
  Then $ab_0$ is fixed by $\Aut(L/K)$, and $a$ is nonzero in $L$ (since $a\notin Q$),
  so $b_0$ is fixed by $\Aut(L/K)$.
  \mathlibok
\end{proof}

Now we upgrade this to elements of $L$ fixed by $D_Q$.
The following lemma will allow us to lift the denominator from $B/Q$ to $A/P$.

\begin{lemma}
  \label{IsAlgebraic.exists_smul_eq_mul}
  \lean{IsAlgebraic.exists_smul_eq_mul}
  \mathlibok
  If $R/S$ is an algebraic extension of integral domains, then any fraction $a/b$ with $a,b\in R$
  can be written as $c/d$ with $c\in R$ and $d\in S$.
\end{lemma}
\begin{proof}
  If $f\in S[X]$ satisfies $f(b)=0$, then $f(0)\in S$ is a multiple of $b$.
  If $f(0)=bx\in S$, then $a/b=(ax)/(bx)$ as desired.
  \mathlibok
\end{proof}

\begin{proposition}
  \label{fixed_of_fixed2}
  \mathlibok
  Let $x \in L$.
  Suppose that $x$ is fixed by the stabilizer subgroup $D_Q$.
  Then $x$ is fixed by the automorphism group $\Aut(L/K)$.
\end{proposition}
\begin{proof}
  \uses{fixed_of_fixed1, Algebra.IsInvariant.isIntegral,
    IsAlgebraic.exists_smul_eq_mul}
  Since $(B/Q)/(A/Q)$ is algebraic by~\ref{Algebra.IsInvariant.isIntegral},
  ~\ref{IsAlgebraic.exists_smul_eq_mul} let's us write $x=b/a$ for $b \in B/Q$ and $a \in A/P$.
  Then $b$ is fixed by the stabilizer subgroup $D_Q$, and it suffices to show that
  $b$ is fixed by the automorphism group $\Aut(L/K)$.
  But this is exactly~\ref{fixed_of_fixed1}.
  \mathlibok
\end{proof}

Combining this with~\ref{FixedPoints.toAlgAut_surjective} finishes the proof.

\begin{proof}[Proof of main theorem]
  \proves{IsFractionRing.stabilizerHom_surjective}
  \uses{fixed_of_fixed2, FixedPoints.toAlgAut_surjective}
  The map $D_Q \to \Aut(L/L^{D_Q})$ is surjective by~\ref{FixedPoints.toAlgAut_surjective}.
  For surjectivity of $\Aut(L/L^{D_Q}) \to \Aut(L/K)$, let $\sigma$ be a field automorphism of $L$
  fixing $K$ pointwise.
  We must show that $\sigma$ automatically fixes $L^{D_Q}$ pointwise.
  But this is exactly~\ref{fixed_of_fixed2}.
  Thus, the composition $D_Q \to \Aut(L/L^{D_Q}) \to \Aut(L/K)$ is surjective.
  \mathlibok
\end{proof}


%%FILE blueprint/src/chapter/FujisakiProject.tex
\chapter{Miniproject: Fujisaki's Lemma}\label{Fujisaki_project}

\section{The goal}

There is an idelic compactness statement which encapsulates both finiteness of the class
group of a number field and Dirichlet's units theorem about the rank of the unit group.
In fact there is even a noncommutative version of this statement. In John Voight's
book~\cite{voightbook} this is Main Theorem 27.6.14(a) and Voight calls it Fujisaki’s lemma.
I know nothing of the history but I'm happy to adopt this name. In the quaternion algebra
miniproject we will use this compactness result to prove finite-dimensionality of a
space of quaternionic modular forms.

\section{Initial definitions}

Let $K$ be a field. A \emph{central simple $K$-algebra} is a $K$-algebra~$B$ (not necessarily
commutative) with centre $K$ such that $B$ has exactly two two-sided ideals, namely ${0}$ and $B$
(or $\bot$ and $\top$, as Lean would call them). We will be concerned
only with central simple $K$-algebras which are finite-dimensional as $K$-vector spaces, and
when $K$ is clear we will just refer to them as central simple algebras. We remark that a
4-dimensional central simple algebra is called a \emph{quaternion algebra}; we will have
more to say about these later on.

Matrix algebras $M_n(K)$ are examples of finite-dimensional central simple $K$-algebras.
If $K=\bbC$ (or more generally if $K$ is algebraically closed)
then matrix algebras are the only finite-dimensional examples
up to isomorphism. There are other examples over the reals: for example Hamilton's quaternions
$\bbH:=\R\oplus\R i\oplus\R j\oplus\R k$ with the usual rules $i^2=j^2=k^2=-1$,
$ij=-ji=k$ etc, are an example of a central simple $\R$-algebra (and a quaternion algebra), and
matrix algebras over $\bbH$ are other central simple $\R$-algebras.
For a general field $K$
one can make an analogue of Hamilton's quaternions $K\oplus Ki\oplus Kj\oplus Kk$ with the
same multiplication rules ($i^2=-1$ and so on) to describe the multiplication, and if the characteristic
of~$K$ isn't 2
then this is a quaternion algebra (which may or may not be isomorphic to $M_2(K)$ in this
generality).

Some central simple algebras~$B$ are \emph{division algebras}, meaning that they are division
rings, or equivalently that every nonzero $b\in B$ has a two-sided inverse. For example
Hamilton's quaternions are a division algebra over $\R$,
because $(x+yi+zj+tk)(x-yi-zj-tk)=x^2+y^2+z^2+t^2$, so the inverse
of a nonzero $x+yi+zj+tk$ is $(x-yi-zj-tk)/(x^2+y^2+z^2+t^2)$.
However $2\times 2$ matrices over a field~$K$, whilst being a central simple algebra
over~$K$, are never a division algebra
(even if $K=\bbC$) because a nonzero matrix with determinant zero such as
$\begin{pmatrix}1&0\\0&0\end{pmatrix}$ has no inverse.

\section{Enter the adeles}

The adeles of a number field are discussed in far more detail
in the adele miniproject \ref{Adele_miniproject}. We just recall here that if $K$ is a number field
then there are two huge commutative topological $K$-algebras called the \emph{finite adeles}
$\A_K^\infty$ and the \emph{adeles} $\A_K$ of $K$, and that they're both locally compact
as topological spaces. We also know from theorem~\ref{NumberField.AdeleRing.baseChangeEquiv}
that $\A_K\cong K\otimes_{\Q}\A_{\Q}K$ (both topologically and algebraically), meaning
that if $R$ is a $K$-algebra then $R_{\A} := R\otimes_K\A_K$ is naturally isomorphic
to $R\otimes_{\Q}\A_{\Q}$. One can
furthermore check that if $R$ is a finite $K$-algebra then the $\A_K$-module topologies and $\A_{\Q}$-module
topologies on $R_{\A}$ coincide. Indeed, the topology on $\A_K$
is the $\A_{\Q}$-module topology, as
$\A_K=\A_{\Q}\otimes_{\Q}K$ as topological $\A_{\Q}$-algebras, where the right hand side
has the $\A_{\Q}$-module topology by definition. So the claim follows from the
fact that if $A$ is a topological ring, $B$ is a topological $A$-algebra
finite as an $A$-module and with the $A$-module topology, and if
$M$ is a topological $B$-module
(and hence a topological $A$-module), then the $A$-module and $B$-module
topologies on~$M$ coincide (this is {\tt moduleTopology.trans} in the repo,
not yet PRed to mathlib).

Let $K$ be a number field and let $D/K$ be a finite-dimensional central simple $K$-algebra
(later on $D$ will be a division algebra (hence the name) but we do not need this yet).
Then $D_{\A}:=D\otimes_K\A_K$ is an $\A_K$-algebra which
is free of finite rank, and if we give $D_{\A}$ the $\A_K$-module topology then it is
a topological ring (by results in mathlib). Furthermore $D_{\A}$ is free of finite
rank over the locally compact topological ring $\A_K$ and is thus also
locally compact. So by the theory of Haar characters (see Chapter~\ref{Haar_char_project})
there is a canonical character $\delta_{D_{\A}}:D_{\A}^\times\to\R_{>0}$ measuring
how left multiplication by an element of $D_{\A}^\times$ changes the additive Haar
measure on $D_{\A}$. Let $D_{\A}^{(1)}$ denote the kernel of $\delta_{D_{\A}}$,
and give it the subspace topology coming from $D_{\A}^\times$.
Corollary~\ref{NumberField.AdeleRing.units_mem_ringHaarCharacter_ker} from the
Haar character miniproject shows that $D^\times$ (regarded as a subgroup of $D_{\A}^\times$
via the map $d\mapsto d\otimes 1$) is contained within $D_{\A}^{(1)}$,
thus the below theorem typechecks.

\begin{theorem}
  \label{NumberField.AdeleRing.DivisionAlgebra.compact_quotient}
  \lean{NumberField.AdeleRing.DivisionAlgebra.compact_quotient}
  \leanok
  If $D$ is a division algebra then
  the quotient $D^\times\backslash D_{\A}^{(1)}$
  with its quotient topology coming from $D_{\A}^{(1)}$, is compact.
\end{theorem}

The rest of this miniproject is devoted to a proof of this theorem.

\section{The proof}

We prove the theorem via a series of lemmas.

\begin{lemma}
  \label{NumberField.AdeleRing.DivisionAlgebra.Aux.existsE}
  \lean{NumberField.AdeleRing.DivisionAlgebra.Aux.existsE}
  \leanok
  There's a compact subset $E$ of $D_{\A}$
  with the property that for all $x\in D_{\A}^{(1)}$,
  the obvious map $xE\to D\backslash D_{\A}$ is not injective.
\end{lemma}

\begin{proof}
  \leanok
  We know that if we pick a $\Q$-basis for $D$
  of size $d$ then this identifies $D$ with $\Q^d$,
  $D_{\A}$ with $\A_{\Q}^d$, and $D\backslash D_{\A}$ with
  $(\Q\backslash\A_{\Q})^d$. Now $\Q$ is discrete in $\A_{\Q}$
  by theorem~\ref{NumberField.AdeleRing.discrete}, and the quotient
  $\Q\backslash \A_{\Q}$ is compact by theorem~\ref{Rat.AdeleRing.cocompact}.
  Hence $D$ is discrete in $D_{\A}$
  and the quotient $D\backslash D_{\A}$ is compact.

    Fix a Haar measure $\mu$ on $D_{\A}$ and push it forward
    to $D\backslash D_{\A}$; by compactness this quotient has finite
    and positive measure, say $m\in\R_{>0}$.
    Choose any compact $E\subseteq D_{\A}$ with measure $> m$
    (for example, choose a $\Z$-lattice $L\cong\Z^d$ in $D\cong\Q^d$,
    define $E_f:=\prod_p L_p\in D\otimes_{\Q}\A_{\Q}^\infty$,
    and define $E_{\infty}\subseteq D\otimes_{\Q}\R\cong\R^n$ to be a huge closed
    ball, large enough to ensure the measure of $E:=E_f\times E_{\infty}$ is bigger than $m$).
    Then $\mu(xE)=\mu(E)>m$ so the map can't be injective.
\end{proof}

\begin{definition}
  \label{NumberField.AdeleRing.DivisionAlgebra.Aux.E}
  \lean{NumberField.AdeleRing.DivisionAlgebra.Aux.E}
  \uses{NumberField.AdeleRing.DivisionAlgebra.Aux.existsE}
  \leanok
  We let $E$ denote any compact set satisfying the hypothesis of the previous lemma.
\end{definition}

\begin{definition}
  \label{NumberField.AdeleRing.DivisionAlgebra.Aux.X}
  \lean{NumberField.AdeleRing.DivisionAlgebra.Aux.X}
  \uses{NumberField.AdeleRing.DivisionAlgebra.Aux.E}
  \leanok
Define $X:=E-E:=\{e-f:e,f\in E\}\subseteq D_{\A}$.
\end{definition}

\begin{definition}
  \label{NumberField.AdeleRing.DivisionAlgebra.Aux.Y}
  \lean{NumberField.AdeleRing.DivisionAlgebra.Aux.Y}
  \uses{NumberField.AdeleRing.DivisionAlgebra.Aux.X}
  \leanok
Define $Y:=X.X:=\{xy:x,y\in X\}\subseteq D_{\A}$.
\end{definition}

\begin{lemma}
  \label{NumberField.AdeleRing.DivisionAlgebra.Aux.X_compact}
  \lean{NumberField.AdeleRing.DivisionAlgebra.Aux.X_compact}
  \uses{NumberField.AdeleRing.DivisionAlgebra.Aux.X}
  \leanok
  $X$ is a compact subset of $D_{\A}$.
\end{lemma}
\begin{proof}
  \leanok
  It's the continuous image of the compact set~$E\times E$.
\end{proof}

\begin{lemma}
  \label{NumberField.AdeleRing.DivisionAlgebra.Aux.Y_compact}
  \lean{NumberField.AdeleRing.DivisionAlgebra.Aux.Y_compact}
  \uses{NumberField.AdeleRing.DivisionAlgebra.Aux.Y}
  \leanok
  $Y$ is a compact subset of $D_{\A}$.
\end{lemma}
\begin{proof}
  \leanok
  It's the continuous image of the compact set~$X\times X$.
\end{proof}

\begin{lemma}
  \label{NumberField.AdeleRing.DivisionAlgebra.Aux.X_meets_kernel}
  \lean{NumberField.AdeleRing.DivisionAlgebra.Aux.X_meets_kernel}
  \uses{NumberField.AdeleRing.DivisionAlgebra.Aux.X, NumberField.AdeleRing.DivisionAlgebra.Aux.E}
  % TODO I used E in the statement of the lemma here
  \leanok
  If $\beta\in D_{\A}^{(1)}$ then
$\beta X\cap D^\times\not=\emptyset$.
\end{lemma}
\begin{proof}
  \leanok
Indeed by lemma~\ref{NumberField.AdeleRing.DivisionAlgebra.Aux.existsE},
the map $\beta E\to D\backslash D_{\A}$
isn't injective, so there are distinct
$\beta e_1,\beta e_2\in \beta E$ with $e_i\in E$ and
$\beta e_1-\beta e_2=b\in D$.
Now $b\not=0$ and $D$ is a division algebra, so $b\in D^\times$.
And $e_1-e_2\in X$ so $b=\beta(e_1-e_2)\in \beta X$, so we're done.
\end{proof}

\begin{lemma}
  \label{NumberField.AdeleRing.DivisionAlgebra.Aux.X_meets_kernel'}
  \lean{NumberField.AdeleRing.DivisionAlgebra.Aux.X_meets_kernel'}
  \uses{NumberField.AdeleRing.DivisionAlgebra.Aux.X}
  \leanok
  Similarly, if $\beta\in D_{\A}^{(1)}$ then
  $X\beta^{-1}\cap D^\times\not=\emptyset$.
\end{lemma}
\begin{proof}
  \leanok
  % TODO does this make the lines dottted? I used E in the proof here
  % TODO does \uses a theorem from another project mess up the colouring?
  \uses{NumberField.AdeleRing.DivisionAlgebra.Aux.E, addHaarScalarFactor.left_mul_eq_right_mul}
  Indeed, $\beta^{-1}\in D_{\A}^{(1)}$, and so left multiplication by $\beta^{-1}$
  doesn't change Haar measure on $D_{\A}$, so neither does right multiplication
  (by theorem~\ref{addHaarScalarFactor.left_mul_eq_right_mul}).
  So the same argument works: $E\beta^{-1}\to D\backslash D_{\A}$ is not
  injective so choose $e_1\beta^{-1}\not=e_2\beta^{-1}$ with difference $b\in D$
  and then $(e_1-e_2)\beta^{-1}\in D-{0}=D^\times$.
\end{proof}

\begin{definition}
  \label{NumberField.AdeleRing.DivisionAlgebra.Aux.T}
  \lean{NumberField.AdeleRing.DivisionAlgebra.Aux.T}
  \uses{NumberField.AdeleRing.DivisionAlgebra.Aux.Y}
  \leanok
  Let $T:=Y\cap D^\times$.
\end{definition}

\begin{lemma}
  \label{NumberField.AdeleRing.DivisionAlgebra.Aux.T_finite}
  \lean{NumberField.AdeleRing.DivisionAlgebra.Aux.T_finite}
  \uses{NumberField.AdeleRing.DivisionAlgebra.Aux.T}
  \leanok
  $T$ is finite.
\end{lemma}
\begin{proof}
  \uses{NumberField.AdeleRing.DivisionAlgebra.Aux.Y_compact}
  \leanok
  % TODO I used T in both the theorem and proof this time!
  It suffices to prove that $Y\cap D$ is finite.
    But $D\subseteq D_{\A}$ is a discrete additive subgroup, and hence closed.
    And $Y\subseteq D_{\A}$ is compact.
    So $D\cap Y$ is compact and discrete, so finite.
\end{proof}

\begin{definition}
  \label{NumberField.AdeleRing.DivisionAlgebra.Aux.C}
  \lean{NumberField.AdeleRing.DivisionAlgebra.Aux.C}
  \uses{NumberField.AdeleRing.DivisionAlgebra.Aux.T,
    NumberField.AdeleRing.DivisionAlgebra.Aux.X}
  \leanok
  Define $C:= (T^{-1}.X) \times X\subset D_{\A}\times D_{\A}$.
\end{definition}

\begin{lemma}
  \label{NumberField.AdeleRing.DivisionAlgebra.Aux.C_compact}
  \lean{NumberField.AdeleRing.DivisionAlgebra.Aux.C_compact}
  \uses{NumberField.AdeleRing.DivisionAlgebra.Aux.C,
    NumberField.AdeleRing.DivisionAlgebra.Aux.T_finite,
    NumberField.AdeleRing.DivisionAlgebra.Aux.X_compact}
  \leanok
  $C$ is compact.
\end{lemma}
\begin{proof}
  \leanok
   $X$ is compact and $T$ is finite.
\end{proof}

\begin{lemma}
  \label{NumberField.AdeleRing.DivisionAlgebra.Aux.antidiag_mem_C}
  \lean{NumberField.AdeleRing.DivisionAlgebra.Aux.antidiag_mem_C}
  \uses{NumberField.AdeleRing.DivisionAlgebra.Aux.C}
  \leanok
  For every $\beta\in D_{\A}^{(1)}$, there exists $b\in D^\times$
  and $\nu\in D_{\A}^{(1)}$ such that $\beta=b\nu$ and $(\nu,\nu^{-1})\in C.$
\end{lemma}
\begin{proof}
  \leanok
  \uses{NumberField.AdeleRing.DivisionAlgebra.Aux.E,
    NumberField.AdeleRing.DivisionAlgebra.Aux.X_meets_kernel,
    NumberField.AdeleRing.DivisionAlgebra.Aux.X_meets_kernel',
    NumberField.AdeleRing.DivisionAlgebra.Aux.C}
  By lemma~\ref{NumberField.AdeleRing.DivisionAlgebra.Aux.X_meets_kernel},
  $\beta X\cap D^\times\not=\emptyset$,
  and lemma~\ref{NumberField.AdeleRing.DivisionAlgebra.Aux.X_meets_kernel'},
  $X\beta^{-1}\cap D^\times\not=\emptyset$,
  so we can write $\beta x_1=b_1$ and $x_2\beta^{-1}=b_2$ with $b_i\in D^\times$ and $x_i\in X$.
  Note that $\beta\in D_{\A}^{(1)}$ and $b_i\in D^{\times}\subseteq D_{\A}^{(1)}$ by
  corollary~\ref{NumberField.AdeleRing.units_mem_ringHaarCharacter_ker}, so $x_i\in D_{\A}^{(1)}$
  as well. In particular $x_i\in D_{\A}^\times$ so $x_1^{-1}$ makes sense.

  Multiplying the equations defining the $x_i$ and $b_i$ we deduce that
  $x_2x_1=b_2b_1\in Y\cap D^\times=T$ (recall that $Y=X.X$ and $T=Y\cap D^\times$
  is finite); call this element $t$. Then $x_1^{-1}=t^{-1}x_2\in T^{-1}.X$,
  and $x_1\in X$, so if we set $\nu=x_1^{-1}\in D_{\A}^{(1)}$
  and $b=b_1\in D^\times$ then we have $\beta=b\nu$ and $(\nu,\nu^{-1})\in C := (T^{-1}.X)\times X$.
  We are done!
\end{proof}

We can now prove Fujisaki's theorem~\ref{NumberField.AdeleRing.DivisionAlgebra.compact_quotient}.

\begin{proof}
  \proves{NumberField.AdeleRing.DivisionAlgebra.compact_quotient}
  \leanok
  \uses{MeasureTheory.ringHaarChar_continuous,
    NumberField.AdeleRing.DivisionAlgebra.Aux.antidiag_mem_C,
    NumberField.AdeleRing.DivisionAlgebra.Aux.C_compact}
  Indeed, if $M$ is the preimage of $C$ under the inclusion $D_{\A}^{(1)} \to D_{\A}\times D_{\A}$
  sending $\nu$ to $(\nu,\nu^{-1})$, then $M$ is a closed subspace
    of a compact
  space so it's compact (note that $\delta_{D_{\A}}$ is continuous,
  by theorem~\ref{MeasureTheory.ringHaarChar_continuous}, so $D_{\A}^{(1)}$ is a closed subset of
  $D_{\A}^\times$ which is itself a closed subset of $D_{\A}\times D_{\A}$).
  Lemma~\ref{NumberField.AdeleRing.DivisionAlgebra.Aux.antidiag_mem_C} shows that $M$ surjects onto
  $D^\times\backslash D_{\A}^{(1)}$ which is thus also compact.
\end{proof}

We note here some useful consequences.

\begin{theorem}
  \label{NumberField.FiniteAdeleRing.DivisionAlgebra.units_cocompact}
  \lean{NumberField.FiniteAdeleRing.DivisionAlgebra.units_cocompact}
  \leanok
  $D^\times\backslash(D\otimes_K\A_K^\infty)^\times$ is compact.
\end{theorem}
\begin{proof}
  \uses{NumberField.AdeleRing.DivisionAlgebra.compact_quotient}
  \leanok
There's a natural map $\alpha$ from $D^\times\backslash D_{\A}^{(1)}$ to
  $D^\times\backslash (D\otimes_K \A_K^\infty)^\times$. We claim that it's
  surjective. Granted this claim, we are home, because if we put the quotient
  topology on $D^\times\backslash (D\otimes_K \A_K^\infty)^\times$ coming from
  $(D\otimes_K \A_K^\infty)^\times$ then it's readily verified that $\alpha$
  is continuous, and the continuous image of a compact space is compact.

  As for surjectivity: say $x\in (D\otimes_K \A_K^\infty)^\times$. We need to extend
  $x$ to an element $(x,y)\in (D\otimes_K \A_K^\infty)^\times\times(D\otimes_K K_\infty)^\times$
  which is in the kernel of $\delta_{D_{\A}}$. Because $\delta_{D_{\A}}(x,1)$ is some positive
  real number, it will suffice to show that if $r$ is any positive real number then we can
  find $y\in (D\otimes_K \A_K^\infty)^\times=(D\otimes_{\Q}\R)^\times$ with $\delta_{D_{\A}}(1,y)=r$,
  or equivalently (setting $D_{\R}=D\otimes_{\Q}\R$) that $\delta_{D_{\R}}(y)=r$.
  But $D\not=0$ as it is a division algebra,and hence $\Q\subseteq D$, meaning
  $\R\subseteq D_{\R}$, and if
  $x\in\R^\times\subseteq D_{\R}^\times$ then $\delta(x)=|x|^d$ with $d=\dim_{\Q}(D)$,
  as multiplication by $x$ is just scaling by a factor of $x$ on $D_{\R}\cong\R^d$.
  In particular we can set $x=y^{1/d}$.
\end{proof}
\begin{remark} In this generality the quotient might not be Hausdorff.
\end{remark}

\begin{theorem}
  \label{NumberField.FiniteAdeleRing.DivisionAlgebra.finiteDoubleCoset}
  \lean{NumberField.FiniteAdeleRing.DivisionAlgebra.finiteDoubleCoset}
  \uses{NumberField.FiniteAdeleRing.DivisionAlgebra.units_cocompact}
  \leanok
  If $U$ is an open subgroup of $(D\otimes_K \A_K^\infty)^\times$
  then the double coset space $D^\times\backslash(D\otimes_K \A_K^\infty)^\times/U$ is finite.
\end{theorem}
\begin{proof}
  \leanok
  The double cosets give a disjoint open cover of $(D\otimes_K \A_K^\infty)$
  which descends to a disjoint open cover of the quotient space
  $D^\times\backslash(D\otimes_K \A_K^\infty)^\times$. However this space is compact
  by theorem~\ref{NumberField.FiniteAdeleRing.DivisionAlgebra.units_cocompact}.
\end{proof}


%%FILE blueprint/src/chapter/HaarCharacterProject.tex
\chapter{Miniproject: Haar Characters}\label{Haar_char_project}

\section{The goal}

The goal of this miniproject is to develop the theory (i.e., the basic API) of Haar characters.
``Haar character'' is a name I've made up to describe a certain character of the units of a locally
compact topological ring. The main result we need here is that if $B$ is a finite-dimensional
algebra over a number field~$K$, then $B^\times$ is in the kernel of the Haar character
of $B\otimes_K\A_K$, where $\A_K$ is the ring of adeles of~$K$. Most if not all of this
should probably be in mathlib.

KMB would like to heartily thank S\'ebastien Gou\"ezel for the help he gave during the preparation
of this material.

\section{Initial definitions}

\subsection{Scaling Haar measure on a group}

Let $A$ be a locally compact topological additive abelian group. There's then a regular additive
Haar measure $\mu$ on $A$, unique up to a positive scalar factor. If $\phi:(A,+)\cong(A,+)$ is a
homeomorphism and an additive automorphism of $A$, then we can push forward $\mu$
along $\phi$ to get a second measure $\phi_*\mu$ on $A$, with the property that
$(\phi_*\mu)(X)=\mu(\phi^{-1}X)$ for any Borel subset $X$ of $A$.

Now $\phi_*\mu$ is a translation-invariant and regular measure,
and hence also a Haar measure on $A.$ It must thus differ from
$\mu$ by a positive scalar factor, which we call $d_A(\phi)$.
There is a choice of normalization here between $d_A(\phi)$
and $d_A(\phi)^{-1}$, so let us be more precise.

\begin{definition}
  \label{MeasureTheory.addEquivAddHaarChar}
  \lean{MeasureTheory.addEquivAddHaarChar}
  \leanok
  If $A$ is a locally compact topological additive abelian group,
  if $\mu$ is a regular Haar measure on $A$, and if $\phi:A\to A$ is an
  additive homeomorphism, then we let $d_A(\phi)$ denote the unique positive
  real number such that $\mu(X)=d_A(\phi)(\phi_*\mu)(X)$ for any Borel set~$X$.
\end{definition}

To give an example, if $\phi$ is multiplication by $2$ on the real numbers,
if $X=[0,1]$, and if $\mu$ is Lebesgue measure on the Borel subsets of $\R$,
we have that $\phi_*\mu(X)=\mu(\phi^{-1}(X))=\mu([0,1/2])=1/2$,
so $1=d_A(\phi)/2$ meaning that $d_A(\phi)=2$. Similarly if $\phi$ is multiplication
by $-2$ and $X=[0,1]$ then $\phi^{-1}(X)=[-1/2,0]$ which again has measure $1/2$,
so $d_A(\phi)$ is 2 again.

Strictly speaking our definition of $d_A(\phi)$ depends on the choice of regular Haar
measure $\mu$. Note that {\tt mathlib} offers a fixed Borel regular Haar measure
{\tt MeasureTheory.Measure.haar} on any locally compact topological group
and the actual definition of $d_A$ in the code uses this definition.
Note also that the code defines everything for multiplicative groups
and uses {\tt @[to\_additive]} to deduce the corresponding results
for additive groups.

Here are some basic results about this construction. In all of them,
$A$ is a locally compact topological group and $\phi:A\to A$ is
a group isomorphism and a homeomorphism. The first lemma shows that
the definition of $d_A(\phi)$ is indeed independent of the choice of Haar measure.

\begin{lemma}
  \label{MeasureTheory.addEquivAddHaarChar_eq}
  \lean{MeasureTheory.addEquivAddHaarChar_eq}
  \uses{MeasureTheory.addEquivAddHaarChar}
  \discussion{508}
  \leanok
  $d_A(\phi)$ is independent of choice of regular Haar measure.
\end{lemma}
\begin{proof}
  \leanok
  If $\mu'$ is a second choice then $\mu'=\lambda\mu$ for some
  positive real $\lambda$, and the $\lambda$s on each side of
  $\mu'(X)=d_A(\phi)(\phi_*\mu')(X)$ cancel.
\end{proof}

\begin{lemma}
  \lean{MeasureTheory.addEquivAddHaarChar_smul_map}
  \label{MeasureTheory.addEquivAddHaarChar_smul_map}
  \uses{MeasureTheory.addEquivAddHaarChar_eq}
  \leanok
  If $\mu$ is any regular Haar measure on $A$ then
  $d_A(\phi)(\phi_*\mu) = \mu.$
\end{lemma}
\begin{proof}
  \leanok
  This is a restatement of the previous result.
\end{proof}

We can of course also pull a Haar measure $\mu$ back along a homeomorphism $\phi$,
giving a measure $\phi^*\mu$ such that $\phi^*\mu(X)=\mu(\phi(X))$.

\begin{corollary}
  \lean{MeasureTheory.addEquivAddHaarChar_smul_eq_comap}
  \label{MeasureTheory.addEquivAddHaarChar_comap}
  \uses{MeasureTheory.addEquivAddHaarChar_smul_map}
  \leanok
  If $\mu$ is any regular Haar measure on $A$ then
  $d_A(\phi)\mu = \phi^*\mu.$
\end{corollary}
\begin{proof}
  \leanok
  This follows from lemma~\ref{MeasureTheory.addEquivAddHaarChar_smul_map}
  applied to the regular Haar measure $\phi^*\mu$ and the fact that $\phi_*\phi^*\mu=\mu$.
\end{proof}

Now let $\mu$ be any regular additive Haar measure on $A$.

\begin{lemma}
  \label{MeasureTheory.addEquivAddHaarChar_smul_preimage}
  \lean{MeasureTheory.addEquivAddHaarChar_smul_preimage}
  \uses{MeasureTheory.addEquivAddHaarChar}
  \discussion{509}
  \leanok
  If $X$ is a Borel set then $\mu(X)=d_A(\phi)\mu(\phi^{-1}X)$.
\end{lemma}
\begin{proof}
  \leanok
  \uses{MeasureTheory.addEquivAddHaarChar_smul_map}
  This follows immediately from lemma~\ref{MeasureTheory.addEquivAddHaarChar_smul_map}
  and the definition of the pushforward of a measure.
\end{proof}

\begin{lemma}
  \label{MeasureTheory.addEquivAddHaarChar_smul_integral_map}
  \lean{MeasureTheory.addEquivAddHaarChar_smul_integral_map}
  \uses{MeasureTheory.addEquivAddHaarChar}
  \discussion{510}
  \leanok
  If $f:A\to\R$ is a Borel measurable function then
  $d_A(\phi)\int f(x)d\phi_*\mu(x)=\int f(x)d\mu(x)$.
\end{lemma}
\begin{proof}
  \leanok
  \uses{MeasureTheory.addEquivAddHaarChar_smul_map}
  This also follows immediately from lemma~\ref{MeasureTheory.addEquivAddHaarChar_smul_map}.
\end{proof}

We also have the following variant:

\begin{lemma}
  \label{MeasureTheory.addEquivAddHaarChar_smul_integral_comap}
  \lean{MeasureTheory.integral_comap_eq_addEquivAddHaarChar_smul}
  \uses{MeasureTheory.addEquivAddHaarChar_comap}
  \leanok
  If $f:A\to\R$ is a Borel measurable function then
  $d_A(\phi)\int f(x)d\mu(x)=\int f(x)d\phi^*\mu(x)$.
\end{lemma}
\begin{proof}
  \leanok
  This is immediate from corollary~\ref{MeasureTheory.addEquivAddHaarChar_comap}.
\end{proof}

Note that as a consequence of lemma~\ref{MeasureTheory.addEquivAddHaarChar_smul_preimage},
if $X$ is a Borel subset of $A$ with positive finite measure then we can read
off $d_A(\phi)$ by $d_A(\phi)=\mu(X)/\mu(\phi^{-1}(X))$, and hence also by
$d_A(\phi)=\mu(\phi(X))/\mu(X)$. A nice special case is when
$\mu(X)=1$, in which case we have $d_A(\phi)=\mu(\phi(X))$ for all $\phi$,
or $d_A(\phi)=1/\mu(\phi^{-1}(X))$.
Similarly, by lemma~\ref{MeasureTheory.addEquivAddHaarChar_smul_integral_map},
if $f$ is a measurable function with $0<\int f(x)d\mu(x)<\infty$ then
we can read off $d_A(\phi)$ by $d_A(\phi)=(\int f(x)d\mu(x))/(\int f(x)\phi_*\mu(x))$.
Note that {\tt mathlib} supplies such $f$ with the function {\tt exists\_continuous\_nonneg\_pos}.
The following are also straightforward.

\begin{lemma}
  \label{MeasureTheory.addEquivAddHaarChar_refl}
  \lean{MeasureTheory.addEquivAddHaarChar_refl}
  \uses{MeasureTheory.addEquivAddHaarChar}
  \leanok
  $d_A(id)=1.$
\end{lemma}
\begin{proof}
  \uses{MeasureTheory.addEquivAddHaarChar_eq}
  \leanok
  Clear.
\end{proof}

\begin{lemma}
  \label{MeasureTheory.addEquivAddHaarChar_trans}
  \lean{MeasureTheory.addEquivAddHaarChar_trans}
  \uses{MeasureTheory.addEquivAddHaarChar}
  \leanok
  \discussion{511}
  $d_A(\phi\circ\psi)=d_A(\phi)d_A(\psi).$
\end{lemma}
\begin{proof}
  \leanok
  \uses{MeasureTheory.addEquivAddHaarChar_smul_preimage}
  Here's one way: it suffices to prove that
  $d_A(\phi\circ\psi)(\phi\circ\psi)_*\mu=d_A(\phi)d_A(\psi)(\phi\circ\psi)_*\mu$
  (because there exists a compact set with positive finite measure)
  and using lemma~\ref{MeasureTheory.addEquivAddHaarChar_smul_map}
  and the fact that $(\phi\circ\psi)_*\mu=\phi_*(\psi_*\mu)$
  one can simplify both sides to $\mu$.
\end{proof}

\subsection{Scaling Haar measure on a ring}

Now let $R$ be a locally compact topological ring. The \emph{Haar character} of $R$,
or more precisely the \emph{left Haar character} of $R$, is a group homomorphism
$R^\times\to\R^\times$ defined in the following way. If $u\in R^\times$ then left multiplication
by $u$, namely the map $\ell_u:(R,+)\to(R,+)$ defined by $\ell_u(r)=ur$, is a homeomorphism and
an additive automorphism of $(R,+)$, so the preceding theory applies to $\ell_u$.

\begin{definition}
  \label{MeasureTheory.ringHaarChar}
  \lean{MeasureTheory.ringHaarChar}
  \uses{MeasureTheory.addEquivAddHaarChar,MeasureTheory.addEquivAddHaarChar_refl,
  MeasureTheory.addEquivAddHaarChar_trans}
  \leanok
  We define $\delta_R(u)$ (or just $\delta(u)$ when the ring $R$ is clear) to be $d_R(\ell_u)$.
\end{definition}

Lemmas~\ref{MeasureTheory.addEquivAddHaarChar_refl}
and~\ref{MeasureTheory.addEquivAddHaarChar_trans} immediately imply that $\delta_R$ is a group
homomorphism from $R^\times$ to $\R_{>0}$.
Also immediate from previous lemmas is

\begin{lemma}
  \label{MeasureTheory.ringHaarChar_mul_integral}
  \lean{MeasureTheory.ringHaarChar_mul_integral}
  \uses{MeasureTheory.ringHaarChar}
  \leanok
  \discussion{514}
  If $f:R\to\R$ is a Borel measurable function and $u\in R^\times$ then
  $\delta_R(u)\int f(ux)d\mu(x)=\int f(x)d\mu(x)$.
\end{lemma}
\begin{proof}
  \uses{MeasureTheory.addEquivAddHaarChar_smul_integral_map}
  \leanok
  A short calculation using lemma~\ref{MeasureTheory.addEquivAddHaarChar_smul_integral_map}.
\end{proof}

\begin{lemma}
  \label{MeasureTheory.ringHaarChar_mul_volume}
  \lean{MeasureTheory.ringHaarChar_mul_volume}
  \uses{MeasureTheory.ringHaarChar,MeasureTheory.addEquivAddHaarChar_smul_preimage}
  \leanok
  \discussion{515}
  If $X$ is a Borel subset of $R$ and $r\in R^\times$ then $\mu(rX)=\delta_R(r)\mu(X)$.
\end{lemma}
\begin{proof}
  \leanok
  \uses{MeasureTheory.addEquivAddHaarChar_smul_preimage}
   Immediate from lemma~\ref{MeasureTheory.addEquivAddHaarChar_smul_preimage}.
\end{proof}

The next result lies a little deeper.

\begin{corollary}
  \label{MeasureTheory.ringHaarChar_continuous}
  \lean{MeasureTheory.ringHaarChar_continuous}
  \uses{MeasureTheory.ringHaarChar}
  \leanok
  \discussion{516}
  The function $\delta_R:R^\times\to\R_{>0}$ is continuous.
\end{corollary}
\begin{proof}
  \leanok
  \uses{MeasureTheory.ringHaarChar_mul_integral}
  Fix a Haar measure $\mu$ on $R$ and a continuous real-valued function
  $f$ on $R$ with compact support and such that $\int f(x) d\mu(x)\not=0$.
  Then $r \mapsto \int f(rx) d\mu(x)$ is a continuous function
  from $R\to\R$ (because a continuous function with compact support is uniformly
   continuous) and thus gives a continuous function $R^\times\to\R$.
   Thus the function $u\mapsto (\int f(ux) d\mu(x))/(\int f(x)d\mu(x))$ is
   a continuous function from $R^\times$ to $\R$ taking values in $\R_{>0}$.
   Hence $\delta_R^{-1}$ is continuous,
   from lemma~\ref{MeasureTheory.ringHaarChar_mul_integral},
   and thus $\delta_R$ is too.
\end{proof}

\section{Examples}

We discuss some examples of Haar characters.

\begin{lemma}
  \label{MeasureTheory.ringHaarChar_real}
  \lean{MeasureTheory.ringHaarChar_real}
  If $R=\R$ then $\delta_R(u)=|u|$.
  \leanok
\end{lemma}
\begin{proof}
  \uses{MeasureTheory.ringHaarChar_mul_volume}
  \leanok
  Take $\mu$ to be Lebesgue measure and $X=[0,1]$.
We have $\delta(u)=\mu(uX)$. If $u>0$ then $u[0,1]=[0,u]$ which has measure $u=|u|$,
and if $u<0$ then $u*[0,1]=[u,0]$ which has measure $-u=|u|$.
\end{proof}

\begin{lemma}
  \label{MeasureTheory.ringHaarChar_complex}
  \lean{MeasureTheory.ringHaarChar_complex}
  \leanok
  If $R=\bbC$ then $\delta_R(u)=|u|^2$.
\end{lemma}
\begin{proof}
  \uses{MeasureTheory.ringHaarChar_mul_volume}
  \leanok
  Multiplication by a positive real $r$ sends a unit square to a square of area $r^2=|r|^2$.
  Multiplication by $e^{i\theta}$ is a rotation and thus does not change area.
  The general case follows.
\end{proof}

\begin{lemma}
  \label{MeasureTheory.ringHaarChar_padic}
  \lean{MeasureTheory.ringHaarChar_padic}
  \leanok
  If $R=\Q_p$ then $\delta_R(u)=|u|_p$, the usual $p$-adic norm.
\end{lemma}
\begin{proof}
  \uses{MeasureTheory.ringHaarChar_mul_volume}
  \leanok
  Normalise Haar measure so that $\mu(\Z_p)=1$.
  If $u$ is a $p$-adic unit then $u\Z_p=\Z_p$ so multiplication by $u$ didn't change
  Haar measure. If however $u=p$ then $u\Z_p$ has index $p$ in $\Z_p$ and, because
  $\mu(i+p\Z_p)=\mu(p\Z_p)$ we have that $\mu(\Z_p)=p\mu(p\Z_p)$ and thus $\delta(p)=p^{-1}$.
  These elements generate $\Q_p^\times$ and two characters which agree on generators
  of a group must agree on the group.
\end{proof}

\begin{remark}
If $R$ is a finite extension of $\Q_p$ then $\delta_R(u)$
is the norm on $R$ normalised in the following way:
$\delta_R(\varpi)=q^{-1}$, where $\varpi$ is a uniformiser
and $q$ is the size of the (finite) residue field. In fact the same is true
for any nonarchimedean local field. The proof is
the same as for $\Q_p$. Right now this is difficult to state in Lean because
there is still some discussion about the definition of a nonarchimedean local field.
\end{remark}

\section{Algebras}

  Say $F$ is a locally compact topological ring (for example $\R$ or $\bbC$ or $\Q_p$,
  or the adeles of a number field), $V$
  is a finite free $F$-module, and $\phi:V\to V$ is an invertible $F$-linear map.
  Then $V$ with its module topology (which is the product topology if one picks a basis)
  is a locally compact topological abelian group, and $\phi$ is additive.
  One can check that linearity implies continuity (this is
  {\tt IsModuleTopology.continuous\_of\_linearMap} in mathlib),
  so in fact $\phi$ is a homeomorphism
  and our theory applies. The following lemma gives a formula for the scale factor $d_V(\phi)$.

\begin{lemma}
  \label{MeasureTheory.addEquivAddHaarChar_eq_ringHaarChar_det,
    MeasureTheory.addEquivAddHaarChar_eq_ringHaarChar_det_of_existsListTransvecEtc}
  \lean{MeasureTheory.addEquivAddHaarChar_eq_ringHaarChar_det,
    MeasureTheory.addEquivAddHaarChar_eq_ringHaarChar_det_of_existsListTransvecEtc}
  \uses{MeasureTheory.ringHaarChar}
  \leanok
  \discussion{517}
  Assume that there's an $F$-basis for $V$ such that $\phi$ is a product of elementary
  and diagonal matrices (note that this is automatic if $F$ is a field and mathlib
  has this). Then
  $d_V(\phi)=\delta_F(\det(\phi))$, where $\det(\phi)\in F$ is the determinant of
  $\phi$ as an $F$-linear map.
\end{lemma}
\begin{proof}
\uses{MeasureTheory.addEquivAddHaarChar}
\leanok
The proof is a generalization of
\href{https://leanprover-community.github.io/mathlib4\_docs/Mathlib/MeasureTheory/Measure/Lebesgue/Basic.html\#Real.map\_matrix\_volume\_pi\_eq\_smul\_volume\_pi}
{\tt Real.map\_matrix\_volume\_pi\_eq\_smul\_volume\_pi},
which crucially uses the induction principle
\href{https://leanprover-community.github.io/mathlib4\_docs/Mathlib/LinearAlgebra/Matrix/Transvection.html\#Matrix.diagonal\_transvection\_induction\_of\_det\_ne\_zero}
{\tt Matrix.diagonal\_transvection\_induction\_of\_det\_ne\_zero}.
One checks it explicitly for diagonal matrices and for matrices which are the identity
except that one off-diagonal entry is non-zero.

Note: we assume that $F$ is second countable
(but it shouldn't be necessary).
\end{proof}

Now say $F$ is a locally compact topological field, and that $R$ is a (possibly
non-commutative) $F$-algebra. Recall that this means that ($R=0$ or) $F$ lies in the centre of $R$.
Assume that $R$ is finite-dimensional as an $F$-vector space. Then if we give $R$ the
$F$-module topology (which is just the product topology if we pick a basis) then it is known
that $R$ becomes a topological ring. Now say $u\in R^\times$, and
recall $\ell_u:R\to R$ is left multiplication by $u$. Then $\ell_u$ is easily checked to be
an $F$-linear homeomorphism.

  \begin{corollary}
  \label{MeasureTheory.algebra_ringHaarChar_eq_ringHaarChar_det}
  \lean{MeasureTheory.algebra_ringHaarChar_eq_ringHaarChar_det}
  \leanok
  If $u\in R^\times$ then $\delta_R(u)=\delta_F(\det(\ell_u))$.
\end{corollary}
\begin{proof}
  \uses{MeasureTheory.addEquivAddHaarChar_eq_ringHaarChar_det}
  \leanok
  Follows immediately from the preceding lemma.
\end{proof}

\section{Left and right multiplication}

If $R$ is a locally compact topological ring, and if multiplication on $R$ is not commutative,
then left and right multiplication by an element of~$R$ can scale Haar measure in different ways.
For example if $R$ is the upper-triangular $2\times 2$ matrices with real
entries, then left multiplication by $\begin{pmatrix}a&0\\0&1\end{pmatrix}$
sends $\begin{pmatrix}x&y\\0&z\end{pmatrix}$ to $\begin{pmatrix}ax&ay\\0&z\end{pmatrix}$
and thus scales $R$'s additive
Haar measure by $|a|^2$, but right multiplication by $\begin{pmatrix}a&0\\0&1\end{pmatrix}$
sends $\begin{pmatrix}x&y\\0&z\end{pmatrix}$ to $\begin{pmatrix}ax&y\\0&z\end{pmatrix}$
and thus scales $R$'s additive Haar measure by a factor of $|a|$.

What's going on here is that if we regard left and right multiplication as $\R$-linear
maps from $R$ to $R$, then their associated matrices with respect to the obvious basis
are $diag(a,a,1)$ and $diag(a,1,1)$, which have different determinants.

However, if $k$ is now any field and if $B$ is a finite-dimensional central
simple algebra over $k$ (for example a quaternion algebra, the case we'll care about later),
and if $u\in B^\times$ then $x\mapsto ux$ and $x\mapsto xu$
are both $k$-linear endomorphisms of $B$, and I claim that they have
the same determinant.

\begin{lemma}
  \label{IsSimpleRing.mulLeft_det_eq_mulRight_det}
  \lean{IsSimpleRing.mulLeft_det_eq_mulRight_det}
  \leanok
  \discussion{518}
  Say $B$ is a finite-dimensional central simple algebra over a field~$k$,
  and $u\in B^\times$. Let $\ell_u:B\to B$ be the $k$-linear mapping $x$ to $ux$ and
  let $r_u:B\to B$ be the $k$-linear map sending $x$ to $xu$. Then
  $\det(\ell_u)=\det(r_u)$.
\end{lemma}
\begin{proof}
  \leanok
  Determinants are unchanged by base extension, so WLOG $k$ is algebraically closed.
  Then it's known that $B$ must be a matrix algebra, say $M_n(k)$. Now $u$ can be thought
  of as a matrix which has its own intrinsic determinant $d$, and $B$ as a left $B$-module
  becomes a direct sum of $n$ copies of $V$, the standard $n$-dimensional representation of $B$.
  Thus $\det(\ell_u)=d^n$. Similarly $\det(r_u)=d^n$ and in particular they are equal.
\end{proof}

\begin{corollary}
  \label{IsSimpleRing.ringHaarChar_eq_addEquivAddHaarChar_mulRight}
  \lean{IsSimpleRing.ringHaarChar_eq_addEquivAddHaarChar_mulRight}
  \leanok
  If $B$ is a central simple algebra over a locally compact field $F$, and if $u\in B^\times$,
  then $d_B(r_u)=\delta_B(u)$ (recall that the latter is defined to mean $d_B(\ell_u)$).
\end{corollary}
\begin{proof}
  \leanok
  \uses{IsSimpleRing.mulLeft_det_eq_mulRight_det, MeasureTheory.addEquivAddHaarChar_eq_ringHaarChar_det}
  If $\ell_u$ and $r_u$ denote left and right multiplication by $u$ on $B$, then we have
  seen in lemma~\ref{MeasureTheory.addEquivAddHaarChar_eq_ringHaarChar_det} that $d_B(r_u)=\delta_F(\det(r_u))$.
  Lemma~\ref{IsSimpleRing.mulLeft_det_eq_mulRight_det} tells
  us that this is $\delta_F(\det(\ell_u))$ and this is $\delta_B(u)$ again by
  corollary~\ref{MeasureTheory.addEquivAddHaarChar_eq_ringHaarChar_det}.
\end{proof}

\section{Finite Products}

Here are two facts which we will need about products.

\begin{lemma}
  \label{MeasureTheory.addEquivAddHaarChar_prodCongr}
  \lean{MeasureTheory.addEquivAddHaarChar_prodCongr}
  \leanok
  \discussion{520}
  If $(A,+)$ and $(B,+)$ are locally compact topological abelian groups,
  and if $\phi:A\to A$ and $\psi:B\to B$ are additive homeomorphisms,
  then $\phi\times\psi:A\times B\to A\times B$ is an additive homeomorphism (this is
  obvious), and
  $d_{A\times B}(\phi\times\psi)=d_A(\phi)d_B(\psi)$.
\end{lemma}
\begin{proof}
  \leanok
  We only need this result in the case where both $A$ and $B$ are second-countable, in which case
  {\tt Prod.borelSpace} can be used to show that Haar measure on $A\times B$ is the product of
  Haar measures on $A$ and $B$, and in this case the result follows easily. Without this assumption,
  the product of these measures may not even be a Borel measure and one has to be more careful.
  The proof in this case is explained by Gou\"{e}zel
  \href{https://leanprover.zulipchat.com/#narrow/channel/116395-maths/topic/Product.20of.20Borel.20spaces/near/487257981}{here}.
  Here is the idea. Let $\rho$ be a Haar measure on $A\times B$. Fix sets $X\subseteq A$ and
  $Y\subseteq B$ which are compact with nonempty interior. We can now pull back $\rho$
  to a measure $\nu$ on the Borel sigma-algebra of $A$ defined as $\nu(s)=\rho(s\times Y)$
  and this is easily checked to be a Haar measure on $A$. Then
  $\delta_{A\times B}(a,0)\nu(X)=
  \delta_{A\times B}(a,0)\rho(X\times Y)=\rho((a,0)(X\times Y))=
  \rho(aX\times Y)=\nu(aX)=\delta_A(a)\nu(X)$
  , so $\delta_{A\times B}(a,0)=\delta_A(a)$.
  Similarly $\delta_{A\times B}(0,b)=\delta_B(b)$ and because $\delta_{A\times B}$ is a group
  homomorphism we're home.
\end{proof}

\begin{lemma}
  \label{MeasureTheory.addEquivAddHaarChar_piCongrRight}
  \lean{MeasureTheory.addEquivAddHaarChar_piCongrRight}
  \leanok
  \discussion{521}
  \uses{MeasureTheory.addEquivAddHaarChar_prodCongr}
  If $A_i$ are a finite collection of locally compact topological abelian groups,
  with $\phi_i:A_i\to A_i$ additive homeomorphisms, then $d_{\prod_i A_i}(\prod_i\phi_i)=\prod_i d_{A_i}(\phi_i)$.
\end{lemma}
\begin{proof}
  \leanok
  Induction on the size of the finite set, using the previous lemma.
\end{proof}

\begin{lemma}
  \label{MeasureTheory.ringHaarChar_prod}
  \lean{MeasureTheory.ringHaarChar_prod}
  \leanok
  \uses{MeasureTheory.addEquivAddHaarChar_prodCongr}
  If $R$ and $S$ are locally compact topological rings, then $\delta_{R\times S}(r,s)=\delta_R(r)\times\delta_S(s)$.
\end{lemma}
\begin{proof}
  \leanok
  Follows immediately from lemma~\ref{MeasureTheory.addEquivAddHaarChar_prodCongr}.
\end{proof}

\begin{lemma}
  \label{MeasureTheory.ringHaarChar_pi}
  \lean{MeasureTheory.ringHaarChar_pi}
  \leanok
  \uses{MeasureTheory.addEquivAddHaarChar_piCongrRight}
  If $R_i$ are a finite collection of locally compact topological rings,
  and $u_i\in R_i^\times$ then $\delta_{\prod_i R_i}((u_i)_i)=\prod_i\delta_{R_i}(u_i)$.
\end{lemma}
\begin{proof}
  \leanok
  Follows immediately from lemma~\ref{MeasureTheory.addEquivAddHaarChar_piCongrRight}.
\end{proof}

\section{Some measure-theoretic preliminaries}

\begin{lemma}
  \label{Topology.IsOpenEmbedding.isHaarMeasure_comap}
  \lean{Topology.IsOpenEmbedding.isHaarMeasure_comap}
  \leanok
  \discussion{507}
  Let $A$ and $B$ be locally compact topological groups
  and let $f:A\to B$ be both a group homomorphism and open embedding.
  The pullback along $f$ of a Haar measure on $B$ is a Haar measure on $A$.
\end{lemma}
\begin{proof}
  \leanok
  Translation-invariance is easy, compact sets are finite because continuous
  image of compact is compact, open sets are bounded because image of open is open.
\end{proof}

\begin{lemma}
  \label{Topology.IsOpenEmbedding.regular_comap}
  \lean{Topology.IsOpenEmbedding.regular_comap}
  \leanok
  \discussion{513}
  The pullback of a regular Borel measure along an open
  embedding is a regular Borel measure.
\end{lemma}
\begin{proof}
  \leanok
  Again this is because the image of compact is compact and the
  image of open is open, so all the properties of being a regular measure
  are easily checked.
\end{proof}

\begin{lemma}
  \label{MeasureTheory.mulEquivHaarChar_eq_one_of_compactSpace}
  \lean{MeasureTheory.mulEquivHaarChar_eq_one_of_compactSpace}
  \leanok
  \discussion{532}
  Say $A$ is a compact topological additive group and $\phi:A\to A$ is an additive isomorphism.
Then $d_A(\phi)=1.$
\end{lemma}
\begin{proof}
  \uses{MeasureTheory.mulEquivHaarChar_smul_preimage}
  \leanok
  We have $d_A(\phi)\mu(A)=\mu(A)$
  from lemma~\ref{MeasureTheory.addEquivAddHaarChar_smul_preimage}
  and $\mu(A)$ is positive and finite because $A$ is open and compact.
\end{proof}

\begin{lemma}
  \label{MeasureTheory.addEquivAddHaarChar_eq_addEquivAddHaarChar_of_isOpenEmbedding}
  \lean{MeasureTheory.addEquivAddHaarChar_eq_addEquivAddHaarChar_of_isOpenEmbedding}
  \leanok
  \discussion{551}
  If $f:A\to B$ is a group homomorphism and open embedding between locally compact
  topological additive groups and if $\alpha:A\to A$ and $\beta:B\to B$ are additive
  homeomorphisms such that the square commutes (i.e., $f\circ\alpha=\beta\circ f$)
  then $d_A(\alpha)=d_B(\beta)$.
\end{lemma}
\begin{proof}
  \uses{Topology.IsOpenEmbedding.isHaarMeasure_comap,
    Topology.IsOpenEmbedding.regular_comap,
    MeasureTheory.addEquivAddHaarChar_smul_integral_comap}
  \leanok
  Choose a regular Haar measure $\mu_B$ on $B$. We just saw
  in lemmas~\ref{Topology.IsOpenEmbedding.isHaarMeasure_comap}
  and~\ref{Topology.IsOpenEmbedding.regular_comap} that its pullback
  $\mu_A:=f^*\mu_B$ to $A$ is also a regular Haar
  measure. Now fix a continuous compactly-supported function $g$ on $A$ with
  $0<\int g(a)d\mu(a)<\infty$. Then $d_A(\alpha)\int g(a)d\mu_A(a)=\int g(a)d(\alpha^*\mu_A)(a)$
  by lemma~\ref{MeasureTheory.addEquivAddHaarChar_smul_integral_comap}.
  This equals $\int g(a)d(\alpha^* f^*\mu_B)(a)$ by definition,
  which is $\int g(a)d(f^*\beta^*\mu_B)(a)$ because pullback of pullback is pullback.
  This equals $d_B(\beta)\int g(a) d(f^*\mu_B)(a)$ by
  corollary~\ref{MeasureTheory.addEquivAddHaarChar_comap}
  which is $d_B(\beta)\int g(a)d\mu_A(a)$ by definition,
  and so $d_A(\alpha)=d_B(\beta)$ as required.
\end{proof}

\section{Restricted products}

Now say $A=\prod'_i A_i$ is the restricted product of a collection of types $A_i$
  with respect to the subsets $C_i$. Recall that this is the subset of $\prod_i A_i$
  consisting of y Say $B=\prod'_i B_i$ is the restricted
  product of types $B_i$ over the same index set, with respect to
  subsets $D_i$. Say $\phi_i:A_i\to B_i$ are functions
  with the property that $\phi_i(C_i)\subseteq D_i$ for all but finitely many $i$.
It is easily checked that the $\phi_i$ induce a function $\phi:=\prod'_i\phi_i:A\to B$. It is also easily
checked that if all the $A_i$ and $B_i$ are groups or rings or $R$-modules, the $C_i$ and $D_i$
are subgroups or subrings or submodules,
and the $\phi_i$ are group or ring or module homomorphisms, then $\phi$ is a group or ring or
module homomorphism. However topological facts lie a little deeper.

\begin{lemma}
  \lean{Continuous.restrictedProduct_congrRight}
  \label{Continuous.restrictedProduct_congrRight}
  \leanok
  \discussion{531}
  If the $A_i$ and $B_i$ are topological spaces and the $\phi_i$ are continuous functions,
  then the restricted product $\phi = \prod'_i\phi_i$ is a continuous function.
\end{lemma}
\begin{proof}
  \leanok
  We use the universal property {\tt RestrictedProduct.continuous\_dom} of the
  topology in mathlib to reduce to the claim that for all finite $S$,
  the induced map $A_S:=\prod_{i\in S}A_i\times\prod_{i\notin S}C_i\to B$ is continuous.
  Because the inclusion $A_S\to A_T$ is continuous for $S\subseteq T$ we are reduced
  to checking this claim for $S$ sufficiently large that it contains all of the $i$
  for which $\phi(C_i)\not=D_i$. For such $S$, this map $A_S\to B$ factors as $A_S\to B_S\to B$
  and $B_S\to B$ is continuous, so it suffices to prove that $A_S\to B_S$ is continuous, but
  this is just a product of continuous maps.
\end{proof}

We now focus on the case that $B_i=A_i$ are locally compact groups, $D_i=C_i$ are compact
open subgroups, and $\phi_i:A_i\to A_i$ are group isomorphisms and homeomorphisms sending
$C_i$ onto $C_i$ for all but finitely many $i$. Then the restricted product $A:=\prod'A_i$
of the $A_i$ with respect to the $C_i$ is also a locally compact topological group, and the
restricted product $\phi=\prod'\phi_i$ of the $\phi_i$ is a group isomorphism and homeomorphism,
so we can ask how $d_A(\phi)$ compares to the $d_{A_i}(\phi_i)$.

First note that $d_{A_i}(\phi_i)=1$ for all the $i$ such that $\phi_i(C_i)=C_i$, as
$d_{A_i}(\phi_i)$ can be computed as $\mu(\phi_i(C_i))/\mu(C_i)$ and $\mu(C_i)$ is
guaranteed to have positive finite measure as it is open and compact. Hence the product
$\prod_id_{A_i}\phi_i$ is a finite product, in the sense that all but finitely many terms are 1.
The following theorem shows that the value of this product is $d(\phi)$.

\begin{theorem}
  \label{MeasureTheory.addEquivAddHaarChar_restrictedProductCongrRight}
  \lean{MeasureTheory.addEquivAddHaarChar_restrictedProductCongrRight}
  \leanok
  \discussion{552}
  With $A$, $A_i$, $C_i$, $\phi_i$, $\phi$ defined as above, we have
  $\delta_A(\phi)=\prod_i\delta_{A_i}(\phi_i)$.
\end{theorem}
\begin{remark} In the Lean file we make the additional assumption that the index set over
  which we're taking the product, is countable, and that the $A_i$ are second countable.
  This is because in this proof we make use
  of an infinite product $\prod_{i\notin S}C_i$ of topological spaces equipped with Haar measure,
  so each topological space gets equipped with the Borel sigma algebra.
  Mathlib then gives the product a canonical product sigma algebra which, in the case where the
  index set is uncountable, may not be the Borel sigma algebra. However we also want to use Haar
  measure on the product so we need the ``canonical'' sigma algebra on this product to be the
  Borel sigma algebra. Rather than taking the trouble
  to locally turn off this product sigma algebra construction in FLT, we simply restrict to a
  countable index set and add second countability assumptions, because then the product is a
  Borel sigma algebra and this is all we need in our applications to adeles.
\end{remark}
\begin{proof}
  \leanok
  \uses{
  MeasureTheory.addEquivAddHaarChar_eq_addEquivAddHaarChar_of_isOpenEmbedding,
  MeasureTheory.mulEquivHaarChar_eq_one_of_compactSpace,
  MeasureTheory.addEquivAddHaarChar_piCongrRight,Continuous.restrictedProduct_congrRight,
  MeasureTheory.addEquivAddHaarChar_prodCongr}
  Assume $\phi_i(C_i)=C_i$ for all $i\not\in S$, a finite set, and work in the
  open subgroup $U:=\prod_{i\in S}A_i\times\prod_{i\not\in S}C_i$. Then $\phi$ induces
  an automorphism $\phi_S$ of this open subgroup $U$ of $A$, and in particular
  lemma~\ref{MeasureTheory.addEquivAddHaarChar_eq_addEquivAddHaarChar_of_isOpenEmbedding} tells us
  that $\delta(\phi)=\delta_U(\phi_S)$.
  Next note that $\phi_S:U\to U$ can be written as a product of the automorphisms
  $\prod_{i\not\in S}\phi_i|_{C_i}$ of $\prod_{i\not\in S}C_i$ and
  $\prod_{i\in S}\phi_i$ of $\prod_{i\in S}A_i$, so by lemma~\ref{MeasureTheory.addEquivAddHaarChar_prodCongr}
  we have $\delta(\phi)=\delta(\prod_{i\not\in S}\phi_i|_{C_i})\times\delta(\prod_{i\in S}\phi_i).$
  Because $\prod_{i\not\in S}C_i$ is a compact
  group we must have $\delta(\phi_{i\not\in S}\phi_i|_{C_i})=1$
  by lemma~\ref{MeasureTheory.mulEquivHaarChar_eq_one_of_compactSpace}. Finally
  $\delta(\prod_{i\in S}\phi_i)=\prod_{i\in S}\delta(\phi_i)$ by
  lemma~\ref{MeasureTheory.addEquivAddHaarChar_piCongrRight} and we are home.
\end{proof}

As a special case, if $R$ is the restricted product of a collection of topological rings $R_i$
  (not necessarily commutative) each equipped with a compact open subring $C_i$, then
  we have

\begin{corollary}
  \label{MeasureTheory.ringHaarChar_restrictedProduct}
  \lean{MeasureTheory.ringHaarChar_restrictedProduct}
  \leanok
  \discussion{554}
  If $u=(u_i)_i\in R^\times$ then $\delta_R(u)=\prod_i\delta_{R_i}(u_i)$.
\end{corollary}
\begin{proof}
  \leanok
  \uses{MeasureTheory.addEquivAddHaarChar_restrictedProductCongrRight}
  By definition of restricted product we have $u_i\in C_i$ for all but finitely many $i$.
  Note also that $u$ has an inverse $v=(v_i)_i$ with $v_i\in C_i$ for all but finitely many $i$.
  The fact that $u_iv_i=v_iu_i=1$ means that $u_i,v_i\in C_i^\times$ for all but finitely many $i$.
  Thus the previous theorem~\ref{MeasureTheory.addEquivAddHaarChar_restrictedProductCongrRight} applies.
\end{proof}

\section{Adeles}

We finish this miniproject by proving some results about Haar characters for
algebras over adele rings.
So let $K$ be a number field and let $\A_K$ be the adeles of $K$.
We will prove some theorems about $\A_K$-algebras $R$ which are finite
and free as $\A_K$-modules. Such algebras can be given the $\A_K$-module topology
and this makes them into locally compact topological rings. In fact we shall only be concerned
in applications with algebras of the form $B\otimes_K\A_K$ where $B$ is a finite-dimensional
$K$-algebra. So fix such a $B$, and write $B_{\A}$ for $B\otimes_K\A_K$. Let us first
deal with a subtlety. Recall that if $K\subseteq L$ are number fields, then $\A_L$ is a
module-finite $\A_K$-algebra and hence an $\A_K$-module, and
theorem~\ref{NumberField.AdeleRing.baseChange_moduleTopology}
tells us that $\A_L$ has the $\A_K$-module topology. Thus the next lemma applies.

\begin{lemma}
  \label{IsModuleTopology.continuous_bilinear_of_finite_left}
  \lean{IsModuleTopology.continuous_bilinear_of_finite_left}
  \leanok
  Say $R$ and $S$ are topological rings, and $S$ is an $R$-algebra, finite as an $R$-module.
  Assume that the topology
  on $S$ is the $R$-module topology. Now say $M$ is an $S$-module, and give it the induced
  $R$-module structure. Then the $R$-module topology and $S$-module topology on~$M$ coincide.
\end{lemma}
\begin{proof}
  \leanok
  Let $i:R\to S$ denote the structure map. First observe that $S$ has the $R$-module topology
  so the $R$-action map $R\times S\to S$ (explicitly defined by $(r,s)\mapsto i(r)s$)
  is continuous, and restricting to $s=1$ we deduce that $i$ is continuous.

  Now let $M_R$ and $M_S$ denote $M$ with the $R$-module and $S$-module topologies respectively.
  It suffices to prove that the identity maps $M_R\to M_S$ and $M_S\to M_R$ are continuous.
  Equivalently, because the $A$-module topology on an $A$-module is the finest topology
  making it into a topological module, we need to prove that $M_R$ is a topological $S$-module
  and that $M_S$ is a topological $R$-module. We start with the latter claim.

  First observe that $M_S$ is a topological $S$-module, so addition is continuous.
  Next note that the map $R\times M_S\to M_S$ factors through $S\times M_S$ and is hence the
  composite of two continuous maps and thus continuous. Hence $M_S$ is a topological $R$-module.

  It thus remains to check that $M_R$ is a topological $S$-module, or equivalently
  that the map $S\times M_R\to M_R$ is continuous. But this map is $R$-bilinear, and
  by the result {\tt Module.continuous\_bilinear\_of\_finite} in {\tt mathlib}, any
  $R$-bilinear map between modules with the $R$-module topology is automatically continuous
  if one of the source modules is finitely-generated. The result applies because $S$ is
  assumed to be a finite $R$-module and the proof is complete.
\end{proof}

\begin{corollary}
  \label{NumberField.AdeleRing.ModuleBaseChangeContinuousLinearEquiv}
  \lean{NumberField.AdeleRing.ModuleBaseChangeContinuousLinearEquiv}
  \leanok
  If $K$ is a number field and $V$ is an $K$-module, then
  the natural isomorphism $V\otimes_K\A_K=V\otimes_{\Q}\A_{\Q}$ induced by the natural
  isomorphism $\A_K=K\otimes_K\A_{\Q}$ is a homeomorphism if the left hand side has the $\A_K$-module
  topology and the right hand side has the $\A_{\Q}$-module topology.
\end{corollary}
\begin{proof}
  \leanok
  \uses{IsModuleTopology.continuous_bilinear_of_finite_left}
  Lemma~\ref{IsModuleTopology.continuous_bilinear_of_finite_left} tells us that $V\otimes_K\A_K$
  has the $\A_{\Q}$-module topology, and it is easily checked that the isomorphism is
  $\A_{\Q}$-linear and hence automatically continuous.

  Note that in the Lean we prove this for a general extension $L/K$ rather than $K/\Q$.
\end{proof}

As a consequence, if $B$ is a $K$-algebra then we can think of $B_{\A}$ as either $B\otimes_K\A_K$
with the $\A_K$-module topology or as $B\otimes_{\Q}\A_{\Q}$ with the $\A_{\Q}$-module
topology. Note that this isomorphism commutes with the inclusions from $B$ into these rings.
But Lean is picky about these things so we'll have to be careful.

\begin{theorem}
  \label{NumberField.AdeleRing.isCentralSimple_addHaarScalarFactor_left_mul_eq_right_mul}
  \lean{NumberField.AdeleRing.isCentralSimple_addHaarScalarFactor_left_mul_eq_right_mul}
  \leanok
  Let $B$ be a finite-dimensional central simple $K$-algebra.
  Say $u\in B_{\A}^\times$, and define $\ell_u$ and $r_u:B_{\A}\to B_{\A}$ by
  $\ell_u(x)=ux$ and $r_u(x)=xu$. Then $d_{B_{\A}}(\ell_u)=d_{B_{\A}}(r_u)$.
\end{theorem}
\begin{proof}
  \uses{MeasureTheory.addEquivAddHaarChar_restrictedProductCongrRight,
    IsSimpleRing.ringHaarChar_eq_addEquivAddHaarChar_mulRight}
  We think of $B_{\A}$ as $B\otimes_K\A_K$.
  If $u=(u_v)$ as $v$ runs through the places of $K$ then
  $d_{B_{\A}}(\ell_u)=\prod_v d_{B_v}(\ell_{u_v})$ by
  theorem~\ref{MeasureTheory.addEquivAddHaarChar_restrictedProductCongrRight} (and the product is finite).
  By corollary~\ref{IsSimpleRing.ringHaarChar_eq_addEquivAddHaarChar_mulRight}
  this equals $\prod_v d_{B_v}(r_{u_v})$, and again by
  theorem~\ref{MeasureTheory.addEquivAddHaarChar_restrictedProductCongrRight} this is $d_{B_{\A}}(r_u)$.
\end{proof}

The previous theorem only applies to inner forms of matrix algebras, but the below theorem,
a generalization of the adelic product formula, is valid for any finite-dimensional
$K$-algebra. Before we state it let's remind ourselves of the product formula for $\Q$,
and restate it in the language of these Haar characters.

\begin{lemma}
  \label{MeasureTheory.ringHaarChar_adeles_rat}
  \lean{MeasureTheory.ringHaarChar_adeles_rat}
  \leanok
  If $x\in\A_{\Q}^\times$ then $\delta_{\A_{\Q}}(x)=\prod_v|x_v|_v.$
\end{lemma}
\begin{proof}
  \leanok
  \uses{MeasureTheory.addEquivAddHaarChar_prodCongr,MeasureTheory.ringHaarChar_real,
    MeasureTheory.addEquivAddHaarChar_restrictedProductCongrRight,MeasureTheory.ringHaarChar_padic}
  By theorem~\ref{MeasureTheory.addEquivAddHaarChar_prodCongr}
  we have $\delta_{\A_{\Q}}(x)=\delta_{\A_{\Q}^\infty}(x^\infty)\times\delta_{\R}(x_\infty)$.
  By lemma~\ref{MeasureTheory.ringHaarChar_real} we have $\delta_{\R}(x_\infty)=|x|_\infty$, and by
  theorem~\ref{MeasureTheory.addEquivAddHaarChar_restrictedProductCongrRight} we have
  $\delta_{\A_{\Q}^\infty}=\prod_p\delta_{\Q_p}(x_p)$. By lemma~\ref{MeasureTheory.ringHaarChar_padic}
  we have $\delta_{\Q_p}(x_p)=|x_p|_p$ and putting everything together we get the result.
\end{proof}

Now $\A_{\Q}$ is nonzero a $\Q$-algebra and hence we have an inclusion $\Q^\times\to\A_{\Q}^\times$.
Here is our reinterpretation of the product formula.

\begin{lemma}
  \label{MeasureTheory.ringHaarChar_adeles_units_rat_eq_one}
  \lean{MeasureTheory.ringHaarChar_adeles_units_rat_eq_one}
  \leanok
  If $x\in\Q^\times\subseteq\A_{\Q}^\times$ then $\delta_{\A_{\Q}}(x)=1.$
\end{lemma}
\begin{proof}
  \leanok
  \uses{MeasureTheory.ringHaarChar_adeles_rat}
  By lemma~\ref{MeasureTheory.ringHaarChar_adeles_rat} we have $\delta_{\A_{\Q}}(x)=\prod_v|x|_v$.
  But the product formula says that this is 1.
  A quick proof: if $x=\pm\prod_pp^{e_p}$ then $\prod_p|x|_p=\prod_pp^{-e_p}$
  and $|x|_\infty=\prod_pp^{e_p}$ so they cancel.
\end{proof}

  Next we generalize this to finite-dimensional $\Q$-vector spaces.

  So say $V$ is an $N$-dimensional $\Q$-vector space,
  and define $V_{\A}:= V\otimes_{\Q}\A_{\Q}$ with its $\A_{\Q}$-module topology.
  If we choose an isomorphism $V\cong\Q^N$ then $V_{\A}\cong\A_{\Q}^N$
  as an additive topological abelian group. In particular, $V_{\A}$ is locally compact.

  Fix a $\Q$-linear automorphism $\phi:V\to V$. By base extension $\phi$ induces
  an $\A_{\Q}$-linear automorphism $\phi_{\A}$ of $V_{\A}$ which is also a homeomorphism of $V_{\A}$
  if $V_{\A}$ is given the module topology as an $\A_{\Q}$-module. Our goal is

  \begin{theorem}
    \label{MeasureTheory.addHaarScalarFactor_tensor_adeles_eq_one}
    \lean{MeasureTheory.addHaarScalarFactor_tensor_adeles_eq_one}
    \leanok
    In the above situation ($V$ a finite-dimensional $\Q$-vector space, $\phi:V\cong V$ is
    $\Q$-linear, $\phi_{\A}$ the base extension to $V_{\A}:=V\otimes_{\Q}{\A_{\Q}}$, a continuous linear
    endomorphism of $V_{\A}$ with the $\A_{\Q}$-module topology), we have $d_{V_{\A}}(\phi_{\A})=1.$
  \end{theorem}
  \begin{proof}
    \leanok
    \uses{
      MeasureTheory.addEquivAddHaarChar_eq_ringHaarChar_det,
      MeasureTheory.addEquivAddHaarChar_restrictedProductCongrRight,
      NumberField.AdeleRing.ModuleBaseChangeContinuousLinearEquiv
    }
    The original blueprint proof of this was that $\phi_{\mathbb{A}} : V_{\mathbb{A}}\to V_{\mathbb{A}}$
    could be written as a restricted product of $\phi_v:V\otimes_{\Q}{\Q_v} \to V\otimes_{\Q}\Q_v$
    and hence by theorem~\ref{MeasureTheory.addEquivAddHaarChar_restrictedProductCongrRight}
    we have $d_{V_{\A}}(\phi_{\A})=\prod_p d_{V_p}(\phi_p)\times d_{V_\infty}(\phi_\infty)$,
    and then applying Lemma~\ref{MeasureTheory.addEquivAddHaarChar_eq_ringHaarChar_det} this
    is equal to $\prod_v\delta_{\Q_v}(\det(\phi_v))\prod_v\delta_{\Q_v}(\det(\phi))=1$.

    This turned out to be a nightmare to formalize, because commuting the tensor product
    and the restricted product cannot be done naively, as the tensor product is over $\Q$
    and the submodules in the restricted product defining $\A_{\Q}$ are not $\Q$-modules.
    So one has to choose a $\Z$-lattice in $V$ and
    use the isomorphisms $V\otimes_{\mathbb{Q}}\mathbb{A}=\Lambda\otimes_{\mathbb{Z}}\mathbb{A}=
    \Lambda\otimes_{\mathbb{Z}}\mathbb{Q}_\infty\times\prod'_p[\Lambda\otimes_{\mathbb{Z}}\mathbb{Q}_p;\Lambda\otimes_{\mathbb{Z}}\mathbb{Z}_p]=
    V\otimes_{\mathbb{Q}}\mathbb{Q}_\infty\times\prod'_p[V\otimes_{\mathbb{Q}}\mathbb{Q}_p;im(\Lambda\otimes_{\Z}\mathbb{Z}_p)]$
    and check that all of these canonical maps are continuous (and one of these claims
    boils down to yet another claim of the form "if you do something to the factors and then
    take the restricted product, then this is topologically the same as doing it to the restriced product",
    with the thing being $\Lambda\otimes_{\Z}$ in this case, something which needs
    checking).

    So here is the proof which we actually formalized. Say an automorphism of a finite free
    $R$-module is \emph{nice} if it's a product of transvections and diagonal matrices with unit
    entries. Mathlib has the fact that if $R$ is a field then all automorphisms are nice,
    and the base change of a nice morphism is nice. Hence $\phi_{\mathbb{A}}$ is nice,
    and we can simply prove Lemma~\cite{MeasureTheory.addEquivAddHaarChar_eq_ringHaarChar_det}
    for nice endomorphisms over a commutative ring, which gives the result we want immediately
    by the product formula.
  \end{proof}

  \begin{corollary}
    \label{NumberField.AdeleRing.units_mem_ringHaarCharacter_ker}
    \lean{NumberField.AdeleRing.units_mem_ringHaarCharacter_ker}
    \leanok
    If $B$ is a finite-dimensional $\Q$-algebra (for example a number field, or a quaternion algebra over a number field),
    if $B_{\A}$ denotes the ring $B\otimes_{\Q}\A_{\Q}$, and if $b\in B^\times$,
    then $\delta_{B_{\A}}(b)=1$.
  \end{corollary}
  \begin{proof}
    \uses{MeasureTheory.addHaarScalarFactor_tensor_adeles_eq_one}
    \leanok
    Follows immediately from the previous theorem.
  \end{proof}

  \begin{corollary}
    \label{NumberField.AdeleRing.addEquivAddHaarChar_mulRight_unit_eq_one}
    \lean{NumberField.AdeleRing.addEquivAddHaarChar_mulRight_unit_eq_one}
    \leanok
    If $B$ is a finite-dimensional $\Q$-algebra and
    if $b\in B^\times$ then right multiplication by $b$
    does not change Haar measure on $B_{\A}$.
  \end{corollary}
  \begin{proof}
    \leanok
    \uses{MeasureTheory.addHaarScalarFactor_tensor_adeles_eq_one}
    Follows immediately from the previous theorem.
  \end{proof}


%%FILE blueprint/src/chapter/HeckeOperatorProject.tex
\chapter{Miniproject: Hecke Operators}\label{HeckeOperator_project}

\section{Status}

This is an active miniproject. The abstract theory is completely formalized;
at the time of writing the concrete theory has no sorried definitions
but it does have some sorried proofs.

\section{The goal}

The goal of this project is to get sorry-free definitions of Hecke operators
acting on spaces of automorphic forms. These Hecke operators generate
Hecke algebras, which are the rings called $T$ in the modularity lifting theorems,
or $R=T$ theorems, crucially introduced by Wiles in order to prove FLT.

The theory comes in two parts;
the ``abstract'' theory, which is pure algebra, and the ``concrete'' theory
where we apply the abstract constructions to produce endomorphisms of
spaces of automorphic forms. The abstract theory is short (and completely formalized);
the concrete theory still needs some work because to apply the theory to the
adelic groups we care about we need to develop some more API around the theory
of restricted products, and of compact open subgroups of matrix groups.

\section{The abstract theory}

\subsection{Introduction}

The set-up: we have a commutative ring $R$, the coefficient ring, and
all of our spaces which the operators act on will be $R$-modules.

We have a group $G$ acting $R$-linearly on an $R$-module $A$.

We have subgroups $U$ and $V$ of $G$.
We will be particularly interested in the $R$-modules $A^U$ and $A^V$
of invariant elements.

Given an element $g\in G$, then under a certain finiteness hypothesis
we will be able to define an $R$-linear map $T_g$ or $[UgV]$
from $A^V$ to $A^U$. The finiteness hypothesis is that the
double coset $UgV$ can be written as a \emph{finite} union of single
cosets $g_iV$.

\begin{definition}
  \lean{AbstractHeckeOperator.heckeOperator.toFun}
  \label{AbstractHeckeOperator.HeckeOperatorToFun}
  \leanok
  Assuming $UgV$ is a finite union of cosets $g_iV$,
  we define $[UgV]:A^V\to A^U$ to be the map sending $a\in A^V$
  to $\sum_i g_ia.$
\end{definition}

\begin{lemma}
  \lean{AbstractHeckeOperator.heckeOperator}
  \label{AbstractHeckeOperator.HeckeOperator}
  \uses{AbstractHeckeOperator.heckeOperator.toFun}
  \leanok
  This function is well-defined (that is, independent of the
  choice of $g_i$), has image in $A^U$ and is $R$-linear.
\end{lemma}
\begin{proof}
  \leanok
  Well-definedness is because if we change $g_i$ to $g'_i:=g_iv$
  for some $v\in V$ then $g_ia=g_i'a$ because $a\in A^V$.

  The image lands in $A^U$ because left multiplication by $u$
  fixes $UgV$ and hence permutes the cosets $g_iV$.

  Finally $R$-linearity is because the $G$-action is $R$-linear.
\end{proof}

The group $G$ is not in general commutative, and hence if $U=V$
the Hecke operators in this generality do not in general commute
as endomorphisms of $A^U$. Here is a criterion for
them to commute.

\begin{lemma}
  \lean{AbstractHeckeOperator.comm}
  \label{AbstractHeckeOperator.comm}
  \uses{AbstractHeckeOperator.heckeOperator}
  \leanok
  Say $g,h\in G$ and we have $UgU=\coprod_i g_iU$
  and $UhU=\coprod_j h_j$ and we have $g_ih_j=h_jg_i$ for all $i,j$.
  Then $[UgU][UhU]=[UhU][UgU]$, that is, the Hecke operators
  acting on $A^U$ commute.
\end{lemma}
\begin{proof}
  \leanok
  We have $[UgU][UhU]a=\sum_ig_i(\sum_jh_ja)=\sum_{i,j}g_ih_ja$
  and $[UhU][UgU]a=\sum_jh_j\sum_ig_ia=\sum_{j,i}h_jg_ia$ and these
  sums are equal because $g_ih_j=h_jg_i$.
\end{proof}

The finiteness hypothesis that the decomposition $UgV=\coprod_i g_iV$ is
into a finite union is necessary for the theory to work. If $G$ is a topological
group then here is a criterion which gives the finiteness hypothesis for free.

\begin{lemma}
  \lean{QuotientGroup.mk_image_finite_of_compact_of_open}
  \label{QuotientGroup.mk_image_finite_of_compact_of_open}
  \discussion{563}
  \leanok
  If $U$ and $V$ are compact subgroups of a topological group~$G$,
  if $V$ is also open, and if $g\in G$, then the double coset space $UgV$
  is a finite union of left cosets $g_iV$.
\end{lemma}
\begin{proof}
  \leanok
  The subset $UgV$ of $G$ is a continuous image of the compact set $U\times V$
  and is hence compact, and it is covered by the disjoint left cosets $g_iV$;
  this cover must thus be finite.
\end{proof}

\section{Restricted products}

In the concrete example of Hecke operators which we care about, the invariants $A^G$
will be spaces of quaternionic automorphic
forms (by definition). We do not need to worry about the definition of $A$ in this
project at all. However we will need to do various computations with the
specific groups $G$ which we are interested in, and they are restricted products.
So we now develop some theory for restricted products,
starting by recalling the definition.

If $I$ is an index set, if $X_i$ are sets indexed by $i\in I$ and if $Y_i$
are subsets, then the \emph{restricted product} $\prod'_iX_i$ (note the dash) is defined
to be the subset of the full product $\prod_i X_i$ consisting of those
tuples $(x_i)$ such that $x_i\in Y_i$ for all but finitely many~$i$. We suppress
the $Y_i$ from the notation in this document, although in Lean we cannot do this and
the restricted product looks something like $\prod{}^{r} i,[X\ i, Y\ i]$.

It is straightforward to check that if the $X_i$ are groups or rings or $R$-modules,
and the $Y_i$ are subgroups or subrings or submodules, then the restricted product
is a group, ring or $R$-module; indeed the structure is inherited via the
inclusion $\prod'_iX_i\subseteq\prod_iX_i$ (and the fact that arbitrary products
of groups/rings/modules are groups/rings/modules).

More subtle is the theory of topological space structures. If the $X_i$
are topological spaces then we do \emph{not} give $\prod'_iX_i$ the subspace
topology coming from the product topology on $\prod_iX_i$; instead we give
it the finest topology making all of the
natural maps $\prod_{i\in S}X_i\times\prod_{i\notin S}Y_i\to \prod'_iX_i$ continuous,
as $S$ runs through all finite subsets of~$I$; here the product of $X_i$s and $Y_i$s
has the product topology. For example if all of the $Y_i$ are open
then one can check that $\prod_iY_i$ is an open subset of
$\prod'_iX_i$ (this is {\tt RestrictedProduct.isOpen\_forall\_mem} in mathlib),
whereas it is not of the form $\left(\prod'_iX_i\right)\cap U$ for any open subset~$U$
of $\prod_iX_i$ in general; the map from $\prod'_i X_i$ to $\prod_i X_i$ is continuous
but is not in general an embedding.

If you've seen automorphic forms before, then here is an obvious-sounding claim:
because the adeles $\mathbb{A}_F$ of a number field are a restricted product of
completions $F_v$ with respect to the integer rings $\mathcal{O}_v$, then
$GL_2(\mathbb{A}_F)$ is obviously topologically a restricted product of
$GL_2(F_v)$ with respect to $GL_2(\mathcal{O}_v)$. We now spend some time justifying
this claim, which is a little more intricate than it sounds.

\subsection{Products}

Here are some basic facts we need about restricted products.

\begin{lemma}
  \lean{Homeomorph.restrictedProductProd}
  \label{Homeomorph.restrictedProductProd}
  \leanok
  \discussion{568}
  If $A_i$ is a family of topological spaces equipped with open
  subsets $B_i$, and if $C_i$ is a family of topological spaces equipped
  with open subsets $D_i$, and if we equip $A_i\times C_i$ with the open
  subset $B_i\times D_i$, then the natural bijection
  $\prod'_i(A_i\times C_i)=\left(\prod'_iA_i\right)\times\left(\prod'_iB_i\right)$
  is a homeomorphism.
\end{lemma}
\begin{remark} This may well not be true if $B_i$ and $D_i$ are not open, because
  filtered colimits and binary products do not appear in general to commute
  in the category of topological spaces. I don't know an explicit counterexample though.
\end{remark}
\begin{proof}
  \leanok
  \proves{Homeomorph.restrictedProductProd}
  We need to check continuity in both directions. The easy way is
  continuity of the map from the restricted product to the map from the binary
  product; the lemma {\tt RestrictedProduct.continuous\_dom} in mathlib
  tells us that a map from a restricted product is continuous when its restriction
  to $\left(\prod_{i\in S}(A_i\times C_i)\right)\times\left(\prod_{i\notin S}(B_i\times D_i)\right)$
  is continuous for all finite $S\subseteq I$; the universal property of the binary
  product tells us that the map into the binary product is continuous iff the maps into
  the factors are continuous, but the map into $\prod'_iX_i$ is a product of the
  natural maps from $\left(\prod_{i\in S}(A_i\times C_i)\right)\times\left(\prod_{i\notin S}(B_i\times D_i)\right)$
  to $\left(\prod_{i\in S}A_i\right)\times\left(\prod_{i\notin S}B_i\right)$
  and the inclusion, and both are known to be continuous (an arbitrary product of continuous
  maps is continuous, and the other claim is in the restricted product API in mathlib).

  The harder direction is the other way, because we are working against both universal
  properties. The trick is {\tt RestrictedProduct.continuous\_dom\_prod} in mathlib
  (this is where we assume $B_i$ and $D_i$ are open), which tells us that a map out of
  a binary product of restricted products is continuous when its restriction to
  $\left(\left(\prod_{i\in S}A_i\right)\times\left(\prod_{i\notin S}B_i\right)\right)\times
  \left(\left(\prod_{i\in S}C_i\right)\times\left(\prod_{i\notin S}D_i\right)\right)$
  is, for all finite $S$ (note that the $S$ in the mathlib lemma is actually our $I-S$).
  The map from this to the restricted product factors through
  $\left(\prod_{i\in S}(A_i\times C_i)\right)\times\left(\prod_{i\notin S}(B_i\times D_i)\right)$;
  the first map is a homeomorphism (use the fact that $\prod_iX_i\times Y_i$ is homeomorphic
  to $\left(\prod_iX_i\right)\times\left(\prod_iY_i\right)$), and the second is continuous
  by definition of the topology on a restricted product.
\end{proof}

\begin{corollary}
  \lean{Homeomorph.restrictedProductPi}
  \label{Homeomorph.restrictedProductPi}
  \leanok
  \discussion{570}
  Restricted products (with respect to open subspaces) commute with finite products.
  In other words, if $j$ runs through a finite set $J$ and $i$ runs through an arbitrary
  set $I$, and if $X_{ji}$ are topological spaces equipped with open subspaces $Y_{ji}$,
  then the obvious bijection $\prod'_i(\prod_j X_{ji})=\prod_j\left(\prod'_i X_{ji}\right)$
  is a homeomorphism.
\end{corollary}
\begin{proof}
  \leanok
  \uses{Homeomorph.restrictedProductProd}
  Induction on the size of the finite set, using lemma~\ref{Homeomorph.restrictedProductProd}
  to get you started.
\end{proof}

 Let $n$ be a natural and let $M_n(X)$ for a set $X$ denote ``$n\times n$
  matrices with coefficient in $X$'', i.e. $X^{n^2}$. If $X$ is a topological
  spaces then give $M_n(X)$ the product topology.

\begin{corollary}
  \label{Homeomorph.restrictedProductMatrix}
  \lean{Homeomorph.restrictedProductMatrix}
  \leanok
  \discussion{571}
  If $X_i$ are topological spaces and the $Y_i$ are open subspaces,
  then the obvious map $M_n(\prod'_iX_i)=\prod'_iM_n(X_i)$ is a homeomorphism.
\end{corollary}
\begin{proof}
  \leanok
  \uses{Homeomorph.restrictedProductPi}
  Immediate from the previous corollary~\ref{Homeomorph.restrictedProductPi}.
\end{proof}

\subsection{Units}

We now want to move from matrices to invertible matrices whilst keeping track of topology,
so we need to understand units of topological monoids. Openness of the subobject was
crucial in the above arguments, so we need the next lemma before we can get anywhere.

\begin{lemma}
  \lean{Submonoid.units_isOpen}
  \label{Submonoid.units_isOpen}
  \leanok
  \discussion{587}
  If $M$ is a topological monoid and $U$ is an open submonoid, then
  the units $U^\times$ of $U$ are naturally an open subgroup of $M^\times$.
\end{lemma}
\begin{remark} Note that $M^\times$ doesn't get the subspace topology from~$M$,
  it is embedded into $M\times M$ via $g\mapsto (g,g^{-1})$ and gets the subspace
  topology from the product. This makes it into a topological group.
\end{remark}
\begin{proof}
  \leanok
  \proves{Submonoid.units_isOpen}
  We have $U\times U$ is an open subset of $M\times M$, and if we imagine $M^\times$
  embedded in $M\times M$ as explained in the remark above, then the intersection
  of this subgroup with $U\times U$ is open in $M^\times$ and consists of the elements
  of $M^\times$ which are in $U$ and whose inverse is also in $U$, which is easily
  checked to be the copy of $U^\times$ we're talking about.
\end{proof}

Later on, compactness will be key for us, so we record the analogous result
for compactness.

\begin{lemma}
  \lean{Submonoid.units_isCompact}
  \label{Submonoid.units_isCompact}
  \leanok
  \discussion{588}
  If $M$ is a Hausdorff topological monoid and $U$ is a compact submonoid,
  then the units $U^\times$ of $U$ are naturally a compact subgroup of $M^\times$.
\end{lemma}
\begin{remark} Is Hausdorffness necessary?
\end{remark}
\begin{proof}
  \proves{Submonoid.units_isCompact}
  First I claim that $M^\times$ embedded in $M\times M$ via $g\mapsto (g,g^{-1})$
  is a closed subset of $M\times M$. Indeed, if $p:M\times M\to M$ is $(a,b)\mapsto ab$
  and $q:M\times M\to M$ is $(a,b)\mapsto ba$, then $p$ and $q$ are continuous,
  $M^\times\subseteq M\times M$ is the intersection
  $p^{-1}\{1\}\cap q^{-1}\{1\}$, and $\{1\}$ is closed because $M$ is Hausdorff.

  We have $U\times U$ is a compact subset of $M\times M$, and so
  $U^\times=M^\times\cap U\times U$ is a closed subspace of a compact space
  and is thus compact.
\end{proof}

\begin{lemma}
  \lean{ContinuousMulEquiv.piUnits}
  \label{ContinuousMulEquiv.piUnits}
  \leanok
  \discussion{581}
  If $U_i$ are topological monoids then the canonical
  group isomorphism $(\prod_i U_i)^\times=\prod_i(U_i^\times)$ is a homeomorphism.
\end{lemma}
\begin{proof}
  \leanok
  We prove that the maps in both directions are continuous. Let's start
  with the map from left to right.

  A map into a product is continuous when the maps to the factors
  are continuous. A map into the units of a monoid is continuous when the
  two projection maps to the monoid (the inclusion and the map $u\mapsto u^{-1}$)
  are continuous (because $M^\times$ has the topology induced from $M\times M$).
  This reduces us to checking that the maps $(\prod_i U_i)^\times\to U_j$
  sending $(u_i)$ to $u_j$ resp $u_j^{-1}$ are continuous. But the former map
  is the continuous inclusion $(\prod_i U_i)^\times\to\prod_i U_i$ followed
  by the continuous projection to $U_j$, and the latter map is the continuous
  inclusion $(\prod_i U_i)^\times\to\prod_i U_i$ sending $x$ to $x^{-1}$
  followed by the projection.

  To go the other way: because the units have the induced topology it suffices
  to check that the two maps $\prod_i(U_i^\times)\to\prod_i U_i$
  sending $(u_i)$ to $(u_i)$ resp $(u_i^{-1})$ are continuous. A map
  to a product is continuous when the induced maps to the factors are.
  A projection from a project is continuous, and the identity and inverse are
  continuous maps $U_j^\times\to U_j$, and the maps we're concerned with are composites
  of these maps and are hence continuous.
\end{proof}

\begin{theorem}
  \lean{ContinuousMulEquiv.restrictedProductUnits}
  \label{ContinuousMulEquiv.restrictedProductUnits}
  \leanok
  \discussion{582}
  If $M_i$ are a family of topological monoids equipped with open
  submonoids $U_i$, then the canonical map $(\prod'_iM_i)^\times\to\prod'_i(M_i^\times)$
  is a homeomorphism.
\end{theorem}
\begin{proof}
  \leanok
  \uses{ContinuousMulEquiv.piUnits, Submonoid.units_isOpen}
  I don't know a clean way of showing that the map from left to right is continuous,
  so here is a ``direct'' proof that the map is a homeomorphism. It is certainly an abstract group
  isomorphism between topological groups. So to prove that it is a homeomorphism it suffices
  to prove that it is a homeomorphism near the identity, or equivalently that there are open
  neighbourhoods $X$ and $Y$ of the identity elements on each side such that the map induces a
  homeomorphism from $X$ to $Y$. We choose $(\prod_i U_i)^\times$ and $\prod_i (U_i^\times)$.
  Note that the former is open because of lemma~\ref{Submonoid.units_isOpen}.
  The result now follows from the previous lemma~\ref{ContinuousMulEquiv.piUnits}.
\end{proof}

\section{Some local theory}

We could work over a general nonarchimedean normed field but we still do not have them
in mathlib, so we stick to the case of interest which is the completion of a number
field~$K$ at a finite place~$v$. Such a completion is a topological field~$K_v$
equipped with a discrete valuation, a ring of integers $\calO_v$ having a principal
maximal ideal $(\varpi)$, and a finite residue field $k_v:=\calO_v/(\varpi)$.

There is no formal Lean code for the lemmas in this section; I am slightly dragging my
feet because it would seem more sensible to prove them in the right generality,
and we don't have a definition of nonarchimedean local field yet.

\begin{lemma}
  \lean{NumberField.isOpenAdicCompletionIntegers}
  \label{NumberField.isOpenAdicCompletionIntegers}
  \leanok
  $\calO_v$ is an open subring of $K_v$.
\end{lemma}
\begin{proof}
  \leanok
  Openness is already in mathlib.
\end{proof}

$\calO_v$ is also a compact subring of $K_v$; we proved this in the adele
miniproject.

\begin{lemma} $M_2(\calO_v)$ is an open subring of $M_2(K_v)$.
  \lean{IsDedekindDomain.M2.localFullLevel.isOpen}
  \label{M2.localFullLevel.isOpen}
  \leanok
\end{lemma}
\begin{proof}
  \leanok
  \uses{NumberField.isOpenAdicCompletionIntegers}
  Topologically $M_2(\calO_v)\cong \calO_v^4$ as a subset of $K_v^4$ so this
  follows because a product of compacts is compact and a product of opens is open.
\end{proof}

\begin{lemma} $M_2(\calO_v)$ is a compact subring of $M_2(K_v)$.
  \lean{IsDedekindDomain.M2.localFullLevel.isCompact}
  \label{M2.localFullLevel.isCompact}
  \leanok
\end{lemma}
\begin{proof}
  \leanok
  \uses{NumberField.instCompactSpaceAdicCompletionIntegers}
  Topologically $M_2(\calO_v)\cong \calO_v^4$ as a subset of $K_v^4$ so this
  follows because a product of compacts is compact and a product of opens is open.
\end{proof}

\begin{lemma} $GL_2(\calO_v)$ is a compact open subgroup of $GL_2(K_v)$.
  \label{nolean-compactopen-GL2}
\end{lemma}
\begin{proof}
  \uses{nolean-compactopen-matrix, Submonoid.units_isOpen, Submonoid.units_isCompact}
  $K_v$ is known to be Hausdorff, so $M_2(K_v)$ is Hausdorff and the
  results follow from lemmas~\ref{Submonoid.units_isOpen} and~\ref{Submonoid.units_isCompact}.
\end{proof}

Recall that there is a projection $\calO_v\to k_v$ where $k_v$ is the residue
field of~$v$, a finite field. This induces a ring homomorphism $M_2(\calO_v)\to M_2(k_v)$
with kernel $M_2(\varpi\calO_v)$, an ideal~$I$ of $M_2(\calO_v)$ isomorphic to $(\varpi\calO_v)^4$
and hence also compact and open.

Say $\Gamma_v$ is a subgroup of $GL_2(k_v)$. Then $\Gamma_v$ is finite. Consider it as a submonoid of
the multiplicative monoid $M_2(k_v)$. Its preimage~$U_v$ in $M_2(\calO_v)$
is easily checked to be a submonoid of $M_2(\calO_v)$; furthermore it is a finite union of
cosets of $I$ and is hence compact and open as a submonoid of $M_2(\calO_v)$ and hence of $M_2(K_v)$.

\begin{lemma}
  \label{nolean-compactopen-U1p} $U_v$ is a compact open subgroup of $GL_2(K_v)$.
\end{lemma}
\begin{proof}
  \uses{Submonoid.units_isOpen, Submonoid.units_isCompact}
  $\Gamma_v$ is a group and hence its preimage $U_v$ is a subgroup of the monoid
  $M_2(K_v)$. It is compact and open as we just saw. Hence its units (also $U_v$)
  are a subgroup of $GL_2(K_v)$, which is compact and open, again by
  lemmas~\ref{Submonoid.units_isOpen} and~\ref{Submonoid.units_isCompact}.
\end{proof}

Say now $\begin{pmatrix}1&*\\0&1\end{pmatrix}\subseteq\Gamma_v\subseteq\begin{pmatrix}*&*\\0&*\end{pmatrix}$
and let $U:=U_v$ be its preimage in $GL_2(\calO_v)$, considered as a compact open subgroup of $GL_2(K_v)$.
Choose $0\not=\alpha\in\calO_v$ and define $g=\begin{pmatrix}\alpha&0\\0&1\end{pmatrix}\in GL_2(K_v)$.
Let's do an explicit double coset decomposition in preparation for a calculation with Hecke operators.

\begin{lemma}
  \lean{TotallyDefiniteQuaternionAlgebra.WeightTwoAutomorphicForm.HeckeOperator.Local.bijOn_unipotentMulDiagU1}
  \label{bijOn_unipotent_mul_diagU1_U1diagU1}
  \leanok
  The double coset space $UgU$ is the disjoint union of $g_tU$ as $t$ ranges
  through $\calO_v/\alpha\calO_v$ and $g_t:=\begin{pmatrix}\alpha&\tilde{t}\\0&1\end{pmatrix}$,
  where $\tilde{t}$ is any lift of $t$ to $\calO_v$.
\end{lemma}
\begin{proof}
  \leanok
  We first manipulate the statement into a statement about finite groups.
  We have $UgU=\coprod_t g_tU\iff UgUg^{-1}=\coprod_t g_tUg^{-1}=\coprod_t g_tg^{-1}(gUg^{-1})$.
  By the second isomorphism theorem this is true if
  $U=\coprod_t g_tg^{-1}(gUg^{-1}\cap U)$. So when is an element of $U$
  in $gUg^{-1}$? Equivalently, if $x\in U$, when is $g^{-1}xg\in U$? An explicit calculation
  of matrices shows us that this is true iff $g=\begin{pmatrix} a&b\\c&d\end{pmatrix}$ with
  $\alpha\mid b$. Define $U^\alpha$ to be this subgroup of~$U$. We have reduced the question
  to showing that the matrices $h_t:=\begin{pmatrix}1&\tilde{t}\\0&1\end{pmatrix}$
  are a set of left coset representatives for the subgroup $U^\alpha$ of $U$.

  It thus suffices to show that if $u=\begin{pmatrix} a&b\\c&d\end{pmatrix}\in U$
  then $u\in h_tU^\alpha$ iff $b\in\calO_v$ reduces mod $\alpha$ to $t\in\calO_v/\alpha$.
  We do this by computing $h_t^{-1}u=\begin{pmatrix} a-\tilde{t}c&b-\tilde{t}d\\c&d\end{pmatrix}$
  and observing that its top right hand entry mod~$\alpha$ is zero iff $b$ mod $\alpha$ is $t$.
\end{proof}

\section{Adelic groups}

We are finally ready to discuss the group~$G$ and the subgroups~$U$ which we will be
using to define our Hecke operators. Let~$K$ be a number field, let~$D$ a quaternion algebra
over~$K$ and let $\A_K^\infty$ be the finite adeles of~$K$; recall that this is a commutative
topological ring, defined to be the restricted
product of the commutative topological fields~$K_v$ as $v$ runs through the finite places
of~$K$, with respect to the compact open subrings $\calO_v$.

The group~$G$ we are interested in for the rest of this miniproject is the group
$(D\otimes_K\A_K^\infty)^\times$. We want to write down compact open subgroups of this group,
but the first thing we need to do is to find a way of talking about elements of the group.

We will assume that there exists an $\A_K^\infty$-algebra isomorphism
$D\otimes_K\A_K^\infty=M_2(\A_K^\infty)$ and we will fix such an isomorphism $r$
(called a \emph{rigidification} in the Lean code). We give both of these $\A_K^\infty$-algebras
the $\A_K^\infty$-module topology, which is a fancy way of saying the product topology
(they are both free of rank 4 as $\A_K^\infty$-modules); the rigidification is then
a homeomorphism (because all $\A_K^\infty$-module maps between modules with the $\A_K^\infty$-module
topology are continuous).

This means that our group~$G$ is isomorphic (both algebraically and topologically)
to $GL_2(\A_K^\infty)$. Before we go any further,
let say something about matrix rings over complete fields.

\begin{theorem}
  \lean{IsDedekindDomain.FiniteAdeleRing.GL2.restrictedProduct}
  \label{GL2.restrictedProduct}
  \leanok
  $G$ is isomorphic and homeomorphic
  to the restricted product of $GL_2(K_v)$ with respect
  to the compact open subgroups $GL_2(\calO_v)$.
\end{theorem}
\begin{proof}
  \leanok
  \uses{ContinuousMulEquiv.restrictedProductUnits,
    Homeomorph.restrictedProductMatrix, ContinuousMulEquiv.restrictedProductMatrixUnits}
  This follows from lemma~\ref{ContinuousMulEquiv.restrictedProductUnits}
  and lemma~\ref{Homeomorph.restrictedProductMatrix}.
\end{proof}

If~$S$ is a finite set of finite places of~$K$, and for each $v\in S$ we choose
a subgroup $\Gamma_v$ of $GL_2(k_v)$ then we saw in the previous section how to
create a compact open subgroup $\tilde{\Gamma}_v$ of $GL_2(K_v)$. For $v\notin S$
define $\tilde{\Gamma}_v=GL_2(\calO_v)$. Then $\prod_v\tilde{\Gamma}_v$ is a compact
open subgroup of $\prod_vGL_2(\calO_v)$. It is compact subgroups of this form
which we shall be using.

\section{Automorphic forms}

We recall some of the definitions of spaces of automorphic forms, from
the quaternion algebra project, section~\ref{Quat_alg_project}.

We fix a totally real field~$F$, a
totally definite quaternion algebra~$D/F$, and a coefficient (additive abelian) group~$R$.
Set $G=(D\otimes_F\A_F^\infty)^\times$ as in the previous section. Note that $G$
naturally contains copies of $D^\times$ and $(\A_F^\infty)^\times$. Recall from
definition~\ref{TotallyDefiniteQuaternionAlgebra.WeightTwoAutomorphicForm}
that an $R$-valued weight~2 automorphic form
is a function $f:G\to R$ satisfying the following axioms:
\begin{enumerate}
  \item $f(dg)=f(g)$ for all $d\in D^\times\subseteq(D\otimes_F\A_F^\times)$;
  \item There exists a compact open subgroup of~$U$ (the \emph{level} of~$f$)
    such that $f(gu)=f(g)$ for all $g\in G$ and $u\in U$;
  \item $f(gz)=f(g)$ for all $z\in(\A_F^\infty)^\times$.
\end{enumerate}

It can be checked that the collection of all such forms is an additive abelian group,
and if~$R$ is a ring then it is naturally an~$R$-module. Let's call this group~$A$ for short.
Then~$A$ has a left action of $G$, with $g\cdot f$
defined via $(g\cdot f)(x):=f(xg)$. Recall from
definition~\ref{TotallyDefiniteQuaternionAlgebra.WeightTwoAutomorphicFormOfLevel}
that a weight 2 automorphic form of level~$U$ is simply an element of the
fixed points $A^U$. In other words, the forms of level~$U$ are the forms satisfying the three axioms
defining an automorphic form but with the compact open subgroup in the second
axiom being~$U$.

\section{Concrete Hecke operators}

Let~$F$ be a number field. For each finite place $v$ we have the completion $F_v$ of $F$ at $v$,
which is a normed field equipped with its integer ring $\calO_v$, a local ring with
finite residue field~$k_v$.

For $v$ a finite place of~$F$, let $\Delta_v$ be a subgroup of $k_v^\times$ and consider
the subgroup $\Gamma_v$ of $GL_2(k_v)$ consisting of matrices $\begin{pmatrix}a&b\\0&d\end{pmatrix}$
with $a,d\in k_v^\times$ and $a/d\in\Delta_v$. Then

It is easily checked that this is a subgroup, and that
$\begin{pmatrix}1&*\\0&1\end{pmatrix}\subseteq\Gamma_v\subseteq\begin{pmatrix}*&*\\0&*\end{pmatrix},$
so lemma~\ref{nolean-U1-coset-decomposition} applies.
Let $U_{\Delta_v}$ be the preimage of this subgroup in $GL_2(\calO_v)$.
This is a compact open subgroup of $GL_2(\calO_v)$, by the remarks above.

Let~$S$ be a finite set of finite places of~$F$, and define $U_\Delta(S)$ to be the matrices
in $\prod_v GL_2(\calO_v)$ which are in $U_{\Delta_v}$ for all $v\in S$ (we put no condition
at the places $v\notin S$). We can consider
$U_{\Delta}(S)$ as a subgroup of $GL_2(\A_F^\infty)$; it is a product of compact subgroups
and thus compact, and it is a product of opens all but finitely many of which are $GL_2(\calO_v)$
and is thus open. Because the inclusion $\prod_v GL_2(\calO_v)\to GL_2(\A_F^\infty)$ is an
open embedding, we can regard $U_\Delta(S)$ as a compact open subgroup of $GL_2(\A_F^\infty)$.

If we fix $r$ a rigidification, it
induces an isomorphism $GL_2(\A_F^\infty)=(D\otimes_F\A_F^\infty)^\times$
and so we can identify~$U_\Delta(S)$ with its image in $(D\otimes_F\A_F^\infty)$.

We introduce Hecke operators of two types.

First, for $v$ any place not in~$S$ we choose a uniformiser $\varpi_v\in F_v$,
form the invertible $2\times 2$ matrix $\begin{pmatrix}\varpi_v&0\\0&1\end{pmatrix}\in GL_2(F_v)$
and extend this element to an element $g\in G$ by letting its component at all finite places
$w\not=v$ be the identity. We define the Hecke operator $T_v:A^U\to A^U$ to be $[UgU]$, using the
notation defined at the beginning of this section.

For the second kind of Hecke operator we choose
$0\not=\alpha\in\calO_v$ and we consider the $2\times 2$ matrix in $GL_2(\A_K^\infty)$
which is $\begin{pmatrix}\alpha&0\\0&1\end{pmatrix}$ at~$v$ and 1 at all other local components.
Via the rigidification~$r$ we obtain an element $g\in G$.
We define the Hecke operator $U_{v,\alpha}$ to be $[UgU]$.

The Hecke algebra of interest to us will be generated by the Hecke operators $T_v$ for $v\notin S$
and $U_{v,\alpha}$ for $v\in S$.

The big theorem we want in this section is
\begin{theorem}
  \label{nolean-hecke-algebra-commutative-noetherian}
  Say~$R$ is a Noetherian ring. Then the subalgebra of the $R$-linear endomorphisms
  of $A^U$ generated by the Hecke operators $T_v$ for $v\notin S$ and $U_{v,\alpha}$ for $v\in S$
  is a Noetherian commutative ring.
\end{theorem}

\section{Analysis of the Hecke algebra}

First we discuss commutativity of the Hecke operators. First, assume that $v\not\in S$.
Then $U=GL_2(\calO_v)\times U'$ where $U'$ is a subgroup of the restricted product
of $GL_2(F_w)$ for $w\not=v$. We can use {\tt RestrictedProduct.SubmonoidClass.isProductAt} to express this
concept of being an internal direct product. If $g$ is the element of $G$ used to make $T_v$ then
$g$ is also supported at $w$, so the double coset space $UgU$ is just
$(GL_2(\calO_v)\begin{pmatrix}\varpi&0\\0&1\end{pmatrix}GL_2(\calO_v))\times U'$
and in particular can be decomposed into single left $U$-cosets of the form $g_iU$
where $g_i$ is also supported at~$v$. This is {\tt RestrictedProduct.mem\_coset\_and\_mulSupport\_subset\_of\_isProductAt}.

Similarly if $v\in S$, if $0\not=\alpha\in\calO_v$ and if $g_v=\begin{pmatrix}\alpha&0\\0&1\end{pmatrix}$
and is 1 elsewhere, then the double coset space $UgU$ can again be written as $\coprod_i g_iU$
with the $g_i$ supported only at~$v$.

We deduce immediately from lemma~\ref{AbstractHeckeOperator.comm} that two Hecke operators
associated to different finite places of~$F$ commute.
What remains is to check that $U_{\alpha,v}$ and $U_{\beta,v}$ commute. In fact we claim
more, namely that $U_{\alpha,v}U_{\beta,v}=U_{\alpha\beta,v}$. This will suffice
because $\alpha\beta=\beta\alpha$.

\begin{lemma}
  \label{nolean-product-of-U-alpha}
  If $v\in S$ and $\begin{pmatrix}1&*\\0&1\end{pmatrix}\subseteq\Gamma_v\subseteq\begin{pmatrix}*&*\\0&*\end{pmatrix}$
  then $U_{\alpha,v}U_{\beta,v}=U_{\alpha\beta,v}$.
\end{lemma}
\begin{proof} Follows easily from the explicit double coset decomposition proved above.
\end{proof}

The reason that the Hecke algebra is Noetherian is that the main theorem of the Fujisaki
miniproject immediately implies that $A^G$ is a submodule of a finite free $R$-module
and is hence Noetherian. Its endomorphism algebra is hence a Noetherian $R$-module,
so the sub-$R$-algebra generated by the Hecke operators is also a Noetherian $R$-module
and thus a Noetherian ring.


%%FILE blueprint/src/chapter/QuaternionAlgebraProject.tex
\chapter{Miniproject: Quaternion algebras}\label{Quat_alg_project}

\section{The goal}

At the time of writing, {\tt mathlib} still does not have a proof that the space
of classical modular forms of a fixed weight, level and character is finite-dimensional.
The main result in this miniproject is to prove that certain spaces of quaternionic modular forms
are finite-dimensional. We need finiteness results
like this in order to control the Hecke algebras which act on these spaces; these Hecke
algebras are the ``$T$''s which are isomorphic to the ``$R$''s in the $R=T$ theorem which
is the big first target for the FLT project.

\section{Initial definitions}

Our first goal in this miniproject is the definition of these spaces of quaternionic modular forms.
We start with some preliminary work towards this.

Let $K$ be a field. Recall that a \emph{quaternion algebra}
over $K$ is a central simple $K$-algebra of $K$-dimension~4.

A fundamental fact about central simple algebras is that if $D/K$
is a central simple $K$-algebra and $L/K$ is an extension of fields, then $D\otimes_KL$
is a central simple $L$-algebra. In particular if $D$ is a quaternion algebra over $K$
then $D\otimes_KL$ is a quaternion algebra over $L$. Some Imperial undergraduate students have established
this fact in ongoing project work.

A \emph{totally real field} is a number field~$F$ such that the image of every ring
homomorphism $F\to\bbC$ is a subset of $\R$. We fix once and for all a totally real field~$F$ and a
quaternion algebra $D$ over $F$. We furthermore assume that $D$ is \emph{totally definite}, that is,
that for all field embeddings $\tau:F\to\R$ we have $D\otimes_{F,\tau}\R\cong\bbH$. Because $F$
has at least one real place, the totally definite hypothesis is enough to show that $D$
is not a matrix algebra and thus must be a division algebra. Thus Fujisaki's theorem (theorem~\ref{NumberField.FiniteAdeleRing.DivisionAlgebra.units_cocompact}
from the Fujisaki miniproject) applies, and we know that $D^\times\backslash (D\otimes_F\A_F)^{(1)}$
is compact.

The high-falutin' explanation of what is about to happen is that the units $D^\times$ of $D$
can be regarded as a connected reductive algebraic group over $F$, and we are going to define certain spaces
of automorphic forms for this algebraic group. For a general connected reductive algebraic group,
part of the definition of an automorphic form is that it is satisfies some differential
equations (for example modular forms are automorphic forms for the algebraic group $\GL_2$ over
$\Q$, and the definition of a modular form involves holomorphic functions, which are solutions to
the Cauchy--Riemann equations). However under the assumption that $F$ is totally real and $D/F$ is
totally definite, the ``associated symmetric space is a 0-dimensional manifold'', meaning in practice that
the part of the definition of an automorphic form involving differential equations is vacuously satisfied in this setting.
As a consequence, the definitions we're about to give have a far more algebraic flavour. Crucially, in stark
contrast to the general theory, the fact that we do not need to talk about differential equations at all means that one does
not need to assume that our automorphic forms are $\bbC$-valued; our definition makes sense for forms valued in an arbitrary additive commutative group.
In particular, it is possible to talk about mod $p^n$ and $p$-adic automorphic forms in this setting without doing any complicated algebraic geometry.

\section{Brief introduction to automorphic forms in this setting}

Having made assumptions on $D$ which makes the theory of automorphic forms over $D^\times$
far less technical, we will now make it a little more technical by using the modern adelic
approach to the theory. Note that many results about the adeles of a number field are proved
in the adele miniproject (section \ref{Adele_miniproject}). Our automorphic forms will be certain
functions on the units of the ring $D_{\A^\infty}:=D\otimes_F\A_F^\infty\cong D\otimes_{\Q}\A_{\Q}^\infty$.
To prove Fermat's Last Theorem it suffices to work in weight 2 and trivial character, and for simplicity we shall (at least for now)
bake these assumptions into our definitions, even though they would be easy to remove (at the expense of having
to write ``of weight 2 and trivial character'' throughout the proof). We remark again that
there is no analogue of the holomorphicity condition that one sees in the classical theory,
because the symmetric space is a finite set of points rather than the upper half plane. Also
there is no analogue of the cuspidality condition because there are no cusps in this setting.
Other than the number field $F$ and the quaternion algebra $D$, the other variable we will see in our situation
will be the \emph{level} of the forms. The main result in this miniproject will be that the space of weight 2 automorphic forms of level $U$
is finite-dimensional.

\section{Definition of spaces of automorphic forms}

Let us now give some precise definitions. Recall that by $\A_F^\infty$ we mean the finite adeles of the totally real number field $F$.

A \emph{level} is a compact open subgroup~$U$ of $(D\otimes_F\A_F^\infty)^\times$. These are
plentiful. The ring $D_f:=D\otimes_F\A_F^\infty$ is a topological ring,
and hence the units $D_f^\times$ of this ring
are a topological group. This group is locally profinite, and hence has many compact open subgroups;
we will see explicit examples later on.

We regard $\A_F^\infty$ as a subring of $D_{\A^\infty}:=D\otimes_F\A_F^\infty$, which is possible because
$F$ is a subring of $D$. More precisely we embed $\A_F^\infty$ into $D\otimes_F\A_F^\infty$
via the map sending $g$ to $1\otimes g$. Because $F$ is in the centre of~$D$, we have
that $\A_F^\infty$ is in the centre of $D_{\A^\infty}$ (in fact it is the centre, but we do not need this).
As a consequence we can identify $(\A_F^\infty)^\times$ as a subgroup of $(D\otimes_F\A_F^\infty)^\times$.
We may also regard $D$ as a subring of $D\otimes_F\A_F^\infty$ via the map $d\mapsto d\otimes 1$, and hence
we can think of $D^\times$ as a subgroup of $(D\otimes_F\A_F^\infty)^\times$.

Let $R$ be an additive commutative group. Later on $R$ will be a commutative ring but we will not need this for the definition.
\begin{definition}
  \lean{TotallyDefiniteQuaternionAlgebra.WeightTwoAutomorphicForm}
  \label{TotallyDefiniteQuaternionAlgebra.WeightTwoAutomorphicForm}
  \leanok
  The space of $R$-valued \emph{automorphic forms} for $D^\times$ is the set of
  functions $f:D_{\A^\infty}^\times\to R$ satisfying the following axioms:
  \begin{itemize}
    \item $f(dg)=f(g)$ for all $d\in D^\times$ and $g\in D_{\A^\infty}^\times$.
    \item $f(gz)=f(g)$ for all $g\in D_{\A^\infty}^\times$.
    \item There exists a compact open subgroup $U\subseteq (D_{\A^f}^\times)$ such that $f(gu)=f(g)$ for all $g\in D_{\A^\infty}^\times$ and $u\in U$.
  \end{itemize}
\end{definition}

Let $S^D(R)$ denote the set of automorphic forms for $D^\times$. The space $S^D(R)$ is sometimes referred to as a space of
``quaternionic modular forms'' over $R$.
Three basic observations about $S^D(R)$ are

\begin{definition}
  \lean{TotallyDefiniteQuaternionAlgebra.WeightTwoAutomorphicForm.addCommGroup}
  \label{TotallyDefiniteQuaternionAlgebra.WeightTwoAutomorphicForm.addCommGroup}
  \uses{TotallyDefiniteQuaternionAlgebra.WeightTwoAutomorphicForm}
  \leanok
  Pointwise addition $(f_1+f_2)(g):=f_1(g)+f_2(g)$ makes $S^D(R)$ into an
  additive abelian group.
\end{definition}

\begin{definition}
  \lean{TotallyDefiniteQuaternionAlgebra.WeightTwoAutomorphicForm.module}
  \label{TotallyDefiniteQuaternionAlgebra.WeightTwoAutomorphicForm.module}
  \uses{TotallyDefiniteQuaternionAlgebra.WeightTwoAutomorphicForm,
    TotallyDefiniteQuaternionAlgebra.WeightTwoAutomorphicForm.addCommGroup}
  \leanok
  If $R$ is a commutative ring then pointwise scalar multiplication $(r\cdot f)(g):= r\cdot(f(g))$ makes
  $S^D(R)$ into an $R$-module.
\end{definition}

\begin{definition}
  \lean{TotallyDefiniteQuaternionAlgebra.WeightTwoAutomorphicForm.distribMulAction}
  \label{TotallyDefiniteQuaternionAlgebra.WeightTwoAutomorphicForm.distribMulAction}
  \uses{TotallyDefiniteQuaternionAlgebra.WeightTwoAutomorphicForm,
    TotallyDefiniteQuaternionAlgebra.WeightTwoAutomorphicForm.addCommGroup,
    TotallyDefiniteQuaternionAlgebra.WeightTwoAutomorphicForm.module}
  \leanok
  The group $D_{\A^f}^\times$ acts on the additive abelian group $S^D(R)$ by $(g\cdot f)(x)=f(xg)$.
\end{definition}

If $R$ is a commutative ring then the action of $D_{\A^\infty}^\times$ commutes with the $R$-action.

Now let $U$ be a level, that is, a compact open subgroup of $D_{\A^\infty}^\times$.

\begin{definition}
  \lean{TotallyDefiniteQuaternionAlgebra.WeightTwoAutomorphicForm.LevelStruct.form}
  \label{TotallyDefiniteQuaternionAlgebra.WeightTwoAutomorphicFormOfLevel}
  \uses{TotallyDefiniteQuaternionAlgebra.WeightTwoAutomorphicForm,
    TotallyDefiniteQuaternionAlgebra.WeightTwoAutomorphicForm.addCommGroup,
    TotallyDefiniteQuaternionAlgebra.WeightTwoAutomorphicForm.distribMulAction,
    TotallyDefiniteQuaternionAlgebra.WeightTwoAutomorphicForm.module}
  \leanok
  The quaternionic modular forms of level $U$, with notation $S^D(U;R)$, are the $U$-invariants for the $D_{\A^\infty}^\times$-action on $S^D(R)$.
\end{definition}

The Hecke algebras involved in the main modularity lifting
theorems needed in the FLT project will be endomorphisms of these spaces $S^D(U;R)$.

\section{The main result}

The point of this miniproject is the finite-dimensionality result below.
This is an analogue of the result that classical modular forms of a fixed
level, weight and character are finite-dimensional. In fact, by delicate results
of Jacquet and Langlands this result (in the case $k=\bbC$) implies many cases of that classical claim,
although of course the Jacquet--Langlands theorem is much much harder to prove than the classical
proof of finite-dimensionality.


\begin{theorem}
  \lean{TotallyDefiniteQuaternionAlgebra.WeightTwoAutomorphicForm.LevelStruct.instFiniteSubtypeMemSubmoduleFormOfIsSufficientlySmallOfIsFinite}
  \label{TotallyDefiniteQuaternionAlgebra.WeightTwoAutomorphicForm.finiteDimensional}
  \uses{TotallyDefiniteQuaternionAlgebra.WeightTwoAutomorphicForm.LevelStruct.form,
    TotallyDefiniteQuaternionAlgebra.instIsFinite}
  Let $k$ be a field. Then the space $S^D(U;k)$ is a finite-dimensional $k$-vector space.
\end{theorem}



\begin{proof}
  \proves{TotallyDefiniteQuaternionAlgebra.WeightTwoAutomorphicForm.LevelStruct.instFiniteSubtypeMemSubmoduleFormOfIsSufficientlySmallOfIsFinite}
  \uses{NumberField.FiniteAdeleRing.DivisionAlgebra.units_cocompact,
  NumberField.FiniteAdeleRing.DivisionAlgebra.finiteDoubleCoset}
  The finite-dimensionality theorem is in fact an easy consequence of Fujisaki's lemma,
  proved in the Fukisaki miniproject, chapter~\ref{Fujisaki_project}.
  Write $(D\otimes_F\A_F^\infty)^\times$ as a disjoint union of double cosets
  $\coprod_i D^\times g_i U$. This open cover descends to a disjoint open
  cover of $D^\times\backslash (D\otimes_F\A_F^\infty)^\times$,
  and this latter space is compact by theorem~\ref{NumberField.FiniteAdeleRing.DivisionAlgebra.units_cocompact}.
  Hence the cover is finite; write the double coset representatives
  as $g_1,g_2,\ldots,g_n$. We claim that
  the function $S^D(U;k)\to W^n$ sending $f$ to $(f(g_1),f(g_2),\ldots,f(g_n))$
  is injective and $k$-linear, which suffices by finite-dimensionality of $W$.
  $k$-linearity is easy, so let's talk about injectivity.

  Say $f_1$ and $f_2$ are two elements of $S^D(U;k)$ which agree on
  each $g_i$. It suffices to prove that $f_1(g)=f_2(g)$ for all
  $g\in(D\otimes_F\A_F^\infty)^\times$. So say $g\in(D\otimes_F\A_F^\infty)^\times$,
  and write $g=\delta g_iu$ for $\delta \in D^\times$ and $u\in U$.
  Then $f_1(g)=f_1(\delta g_iu)=f_1(g_i)$ (by the definition of $S^D(U;k)$), and similarly $f_2(g)=f_2(g_i)$
  and because $f_1(g_i)=f_2(g_i)$ by assumption, we deduce $f_1(g)=f_2(g)$ as required.
\end{proof}


%%FILE blueprint/src/chapter/biblio.tex
\bibliographystyle{plain} % We choose the "plain" reference style
\bibliography{FLT}




%%FILE blueprint/src/chapter/ch01introduction.tex
\chapter{Introduction}

Fermat's Last Theorem is the statement that if $a,b,c,n$ are positive whole numbers with $n\geq 3$,
then $a^n+b^n\not=c^n$. It is thus the claim that a family of \emph{Diophantine equations}
($a^3+b^3=c^3, a^4+b^4=c^4,\ldots$) has no positive integer solutions.
Diophantus was a Greek mathematician who lived around 1800
years ago, and he would have been able to understand the statement of the theorem (he knew about
positive integers, addition and multiplication).

Fermat's Last Theorem was explicitly raised by Fermat in 1637, and was proved by Wiles (with the
proof completed in joint work with Taylor) in 1994. There are now several proofs but all of them
go broadly in the same direction, using elliptic curves and modular forms.

Lean is an interactive theorem prover; it checks mathematical arguments with super-human accuracy.
Explaining a proof of Fermat's Last Theorem to Lean is in some sense like explaining the proof to
Diophantus; for example, the proof starts by observing that before we go any further it's convenient
to first invent/discover zero and negative numbers, and one can point explicitly at places in Lean's
source code \href{https://github.com/leanprover/lean4/blob/260eaebf4e804c9ac1319532970544a4e157c336/src/Init/Prelude.lean#L1049}{here}
and \href{https://github.com/leanprover/lean4/blob/260eaebf4e804c9ac1319532970544a4e157c336/src/Init/Data/Int/Basic.lean#L45}{here}
where these things happen. However we will adopt a more efficient approach: we will assume all of
the theorems both in core Lean and in its mathematics library
\href{https://github.com/leanprover-community/mathlib4}{\tt mathlib}, and proceed from there.
To give some idea of what this entails: {\tt mathlib} at the time of writing contains most of an
undergraduate mathematics degree and parts of several relevant Masters level courses (for example,
the definitions and basic properties of
\href{https://leanprover-community.github.io/mathlib4_docs/Mathlib/AlgebraicGeometry/EllipticCurve/Weierstrass.html}{elliptic curves}
and \href{https://leanprover-community.github.io/mathlib4_docs/Mathlib/NumberTheory/ModularForms/Basic.html}{modular forms}
are in {\tt mathlib}). Thus our task can be likened to teaching a graduate level course on
Fermat's Last Theorem to a computer. The computer is quite a challenging audience member -- it
will insist on being given all technical details of all arguments, and it will not accept proof by
intimidation or by appeal to higher authority. Most mathematicians know humans who also behave
in this manner. Up until 2026, the computer would only start filling
in details of arguments by itself once the arguments are mathematically utterly obvious.
This made formalization of mathematics a very time-consuming process. However, with the
advent of AI computers are now getting better at filling in arguments themselves.

\section{Which proof is being formalised?}

At the time of writing, these notes do not contain anywhere near a proof of FLT, or even a sketch proof,
although we are currently actively working on fixing this.

From 2024 to 2029 we will be beginning to build a proof of FLT, following a strategy
basically due to Khare, and modified by Taylor, taking into account Buzzard's comments on what
would be easy or hard to do
in Lean. Our strategy uses refinements of the original Taylor--Wiles method by Diamond/Fujiwara,
Khare--Wintenberger, Skinner--Wiles, Kisin, Taylor and others -- one could call it a 21st century
proof of the theorem. During this initial phase of the project, we shall also be \emph{assuming}
many nontrivial theorems without proof, as long as they were published by 31st December 1989.
To get technical for just a second -- we shall for example be assuming the existence of Galois
representations attached to weight~2 Hilbert modular forms, we will assume Langlands' cyclic base
change theorem for $\GL_2$, Mazur's theorem bounding the torsion subgroup of an elliptic curve over
the rationals, and several other nontrivial results which were known by the end of the 1980s.

The upshot of this is that, by 2029 at the end of this first phase, the project should contain a
complete proof that FLT follows from results that were known to humanity in the 1980s. In
particular, one naive way of understanding the goal is that it is a ``formalization of the papers
of Wiles and Taylor--Wiles, assuming the results in the references of those papers''. However,
as noted above, we will actually be taking a slightly different path.

\section{The structure of this blueprint}

This blueprint is a \emph{nonlinear} document, comprising many chapters. The chapters
are not designed to be read in order. Each chapter is self-contained and has a well-defined goal,
typically stated at the top of the chapter.

After this chapter, you should next read chapter~\ref{ch_reductions}, which explains how to reduce FLT
to two highly nontrivial statements about the $p$-torsion in a certain elliptic curve (the Frey curve).
One of these statements was proved in the 1970s by Mazur, and we shall not be concentrating on
it until after the first phase is complete. The other is a theorem of Wiles, and this is what we will
be concentrating on in the remainder of the blueprint.

\section{Remarks}

The actual blueprint currently also contains a lot of disorganised ideas. Currently these
should be disregarded.

Chapter~\ref{ch_overview} is an \emph{extremely} sketchy overview of how
the rest of the proof goes. This is currently being expanded and should be ignored
right now.

All of the remaining chapters are experiments, and most of them are what I am currently
calling ``mini-projects''. A mini-project is a bottom-up project, typically at early graduate
student level, with a concrete goal. The ultimate goal of many of these projects is to actually
get some result into mathlib. We have had one success so far -- the Frobenius mini-project
is currently being PRed to mathlib by Thomas Browning. Several other mini-projects have
been completed in this repo but have not yet been PRed to mathlib. Examples of completed
miniproject which are not yet upstreamed include adeles
mini-project and the quaternion algebra mini-project. These projects do not logically depend
on each other for the most part, and one can pick and choose how one reads them.

There is also an appendix, which is again very sketchy, and comprises mostly of a big
list of nontrivial theorems many of which we will be assuming without proof in the FLT
project.

The next chapter to read, where the proof begins, is chapter~\ref{ch_reductions}.


%%FILE blueprint/src/chapter/ch02reductions.tex
\chapter{First reductions of the problem}\label{ch_reductions}

\section{Goal}

The goal of this chapter is to reduce FLT to a deep theorem of Mazur and a deep theorem of Wiles
about a Galois representation.

\section{Overview}
The proof of Fermat's Last Theorem is by contradiction. We assume that we have a counterexample
$a^n+b^n=c^n$, and manipulate it until it satisfies the axioms of a ``Frey package'', a basic
concept which we will explain below. From the
Frey package we build a Frey curve -- an elliptic curve defined over the rationals. We then look at
a certain representation of a Galois group coming from this elliptic curve, and finally using two
very deep and independent theorems (one due to Mazur, the other due to Wiles) we show that this
representation is both reducible and irreducible, the contradiction we seek.

\section{Reduction to \texorpdfstring{$n\geq5$}{ngeq5} and prime}

\begin{lemma}\label{FermatLastTheorem.of_odd_primes}\lean{FermatLastTheorem.of_odd_primes}\leanok
  If there is a counterexample to Fermat's Last Theorem, then there is a counterexample $a^p+b^p=c^p$
  with $p$ an odd prime.
\end{lemma}
\begin{proof}\leanok
  Note: this proof is \href{https://leanprover-community.github.io/mathlib4_docs/Mathlib/NumberTheory/FLT/Four.html#FermatLastTheorem.of_odd_primes}{in mathlib already};
  we run through it for completeness' sake.

  Say $a^n + b^n = c^n$ is a counterexample to Fermat's Last Theorem. Every positive integer is either
  a power of 2 or has an odd prime factor. If $n=kp$ has an odd prime factor $p$ then
  $(a^k)^p+(b^k)^p=(c^k)^p$ is the counterexample we seek. It remains to deal with the case where
  $n$ is a power of 2, so let's assume this. We have $3\leq n$ by assumption, so
  $n=4k$ must be a multiple of~4, and thus $(a^k)^4=(b^k)^4=(c^k)^4$, giving us a counterexample
  to Fermat's Last Theorem for $n=4$. However an old result of Fermat himself (proved as
  \href{https://leanprover-community.github.io/mathlib4_docs/Mathlib/NumberTheory/FLT/Four.html#fermatLastTheoremFour}{\tt fermatLastTheoremFour}
  in {\tt mathlib}) says that $x^4+y^4=z^4$ has no solutions in positive integers.
\end{proof}

Euler proved Fermat's Last Theorem for $p=3$;

\begin{lemma}\label{fermatLastTheoremThree}\lean{fermatLastTheoremThree}\leanok
\discussion{16}
  There are no solutions in positive integers to $a^3+b^3=c^3$.
\end{lemma}
\begin{proof}
  \leanok
  The proof in mathlib was formalized by a team from the ``Lean For the Curious Mathematician'' conference held in Luminy in March 2024
  (its dependency graph can be visualised \href{https://pitmonticone.github.io/FLT3/blueprint/dep_graph_document.html}{\underline{here}}).
\end{proof}

\begin{corollary}\label{FermatLastTheorem.of_p_ge_5}\lean{FermatLastTheorem.of_p_ge_5}\leanok If there is a counterexample to
  Fermat's Last Theorem, then there is a counterexample $a^p+b^p=c^p$ with $p$ prime and $p\geq 5$.
\end{corollary}
\begin{proof}\uses{fermatLastTheoremThree, FermatLastTheorem.of_odd_primes}\leanok Follows from the
  previous two lemmas.\end{proof}

\section{Frey packages}

For convenience we make the following definition.

\begin{definition}\label{FreyPackage}\lean{FreyPackage}\leanok A \emph{Frey package} $(a,b,c,p)$
is three nonzero pairwise-coprime integers $a$, $b$, $c$, with $a\equiv3\pmod4$ and $b\equiv0\pmod2$,
and a prime $p\geq5$, such that $a^p+b^p=c^p$.\end{definition}

Our next reduction is as follows:

\begin{lemma}
  \label{FreyPackage.of_not_FermatLastTheorem_p_ge_5}
  \lean{FreyPackage.of_not_FermatLastTheoremFor_p_ge_5}
  \leanok
  \discussion{19}
  If Fermat's Last Theorem is false for $p \ge 5$ and prime, then there exists a Frey package.
\end{lemma}
\begin{proof}
  \uses{FreyPackage}
  \leanok
  Suppose we have a counterexample $a^p+b^p=c^p$ for the given $p$; we now build a Frey package from this data.

  If the greatest common divisor of $a,b,c$ is $d$ then $a^p+b^p=c^p$ implies $(a/d)^p+(b/d)^p=(c/d)^p$.
  Dividing through, we can thus assume that no prime divides all of $a,b,c$. Under this assumption we
  must have that $a,b,c$ are pairwise coprime, as if some prime divides two of the integers $a,b,c$
  then by $a^p+b^p=c^p$ and unique factorization it must divide all three of them.
  In particular we may assume that not all of $a,b,c$ are even, and now reducing modulo~2 shows
  that precisely one of them must be even.

  Next we show that we can find a counterexample with $b$ even. If $a$ is the even one then we can
  just switch $a$ and $b$. If $c$ is the even one then we can replace $c$ by $-b$ and $b$ by $-c$
  (using that $p$ is odd).

  The last thing to ensure is that $a$ is 3 mod~4. Because $b$ is even, we know that $a$ is odd,
  so it is either~1 or~3 mod~4. If $a$ is 3 mod~4 then we are home; if however $a$ is 1 mod~4 we
  replace $a,b,c$ by their negatives and this is the Frey package we seek.
\end{proof}

\section{Galois representations and elliptic curves}\label{twopointfour}

To continue, we need some of the theory of elliptic curves over $\Q$. So let $f(X)$ denote any
monic cubic polynomial with rational coefficients and whose three complex roots are distinct,
and let us consider the equation $E:Y^2=f(X)$, which defines a curve in the $(X,Y)$ plane.
This curve (or strictly speaking its projectivisation) is a so-called elliptic curve
(or an elliptic curve over $\Q$ if we want to keep track of the field where the coefficients
of $f(X)$ lie).

If $E:Y^2=f(X)$ is an elliptic curve over $\Q$, and if $K$ is any characteristic zero field (and
hence a $\Q$-algebra), then we write $E(K)$ for the set of solutions to $y^2=f(x)$ with $x,y\in K$,
together with an additional ``point at infinity'' corresponding morally to $x=y=\infty$.
It is an extraordinary fact,
and not at all obvious, that $E(K)$ naturally has the structure of an additive abelian group,
with the point at infinity being the zero element (the identity). Fortunately this fact is
already in {\tt mathlib}. This additive group structure has
the property that three distinct points $P,Q,R\in K^2$ which are in $E(K)$ will sum to zero
if and only if they are collinear.

The group structure behaves well under change of field: if $E$ is an elliptic curve over $\Q$
and if $K\to L$ is a homomorphism of
characteristic zero fields then the induced map $E(K)\to E(L)$ is a group homomorphism.
Thus if $f:K\to L$ is an isomorphism of characteristic zero fields, the induced map $E(K)\to E(L)$
is an isomorphism of groups, with the inverse isomorphism being the map $E(L)\to E(K)$ induced
by $f^{-1}$. This construction thus gives us an action of the multiplicative group $\Aut(K)$
of automorphisms of the field $K$ on the additive abelian group $E(K)$, and hence also
on the $n$-torsion of this group for any positive integer~$n$.
In particular, if $\Qbar$ denotes an algebraic closure of the
rationals (for example, the algebraic numbers in $\bbC$) and if $\GQ$ denotes the group of field
isomorphisms $\Qbar\to\Qbar$, then for any elliptic curve $E$ over $\Q$ we have an action
of $\GQ$ on the additive abelian group $E(\Qbar)$, and hence on its $n$-torsion subgroup
$E(\Qbar)[n]$.

If furthermore $n=p$ is prime, then $E(\Qbar)[p]$ is naturally a vector space over the
field $\Z/p\Z$, and thus it inherits the structure of a mod $p$ representation of $\GQ$.
This is the \emph{mod $p$ Galois representation} attached to the elliptic curve $E$.
It is well-known to be 2-dimensional. We call this representation $\rho_{E,p}$.

In the next section we apply this theory to an elliptic curve coming from a counterexample to
Fermat's Last theorem.

\section{The Frey curve}

Recall that a \emph{Frey package} $(a,b,c,p)$ is simply a prime $p\geq5$ and nonzero
pairwise-coprime integers $a,b,c$ satisfying $a^p+b^p=c^p$ and satisfying the congruences
$a\equiv3\pmod4$ and $b\equiv0\pmod2$. We have shown above that if Fermat's Last Theorem is false,
then a Frey package exists.

\begin{definition}[Frey]
  \label{FreyCurve}
  \lean{FreyPackage.freyCurve}
  \uses{FreyPackage}\discussion{21}
  \leanok
  Given a Frey package $(a,b,c,p)$, the corresponding \emph{Frey curve} (considered by Frey and,
  before him, Hellegouarch) is the elliptic curve over $\Q$ defined by the equation
  $Y^2=X(X-a^p)(X+b^p).$
\end{definition}

Note that the roots of the cubic $X(X-a^p)(X+b^p)$ are distinct because $a,b,c$ are nonzero and
$a^p+b^p=c^p$.

Given a Frey package $(a,b,c,p)$ with corresponding Frey curve $E$, the mod $p$ Galois
representation $\rho_{E,p}$ associated to this package is the 2-dimensional representation of
$\GQ$ on $E(\Qbar)[p]$ described above. Frey's observation is that this mod $p$ Galois
representation has some very surprising properties. We will make this remark more explicit
in the next chapter. Here we shall show how these properties can be used to finish the job.

\section{Reduction to two big theorems.}

Recall that a representation of a group $G$ on a vector space $W$ is said to be \emph{irreducible}
if there are precisely two $G$-stable subspaces of $W$, namely $0$ and $W$.
The representation is said to be \emph{reducible} otherwise.

Now say $(a,b,c,p)$ is a Frey package.
Consider the mod $p$ representation of $\GQ$ coming from the $p$-torsion in the Frey
curve $Y^2=X(x-a^p)(X+b^p)$ associated to the package. Let's call this representation $\rho$,
and we say that $\rho$ is the mod $p$ representation associated to the Frey package $(a,b,c,p)$.
Is it irreducible or not?

\begin{theorem}[Mazur]
  \label{Mazur_Frey}
  \lean{FreyPackage.mazur}
  \uses{FreyCurve}
  \leanok
  If $\rho$ is the mod $p$ Galois representation associated to a Frey package $(a,b,c,p)$ then
  $\rho$ is irreducible.
\end{theorem}
\begin{proof}
  \notready
  This follows from a profound and long result of Mazur \cite{mazur-torsion} from 1977,
  namely the fact that the torsion subgroup of an elliptic curve over $\Q$ can have size at most~16.
  In fact there is still a little more work which needs to be done to deduce the theorem from
  Mazur's result. A pre-1990 reference for the full proof of this claim is Proposition~6
  in~\S4.1 of~\cite{serreconj}.
\end{proof}

Note that in the first (pre-2029) phase of the FLT project, we will not be working on
a formalization of this result, as it was known in the 1980s. We will however be thinking
a lot about the next result, which says the exact opposite.

\begin{theorem}[Wiles,Taylor--Wiles, Ribet,\ldots]
  \label{Wiles_Frey}
  \lean{FLT.Bosses.B4_proof}
  \uses{FreyCurve}
  \leanok
  If $\rho$ is the mod $p$ Galois representation associated to a Frey package $(a,b,c,p)$ then
  $\rho$ is reducible.
\end{theorem}
\begin{proof}
%  \uses{modularity_lifting_theorem,frey_curve_hardly_ramified,moret-bailly}
  This is the main content of Wiles' magnum opus.
  We omit the argument for now, although later on in this project
  we will have a lot to say about a proof of this.
\end{proof}

\begin{corollary}
  \label{FreyPackage.false}
  \lean{FLT.Bosses.B3_proof}
  \uses{Mazur_Frey, Wiles_Frey}
  \leanok
  There is no Frey package.
\end{corollary}
\begin{proof} Follows immediately from the previous two
  theorems~\ref{Mazur_Frey} and~\ref{Wiles_Frey}.
\end{proof}

We deduce

\begin{corollary}
  \label{FLT}
  \lean{flt}
  \leanok
  Fermat's Last Theorem is true. In other words, there are no positive integers $a,b,c$ and
  natural $n\geq3$ such that $a^n+b^n=c^n$.
\end{corollary}
\begin{proof}
  \uses{FermatLastTheorem.of_p_ge_5, FreyPackage.false, FreyPackage.of_not_FermatLastTheorem_p_ge_5}
  Assume there is a there is a counterexample $a^n+b^n=c^n$.
  By Corollary \ref{FermatLastTheorem.of_p_ge_5} we may assume that there is also a counterexample
  $a^p+b^p=c^p$ with $p\geq 5$ and prime.
  Then there is a Frey package $(a,b,c,p)$ by~\ref{FreyPackage.of_not_FermatLastTheorem_p_ge_5},
  contradicting Corollary~\ref{FreyPackage.false}.
\end{proof}

Because we are (for now at least) assuming Mazur's theorem, we now need to turn our attention
to a proof of theorem~\ref{Wiles_Frey}. We start on this proof in Chapter~\ref{ch_freyreduction}.


%%FILE blueprint/src/chapter/ch03freyold.tex
\chapter{Elliptic curves, and the Frey Curve}\label{ch_frey}

\section{Overview}

In Chapter~\ref{ch_reductions} we explained how, given a counterexample to Fermat's Last Theorem, we could construct a Frey package and thus a Frey curve, which is an elliptic curve with some interesting properties. In this chapter we start with an overview of parts of the theory of the arithmetic of elliptic curves. Following
this we sketch proofs of the two main results of this chapter: firstly that the $p$-torsion $\rho$ in the Frey curve is ``hardly ramified'', and secondly that Mazur's result on the possible torsion of elliptic curves implies that $\rho$ must be irreducible. Everything here follows from standard results about elliptic curves, however almost none of these results are in {\tt mathlib} as I am writing this, so there is plenty to be done here.

\section{The arithmetic of elliptic curves}

We give an overview of the results we need, citing the literature for proofs. Everything here is
 standard, and most of it dates back to the 1970s or before.

\begin{theorem}\label{WeierstrassCurve.n_torsion_card}\lean{WeierstrassCurve.n_torsion_card}\notready
  Let $n$ be a positive integer, let $F$ be a separably closed
  field with $n$ nonzero in $F$, and let $E$ be an elliptic curve over $F$. Then the $n$-torsion $E(K)[n]$
  in the $F$-points of $E$ is a finite group of size $n^2$.
\end{theorem}
\begin{proof}
  There are several proofs in the textbooks. The proof being worked on uses the theory of division polynomials; the formalisation is ongoing work of David Angdinata, and it will be part of his PhD thesis.
\end{proof}

This theorem actually tells us the structure of the $n$-torsion, because of the following
purely group-theoretic result:
\begin{lemma}\label{group_theory_lemma}\lean{group_theory_lemma}
  Say $n$ is a positive integer, $r$ is a natural, and $A$ is an abelian group. Assume that for all $d\mid n$, the $d$-torsion $A[d]$ of $A$ has size $d^r$. Then $A[n]\cong (\Z/n\Z)^r$.
\end{lemma}
\begin{proof}
  The result is obvious if $n=1$, so we may assume $n>1$. One proof would be to write $A$ as $\prod_{i=1}^t(\Z/a_i\Z)$
  with $a_i\mid a_{i+1}$ (this is possible by the structure theorem for finite abelian groups), and then to apply our hypothesis firstly with $d=a_1$ to deduce $t=r$ and then with $d=a_t$ to deduce $a_1=a_t$.
\end{proof}

\begin{corollary}\label{Elliptic_curve_n_torsion_2d}\lean{WeierstrassCurve.n_torsion_dimension}\leanok
  Let $n$ be a positive integer, let $F$ be a separably closed
  field with $n$ nonzero in $F$, and let $E$ be an elliptic curve over $F$. Then the $n$-torsion $E(F)[n]$
  in the $F$-points of $E$ is a finite group isomorphic to $(\Z/n\Z)^2$.
\end{corollary}
\begin{proof}\uses{group_theory_lemma,WeierstrassCurve.n_torsion_card}\leanok
  This follows from the previous group-theoretic lemma~\ref{group_theory_lemma} and
  theorem~\ref{WeierstrassCurve.n_torsion_card}.
\end{proof}

We saw in section~\ref{twopointfour} that if $E$ is an elliptic curve over a field $k$ and if $k^{\sep}$ is a separable closure of~$k$, then the group $\Gal(k^{\sep}/k)$ acts on $E(k^{\sep})[n]$. Now let $n$ be a positive integer which is nonzero in $k$. We have just seen that $E(k^{\sep})[n]$ is isomorphic to $(\Z/n\Z)^2$, and it inherits an action of $\Gal(k^{\sep}/k)$. If we fix an isomorphism $E(k^{\sep})[n]\cong(\Z/n\Z)^2$ then we get a representation $\Gal(k^{\sep}/k)\to\GL_2(\Z/n\Z)$. A fundamental fact about this Galois representation is that its determinant is the cyclotomic character.

\begin{theorem}\label{Elliptic_curve_det_n_torsion}\uses{Elliptic_curve_n_torsion_2d} If $E$ is an elliptic curve over a field $k$, and $n$ is a positive integer which is nonzero in $k$, then the determinant of the 2-dimensional representation of $\Gal(k^{\sep}/k)$ on $E(k^{\sep})[n]$ is the mod $n$ cyclotomic character.
\end{theorem}
\begin{proof}
  This presumably should be done via the Weil pairing. I have not yet put any thought into a feasible way to formalise this.
\end{proof}

\section{Good reduction}

We give a brief overview of the theory of good and multiplicative reduction of elliptic curves.
For more details one can consult the standard sources such as~\cite{silverman1}. We stick with the low-level approach, thinking of elliptic curves as plane cubics; whilst we cannot do this forever, it will suffice for these initial results.

\begin{definition}\label{EllipticCurve.GoodReduction} Let $E$ be an elliptic curve over the field of fractions $K$ of a valuation ring $R$ with maximal ideal $\m$. We say $E$ has \emph{good reduction over $R$} if $E$ has a model with
coefficients in $R$ and the reduction mod $\m$ is still non-singular. If $E$ is an elliptic curve
over a number field $N$ and $P$ is a maximal ideal of its integer ring $\calO_N$, then one says that $E$ has \emph{good reduction at $P$} if $E$ has good reduction over the $\calO_{N,P}$, the localisation of $\calO_N$ at $P$.
\end{definition}

\begin{remark} From this point on, our Frey curves and Frey packages will use notation $(a,b,c,\ell)$, with $\ell\geq 5$ a prime number, rather than $p$. This frees up $p$ for use as another prime.
\end{remark}

\begin{lemma}\label{Frey_curve_good_reduction}\uses{FreyCurve, EllipticCurve.GoodReduction} If $E$ is the Frey curve $Y^2=X(X-a^\ell)(X+b^\ell)$ associated to a
  Frey package $(a,b,c,\ell)$, and if $p$ is a prime
  not dividing $abc$ (and in particular if $p>2$), then $E$ has good reduction at~$p$.
\end{lemma}
\begin{proof} The reduction mod $p$ of the equation defining the Frey curve is still a smooth
  plane cubic, because the three roots $0$, $a^\ell$ and $-b^\ell$ are distinct modulo $p$
  (note that the difference between $a^\ell$ and $-b^\ell$ is $c^\ell$).
\end{proof}

If $E$ is an elliptic curve over a number field $N$ and if $\rho$ is the representation
of $\Gal(\overline{N}/N)$ on the $n$-torsion of $E$ then $\rho$ is continuous and its image is finite,
so by the fundamental theorem of (infinite) Galois theory the representation factors through an
injection $\Gal(L/N)\to\GL_2(\Z/n\Z)$ where $L/N$ is a finite Galois extension of
number fields. One says that $\rho$ is \emph{unramified} at a maximal ideal $P$ of $\calO_N$
if the extension $L/N$ is unramified at $P$ (or in other words, if the factorization
of $P\calO_L$ into prime ideals is squarefree).

At some point we will need a theory of finite flat group schemes over an affine base. Here
is a working definition.

\begin{definition}\label{finite_flat_group_scheme} If $R$ is a commutative ring, then
  a \emph{finite flat group scheme} over $R$ is the spectrum of a commutative Hopf algebra $H/R$
  which is finite and flat as an $R$-module.
\end{definition}

(Probably this is not the correct definition in the non-Noetherian case; one should instead ask for locally free, which is equivalent in the Noetherian case and enables you to reduce to the Noetherian case in general)

Some facts we will need are:

\begin{theorem}\label{good_reduction_implies_unramified}\notready If $E$ is an elliptic curve over a number
  field $N$ and $E$ has good reduction at a maximal ideal $P$ of $\calO_N$, and if furthermore
  $n\not\in P$, then the Galois representation on the $n$-torsion of $E$ is unramified.
\end{theorem}
\begin{proof}
  One approach would be by showing that the $n$-torsion in the integral model of $E$ over $\calO_{N,P}$
  is an etale finite flat group scheme. There might be simpler approaches however. It's worth looking to see what Silverman does.
\end{proof}

\begin{theorem}\label{good_reduction_implies_flat}\uses{finite_flat_group_scheme}\notready If $E$ is an elliptic curve over a number field
  $N$ and $E$ has good reduction at a maximal ideal $P$ of $\calO_N$ containing the prime number $p$,
  then the Galois representation on the $p$-torsion of $E$ comes from a finite flat group scheme
  over the localisation $\calO_{N,P}$.
\end{theorem}
\begin{proof}
  Indeed, the kernel of the $p$-torsion on a good integral model is finite and flat.
  Checking this claim formally will probably involve a fair amount of work.
\end{proof}

\section{Multiplicative reduction}

\begin{definition}\label{EllipticCurve.MultiplicativeReduction} Let $E$ be an elliptic curve over the field of fractions $K$ of a valuation ring $R$ with maximal ideal $\m$. We say $E$ has \emph{multiplicative reduction over $R$} if $E$ has a model with coefficients in $R$ and which reduces mod $R/\m$ to a plane cubic with one singularity, which is an ordinary double point.
  We say that the reduction is \emph{split} if the two tangent lines at the ordinary double point
  are both defined over $R/\m$, and \emph{non-split} otherwise.
\end{definition}

  If $E$ is an elliptic curve
  over a number field $N$ and $P$ is a maximal ideal of its integer ring $\calO_N$, then one says that $E$ has \emph{multiplicative reduction at $P$} if $E$ has multiplicative reduction over the $\calO_{N,P}$, the localisation of $\calO_N$ at $P$.

\begin{lemma}\label{Frey_curve_mult_reduction}\uses{EllipticCurve.MultiplicativeReduction} If $E$ is the Frey curve $Y^2=X(X-a^\ell)(X+b^\ell)$ associated to a Frey
  package $(a,b,c,\ell)$, and if $p$ is an odd prime
  which divides $abc$, then $E$ has multiplicative reduction at~$p$.
\end{lemma}
\begin{proof} The hypothesis $p\mid abc$ implies that precisely two of the three roots $0$, $a^\ell$ and $-b^\ell$
  of the cubic are equal mod~$p$. Call $x\in\Z/p\Z$ this common value. Then the reduction mod $p$ of
  the curve is smooth away from the point $(x,0)$,
  and has an ordinary double point at $(x,0)$. Hence the Frey curve has
  multiplicative reduction at $p$.
\end{proof}
\begin{remark} If the third root reduces mod $p$ to $y\not=x$, then the reduction
  is split multiplicative iff $x-y$ is a square mod $p$. We shall not need this fact.
\end{remark}

\begin{lemma}\label{Frey_curve_mult_reduction_at_two}\uses{EllipticCurve.MultiplicativeReduction}  If $E$ is the Frey curve $Y^2=X(X-a^\ell)(X+b^\ell)$ associated to a Frey package
  $(a,b,c,\ell)$ then $E$ has multiplicative reduction at 2.
\end{lemma}
\begin{proof} Indeed, the change of variables $X=4X'$
  and $Y=8Y'+4X'$ transforms the equation to
  $64Y'^2+64X'Y'=64X'^3+16X'^2(b^\ell-a^\ell-1)-4X'a^\ell b^\ell$ and, because $\ell\geq5$,
  $b$ is even and $a=3$ mod 4, we see that the 64s cancel, giving an equation over $\Z$ which reduces mod 2 to
  $Y'^2+X'Y'=X'^3+cX'^2$ for some $c\in\{0,1\}$. This cubic is smooth away from an ordinary
  double point at $(0,0)$. Hence the Frey curve has multiplicative reduction at~2.
\end{proof}
\begin{remark} Note that $E$ has split multiplicative reduction iff $c=0$, which happens iff $a^\ell=7$ mod $8$. We shall not need this fact.
\end{remark}

In particular, the Frey curve associated to a Frey package is \emph{semistable} -- it has good or
 multiplicative reduction at all primes.

The main thing we need about elliptic curves with multiplicative reduction over nonarchimedean
local fields is the \emph{uniformisation theorem}, originally due to Tate.

\begin{theorem}\label{Tate_curve_uniformisation}\uses{EllipticCurve.MultiplicativeReduction}\notready If $E$ is an elliptic curve over a field
  complete with respect to a nontrivial nonarchimedean (real-valued) norm $K$ and if $E$ has split
  multiplicative reduction, then there is a Galois-equivariant injection
  $(K^{\sep})^\times/q^{\mathbb{Z}}\to E(K^{\sep})$, where $q\in K^\times$ satisfies
  $|q|=|j(E)|^{-1}$.
\end{theorem}
\begin{proof}
  See~\cite{silverman2}, Theorems V.3.1, Remark V.3.1.2 (we don't need surjectivity),
  and Theorem V.5.3. This is a lot of work and is a good target for breaking down into many smaller lemmas.
\end{proof}

\begin{corollary}\label{multiplicative_reduction_torsion} If $E$ is an elliptic curve
  over a field $K$ complete with respect to a nontrivial nonarchimedean (real-valued) norm
  and with perfect residue field, and if $E$ has multiplicative reduction, then there's
  an unramified character $\chi$ of $\Gal(K^{\sep}/K)$ whose square is 1, such that for
  all positive integers $n$ with $n\not=0$ in $K$, the
  $n$-torsion $E(K^{\sep})[n]$ is an extension of $\chi$ by $\epsilon\chi$, where $\epsilon$ is the
  cyclotomic character. Furthermore, the element of $K^\times/(K^\times)^\ell$ corresponding
  to this extension is given by the $q$-invariant of the curve.
\end{corollary}
\begin{proof}\uses{Tate_curve_uniformisation} After a quadratic twist we may assume that $E$ has split multiplicative reduction.
  The result then follows from the uniformisation theorem and an explicit computation.
  Note that even if we do not prove surjectivity of Tate's uniformisation, we still know
  that it's surjective on the $n$-torsion, because we know that there are $n^2$ points in the $n$-torsion of $E$ over $K^{\sep}$, and they are all accounted for by the $n$-torsion in $(K^{\sep})^\times/q^{\mathbb{Z}}$.
\end{proof}

\section{Hardly ramified representations}

We make the following definition; this is not in the literature but it is a useful concept for us.

\begin{definition}\label{hardly_ramified}\notready Let $\ell\geq5$ be a prime and let $V$ be a
  2-dimensional vector space over $\Z/\ell\Z$. A representation
  $\rho: \GQ\to \GL(V)$ is said to be \emph{hardly ramified} if it satisfies the following four axioms:
  \begin{enumerate}
  \item $\det(\rho)$ is the mod $\ell$ cyclotomic character;
  \item $\rho$ is unramified outside $2\ell$;
  \item The semisimplification of the restriction of $\rho$ to $\Gal(\Qbar_2/\Q_2)$ is unramified;
  \item The restriction of $\rho$ to $\GQl$ comes from a finite flat group scheme.
  \end{enumerate}
\end{definition}

We are interested in hardly ramified representations for several reasons. One is that by using some
deep theorems, we will be able to prove that all hardly ramified representations are
\emph{potentially automorphic}, which will give us our first foothold into the world of modular
forms. We shall come back to these ideas later. In the next section we shall be concerned with
the following rather simpler result, namely that the $\ell$-torsion in the Frey curve
associated to a Frey package $(a,b,c,\ell)$ is hardly ramified. The proof is standard;
for another reference, see Theorem~2.15 of~\cite{ddt}.

\section[The l-torsion in the Frey curve is hardly ramified.]{The $\ell$-torsion in the Frey curve is hardly ramified.}

Let $(a,b,c,\ell)$ be a Frey package, with associated Frey curve $E$ and mod $\ell$ Galois
 representation $\rho=E[\ell]$. We now work through a proof that $\rho$ is hardly ramified.

 \begin{theorem}\label{Frey_curve_good}\notready If $p\not=\ell$ is a prime not dividing $abc$ then
  $\rho$ is unramified at~$p$.
 \end{theorem}
 \begin{proof}\uses{Frey_curve_good_reduction,good_reduction_implies_unramified} Indeed, $E$ has good reduction at $p$, and hence $\rho$ is unramified at $p$
  by~\ref{good_reduction_implies_unramified}.
 \end{proof}

If however $p$ divides $abc$ then $E$ has multiplicative
reduction at $p$, and we can use the theory of the Tate curve to analyse $\rho$ at $p$.

\begin{theorem}\label{Frey_curve_j}\lean{FreyCurve.j}\leanok\uses{FreyCurve} If $(a,b,c,\ell)$ is a Frey package then the $j$-invariant of the corresponding Frey curve is $2^8(C^2-AB)^3/A^2B^2C^2$, where $A=a^\ell$, $B=b^\ell$ and $C=c^\ell$.
\end{theorem}
\begin{proof}
  \leanok
  Apply the explicit formula (presumably already in mathlib)
\end{proof}

\begin{corollary}
  \label{FreyCurve.j_valuation_of_bad_prime}
  \lean{FreyCurve.j_valuation_of_bad_prime}
  \uses{Frey_curve_j}\leanok
  If $(a,b,c,\ell)$ is a Frey package and the $j$-invariant of the corresponding Frey curve is $j$, and if $2<p\mid abc$, then the $p$-adic valuation $v_p(j)$ of $j$ is a multiple of $\ell$.
\end{corollary}
\begin{proof}
  \uses{Frey_curve_j}\leanok
  Indeed $p$ does not divide $2^8$ as $p>2$, and (using the notation of the previous theorem) $p$ does not divide $C^2-AB$ either, because it divides precisely one of $A$, $B$ and $C$. Hence $v_p(j)=-2v_p(a^\ell b^\ell c^\ell)=-2\ell v_p(abc)$ is a multiple of $\ell$.
\end{proof}

\begin{corollary}\label{Frey_curve_unram}\notready If $(a,b,c,\ell)$ is a Frey package, if $2<p\mid abc$
  is a prime with $p\not=\ell$, then the $\ell$-torsion in the Frey curve is unramified
  at $p$.
\end{corollary}
\begin{proof}
  \uses{multiplicative_reduction_torsion, FreyCurve.j_valuation_of_bad_prime}
  After an unramified quadratic twist we may assume the curve is split at $p$.
  The theory of the Tate curve tells us that the extension of $\Q_p$ cut out by the $\ell$-torsion
  of the Frey curve is $\Q_p(\mu_\ell,\sqrt[\ell]{q})$, with $\mu_\ell$ the set of $\ell$th roots
  of unity in $\Qpbar$. Because $\ell\not=p$ the extension
  $\Q_p(\mu_\ell)$ is unramified at $p$. And because $p\not=2$ divides $abc$,
  theorem~\ref{FreyCurve.j_valuation_of_bad_prime} shows us that the $j$-invariant
  of the Frey curve has $p$-adic valuation a multiple of $\ell$.
  Thus the extension can be written $\Q_p(\mu_\ell,\sqrt[\ell]{u})$,
  where $u\in\Q_p^\times$ is a unit. The extension is hence unramified (because, for example,
  Hensel's Lemma shows that the $\ell$th root of $u$ is in the maximal unramified
  extension of $\Q_p$).
\end{proof}

\begin{corollary}\label{frey_curve_unramified}\notready If $(a,b,c,\ell)$ is a Frey package, then the $\ell$-torsion in the Frey curve is unramified at all primes $p\not=2,\ell$.
\end{corollary}
\begin{proof}\uses{Frey_curve_good, Frey_curve_unram} Follows from~\ref{Frey_curve_good} and~\ref{Frey_curve_unram}.
\end{proof}

This analysis needs to be slightly modified if $p=2$, because the $j$-invariant of
the Frey curve may not have 2-adic valuation a multiple of $\ell$. We obtain the
following weaker result.

\begin{corollary}\label{frey_curve_at_2}\notready If $(a,b,c,\ell)$ is a Frey package, then the
  semisimplification of the restriction of the $\ell$-torsion $\rho$ in the associated Frey curve
  to $\Gal(\Qbar_2/\Q_2)$ is unramified.
\end{corollary}
\begin{proof}\uses{multiplicative_reduction_torsion} After a quadratic twist to make the curve have split multiplicative
  reduction, the theory of the Tate curve shows us that $\rho$ is an extension
  of the trivial character by the cyclotomic character. Hence the semisimplification of this representation is the direct sum of two
  unramified characters and is hence unramified.
\end{proof}

\begin{theorem}\label{Frey_curve_mod_ell_rep_at_ell}\uses{finite_flat_group_scheme}\notready Let $\rho$ be the $\ell$-torsion in the
  Frey curve associated to a Frey package $(a,b,c,\ell)$. Then the restriction of $\rho$ to $\GQl$ comes from a finite flat group scheme.
\end{theorem}
\begin{proof}\uses{good_reduction_implies_flat, multiplicative_reduction_torsion} The Frey curve either has good reduction at $\ell$ (case 1 of FLT) or multiplicative
  reduction at $\ell$ (case 2 of FLT). In the first case the $\ell$-torsion is finite and flat
  at $\ell$ by theorem~\ref{good_reduction_implies_flat}. In the second case the theory of the Tate
   curve shows that the $\ell$-torsion is (up to quadratic twist) an extension of
    the trivial character by the cyclotomic character corresponding (via Hilbert 90)
    to the $\ell$th power of an $\ell$-adic unit. This extension is known to be finite and flat;
  see for example Proposition~8.2 of~\cite{edix}. Note that the proof in~\cite{edix} uses
  fppf cohomology, although one can write down a much more elementary proof of this using
  arguments in~\cite{katz-mazur}.
\end{proof}

We have now proved the first main result of this chapter.

\begin{theorem}\label{frey_curve_hardly_ramified}\uses{hardly_ramified}\notready Let $\rho$ be the Galois representation on the
  $\ell$-torsion of the Frey curve coming from a Frey package $(a,b,c,\ell)$. Then $\rho$ is hardly
  ramified.
\end{theorem}
\begin{proof}\uses{Elliptic_curve_det_n_torsion, frey_curve_unramified, frey_curve_at_2, Frey_curve_mod_ell_rep_at_ell}
  This follows from the results above. The fact that $\ell\geq 5$ follows from the definition of
  a Frey package. The first condition is theorem~\ref{Elliptic_curve_det_n_torsion},
  and the second is theorem~\ref{frey_curve_unramified}. The third condition is
  theorem~\ref{frey_curve_at_2}, and the fourth is theorem~\ref{Frey_curve_mod_ell_rep_at_ell}.
\end{proof}

\section[The l-torsion in the Frey curve is irreducible.]{The $\ell$-torsion in the Frey curve is irreducible.}

We finish this chapter by showing that Mazur's theorem implies that the $\ell$-torsion in the Frey
 curve is irreducible. We start by stating Mazur's theorem.

\begin{theorem}\label{mazur}\notready Let $E$ be an elliptic curve over $\Q$. Then the torsion subgroup
  of $E$ has size at most 16.
\end{theorem}
\begin{proof}\notready
  This is the main theorem of~\cite{mazur-torsion}. Formalising this result will be a highly
  non-trivial project; note that this theorem is used in all known proofs of FLT, so there
  seems to be no way around it.
\end{proof}

Let $(a,b,c,\ell)$ be a Frey package, with associated Frey curve $E$ and mod $\ell$ Galois
 representation $\rho=E[\ell]$. We know that $\rho$ is 2-dimensional; let's suppose for a
 contradiction that that it is
 reducible, so in particular its semisimplification is the direct sum of two characters $\alpha$
 and $\beta$.

 The next two results are Lemme 6 on p307 of~\cite{serrepropgal}.

\begin{theorem}\label{Frey_characters_are_unramified}\notready With notation as above, the characters
  $\alpha$ and $\beta$ are unramified at $p$ for all primes $p\not=\ell$.
\end{theorem}
\begin{proof}\uses{frey_curve_unramified, frey_curve_at_2} We have seen in theorem~\ref{frey_curve_unramified} that $\rho$ is unramified at all
  primes $p\not=2,\ell$, so the
  characters $\alpha$ and $\beta$ are unramified at all such primes. If $p=2$ then the
  semisimplification of the restriction of $\rho$ to $\Gal(\Qbar_2/\Q_2)$ is unramified
  by corollary~\ref{frey_curve_at_2}, so $\alpha$ and $\beta$ are unramified at 2.
\end{proof}

\begin{remark} Does this innocuous-looking proof above use some form of the Brauer-Nesbitt theorem?
\end{remark}

\begin{theorem}\label{Frey_characters_at_ell}\notready One of $\alpha$ and $\beta$ is unramified at $\ell$.
\end{theorem}
\begin{proof}\uses{multiplicative_reduction_torsion, Frey_curve_mod_ell_rep_at_ell}
In the multiplicative case this follows immediately from the theory of the Tate curve.
In the good reduction case, the $\ell$-torsion is finite and flat at $\ell$ by theorem~\ref{Frey_curve_mod_ell_rep_at_ell}, so we now need to understand what such representations
look like. If the reduction is supersingular, then $\rho$ is necessarily irreducible,
contradicting our assumption. If however the reduction is ordinary, then the theory of the
canonical subgroup shows that the $\ell$-torsion is an extension of an unramified character by
an unramified twist of the cyclotomic character (see Proposition 11 on p273 of~\cite{serrepropgal}).
\end{proof}

\begin{corollary}\label{Frey_characters_trivial} One of $\alpha$ and $\beta$ is trivial.
\end{corollary}
\begin{proof}\uses{Frey_characters_are_unramified, Frey_characters_at_ell}
The previous two theorems show that one of $\alpha$ and $\beta$ is a character unramified at all
primes, and hence cuts out an extension unramified at all primes, so by Minkowski's theorem this character
is trivial.
\end{proof}

To summarise, we have shown the following.
\begin{theorem}\label{Frey_curve_reducible_structure}\notready If $\rho$ is reducible, then either $\rho$
  has a trivial 1-dimensional submodule or a trivial 1-dimensional quotient (here ``trivial'' means
  that the Galois group $\GQ$ acts trivially).
\end{theorem}
\begin{proof}\uses{Frey_characters_trivial} Follows from the above.
\end{proof}

We now split into two cases, depending on whether $\rho$ has a trivial submodule or a trivial quotient.

\begin{lemma}\label{Frey_curve_trivial_submodule} If $\rho$ has a trivial 1-dimensional submodule then the
  Frey curve has a non-trivial point of order $\ell$.
\end{lemma}
\begin{proof} Indeed, the trivial 1-dimensional submodule is a Galois-invariant subgroup of $E[\ell]$, so
  it corresponds to a Galois-stable point of order $\ell$.
\end{proof}

\begin{corollary}\label{Frey_curve_no_trivial_submodule} $\rho$ cannot have a trivial 1-dimensional submodule.
\end{corollary}
\begin{proof}\uses{mazur, Frey_curve_trivial_submodule}
  We have just seen that in this case, the Frey curve has a point of order $\ell$.
  It also has three points of order 2, meaning that its torsion subgroup has order at least
  $4\ell\geq 20$, contradicting Mazur's theorem~\ref{mazur}.
\end{proof}

It remains to rule out the case where $\rho$ is reducible and has a trivial quotient. To do this, we need
to quotient out $\rho$ by its 1-dimensional Galois-stable submodule.

\begin{theorem}\label{Elliptic_curve_quotient_by_finite_subgroup}\notready If $p$ is a prime and
  if $E$ is an elliptic curve over a field $K$ of characteristic not equal to $p$,
   and if $C\subseteq E(K^{\sep})[p]$ is a Galois-stable
  subgroup of order $p$, then there's an elliptic curve $E':=$``$E/C$'' over $K$ and an isogeny of elliptic
  curves $E\to E'$ over $K$ inducing a Galois-equivariant surjection $E(K^{\sep})\to E'(K^{\sep})$
  with kernel precisely $C$.
\end{theorem}
\begin{proof}
  Brian Conrad suggested the following approach, applicable as well for abelian schemes $A\to S$
  over a base.  Let $G$ be a finite locally free $S$-subgroup of $A$, say $G$ with constant
  rank $n > 0$ by working locally on the base, so $G$ is contained in $A[n]$.  Then
  $n: A \to A$ is the fppf quotient of the source by $A[n]$, so it expresses $A$ as an
  $A[n]$-torsor over itself.  The problem of building $A/G$ as an abelian scheme is then
  seen to be the “same” as that of constructing the quotient of this $A[n]$-torsor by the G-action.

  In other words, the problem them becomes one having nothing specific to do with abelian schemes,
  at the cost of working over a base (such as the original target $A$) even when $S$ was the
  spectrum of a field in the application. The question is now: for a finite locally free
  commutative $S$-group $H$ and a closed locally free $S$-subgroup $G$, build a reasonable quotient
  $H/G$. One approach is to look at the Cartier dual $H^\vee\to G^\vee$, show that it's faithfully
  flat, and then deduce that the Cartier dual of the kernel of this map does the job. Note that
  one input for this proof is the claim that inclusions of Hopf algebras over fields are flat,
  proved nicely in Waterhouse’s book.
\end{proof}

I suspect that the proof above is a no-go right now; there will presumably be a much easier proof of this result in Silverman though. Note also that this approach will not
give us a plane cubic, but rather a smooth proper group scheme so we would need Riemann-Roch
to turn it into a plane cubic, although it's unlikely that one will be able to prove Mazur's
theorem without developing all of this machinery and much more.

\begin{corollary}\label{Frey_curve_no_trivial_quotient} $\rho$ cannot have a trivial 1-dimensional quotient.
\end{corollary}
\begin{proof}\uses{Elliptic_curve_quotient_by_finite_subgroup, mazur} $\rho$ has a Galois-stable submodule $C$. The quotient curve $E/C$ now has
  a trivial 1-dimensional submodule, and also three points of order~2 (the images of the three
  2-torsion points in $E$). Hence the torsion subgroup of $E/C$ has order at least $4\ell\geq 20$,
  again contradicting Mazur's theorem.
\end{proof}

\begin{theorem}\label{Frey_curve_irreducible} The $\ell$-torsion in the Frey curve associated to a Frey package $(a,b,c,\ell)$ is irreducible.
\end{theorem}
\begin{proof}\uses{Frey_curve_reducible_structure, Frey_curve_no_trivial_submodule, Frey_curve_no_trivial_quotient,Elliptic_curve_n_torsion_2d} Follows from theorem~\ref{Frey_curve_reducible_structure}, corollary~\ref{Frey_curve_no_trivial_submodule}
  and corollary~\ref{Frey_curve_no_trivial_quotient}.
\end{proof}


%%FILE blueprint/src/chapter/ch03freyreduction.tex
\chapter{Reducibility of {\it p}-torsion of the Frey curve}\label{ch_freyreduction}

\section{Overview}

In chapter~\ref{ch_reductions} we reduced FLT, modulo a hard theorem from the 1970s,
to Theorem~\ref{Wiles_Frey}, the assertion that $p$-torsion in the Frey
curve is reducible. In this chapter we deduce this assertion from three more complex claims
about ``hardly ramified'' Galois representations. It is relatively straightforward
to reduce one of these three claims to a result of Fontaine proved in the 1980s in his
paper on the nonexistence of nontrivial abelian schemes over $\Z$. The other two
claims lie deeper, and their proofs use techniques initially developed by Wiles in
the 1990s.

\section{Hardly ramified representations}

Let $(a,b,c,p)$ be a Frey package (so in particular $p\geq5$ is prime and $a^p+b^p=c^p$),
let $E$ be the corresponding Frey curve over $\Q$, and let $\rho:\GQ\to\Aut(E(\Qbar)[p])$
be the 2-dimensional Galois representation on the $p$-torsion of~$E$. Recall that our goal
is to prove that $\rho$ is reducible.

What we need to leverage is the fact that $\rho$ has very little ramification. To give
a toy example before we start: if $K$ is a number field (i.e., a finite extension of $\Q$)
and if the extension $K/\Q$ is unramified at all primes, then an old theorem of
Minkowski tells us that $K=\Q$. We want to prove a theorem in a similar vein, namely
that if a 2-dimensional mod $p$ Galois representation is ``hardly ramified'', then it is reducible.
Below, we give a precise
definition of what it means for a continuous 2-dimensional representation $\GQ\to\GL_2(R)$
to be hardly ramified. Before we do that, we need to say precisely which topological rings~$R$
we will allow. We say that a topological ring is emph{pro-Artinian} if it is a projective
limit of Artin local rings each equipped with the discrete topology, and if it has the
projective limit topology. We are only concerned with local pro-Artinian rings with finite
residue field; such things can be checked to be the same thing as topological
local rings with finite residue field whose underlying topological space is profinite,
and such that additive translates of open ideals form a basis for the topology.
Let us call such rings ``coefficient rings'' for now.

\begin{remark} We make some remarks to orient the reader.
  \begin{itemize}
    \item Any complete local Noetherian ring with finite residue field is a coefficient ring,
      if the ring is equipped with the $\m$-adic topology where $\m$ is the maximal ideal.
      In this case, all powers of $\m$ are open.
    \item In particular finite fields, and integer rings of finite extensions
      of $\Q_p$, are coefficient rings.
    \item If $R$ is a coefficient ring then $R$ is isomorphic to the projective limit
      of the finite rings $R/I$ as $I$ runs over the open ideals of~$R$.
    \item A non-Noetherian example of a coefficient ring is the projective limit over $n$ of
the rings $\Z/p\Z[\varepsilon_1,\ldots,\varepsilon_n]/(\forall i,j,\varepsilon_i\varepsilon_j=0)$;
these rings are convenient to include as coefficient rings for technical reasons; they make
representability theorems easier.
    \item The category of coefficient rings is equivalent to the pro-category of the
      category of finite local rings.
    \item A coefficient ring is pseudocompact in the sense of Grothendieck. A pseudocompact
      local ring is however a more general concept as such a thing may have an infinite
      residue field and would thus not be profinite.
    \item If $R$ is a coefficient ring with residue field of characteristic $\ell$,
  then there is a unique continuous map $\Z_\ell\to R$. Indeed, it suffices to prove that there
  is a unique continuous map $\Z_\ell\to R/I$ for each
  open ideal~$I$, but $R/I$ is a finite local ring with residue field of characteristic $\ell$.
  $R/I$ is hence Artinian, so some power of the maximal ideal is zero by Nakayama. This means
  that $\ell^N=0$ for some sufficiently large $N$, and hence $R/I$ is a $\Z/\ell^N\Z$-algebra
  and thus admits admits a unique map from $\Z_\ell$.
    \item It will be more convenient to fix once and for all the integer $\mathcal{O}$ in a
  finite extension of $\Q_\ell$ and consider ``coefficient $\mathcal{O}$-algebras'', namely
  coefficient rings~$R$ equipped with a continuous map $\mathcal{O}\to R$ which is a local
  homomorphism inducing an isomorphism on residue fields.
  \end{itemize}
\end{remark}

Because a coefficient ring~$R$ with residue field of characteristic $\ell$ is naturally
a $\Z_\ell$-algebra, we can talk about the $\ell$-adic cyclotomic character $\GQ\to R^\times$.
We are now ready to define hardly ramified representations.

\begin{definition}
  \lean{GaloisRepresentation.IsHardlyRamified}
  \label{hardly_ramified}
  \leanok
  Let $R$ be a coefficient ring with finite residue field of characteristic $\ell\geq3$.
  Let $V$ be a finite free $R$-module of rank~2, equipped with the product topology. A
  continuous representation $\rho: \GQ\to \GL_R(V)$ is said to be \emph{hardly ramified} if it
  satisfies the following four conditions:
  \begin{enumerate}
  \item $\det(\rho):\GQ\to R^\times$ is the cyclotomic character;
  \item $\rho$ is unramified outside $2\ell$;
  \item The restriction of $\rho$ to $\Gal(\Qbar_2/\Q_2)$ is reducible (more precisely,
  there is a short exact sequence $0\to R\to V\to R\to 0$ which is stable
  under the action of $\Gal(\Qbar_2/\Q_2)$) and the Galois action on the 1-dimensional
  quotient is an unramified representation of $\Gal(\Qbar_2/\Q_2)$ whose square is trivial;
  \item The restriction of $\rho$ to $\GQl$ is flat, by which we mean that for all open
  ideals $I$ of $R$, the (finite image) representation $\rho$ mod $I:\GQl\to \GL_{R/I}(V/I)$
  comes from a finite flat group scheme.
  \end{enumerate}
\end{definition}

A well-known result, which basically goes back to Frey, is the following:

\begin{theorem}
  \lean{FreyCurve.torsion_isHardlyRamified}
  \label{Frey_curve_hardly_ramified}
  \uses{hardly_ramified}
  \leanok
  The $\ell$-torsion $\rho:\GQ\to\GL_2(\Z/\ell\Z)$ in the Frey curve associated to a Frey
  package $(a,b,c,\ell)$ is hardly ramified.
\end{theorem}
\begin{proof}
  This was well-known in the 1980s. A proof sketch is as follows.
  First note that $\ell\geq5>3$ by definition of a Frey package. Let $\rho$
  denote the Galois representation on the $\ell$-torsion of the Frey curve.
  The fact that $\rho$ is 2-dimensional is~Corollary
  III.6.4(b) of~\cite{silverman1}, and the fact that its determinant is
  cyclotomic is Proposition~III.8.3 of the same reference. These results hold for elliptic curves
  in general. The remaining claims are specific to the Frey curve and lie
  deeper. The fact that $\rho$ is unramified outside $2\ell$ is a consequence
  of (4.1.12) and (4.1.13) of~\cite{serreconj}. The fact that $\rho$ at 2
  has an unramified 1-dimensional quotient of order at most 2 follows from
  the fact that the Frey curve is semistable at~2 (see (4.1.5) of~\cite{serreconj})
  and the theory of the Tate curve. Finally, the claim that $\rho$ is flat at $\ell$
  is Proposition~5 and (4.1.13) of~\cite{serreconj}.
\end{proof}


Note that irreducibility and absolute irreducibility for hardly ramified mod $\ell$ representations
are the same, because our assumptions that $\ell\geq3$
and that the determinant is cyclotomic imply that the image of complex conjugation
has distinct eigenvalues defined over the ground field.

The key theorem about hardly ramified representations is the following.

\begin{theorem}
  \lean{FreyCurve.torsion_not_isIrreducible}
  \label{hardly_ramified_reducible}
  \leanok
  If $\ell\geq 3$ is a prime and $\rho:\GQ\to\GL_2(\Z/\ell\Z)$ is hardly ramified,
  then $\rho$ is reducible.
\end{theorem}

Note that this (deep) claim is a consequence of Serre's conjecture~\cite{serreconj},
now a theorem of Khare and Wintenberger~\cite{kwII}, and indeed we shall use
methods introduced by Khare and Wintenberger to prove this special case of
Serre's conjecture. Given this result, we can deduce Theorem~\ref{Wiles_Frey}
(which we restate here) easily:

\begin{theorem}
  \label{Wiles_Frey_again}
  \lean{FLT.Bosses.B4_proof}
  \leanok
  If $\overline{\rho}$ is the mod $p$ Galois representation associated to a Frey package $(a,b,c,p)$ then
  $\overline{\rho}$ is reducible.
\end{theorem}
\begin{proof}
  \uses{Frey_curve_hardly_ramified,hardly_ramified_reducible}
  Indeed, $\rho$ is hardly ramified by theorem~\ref{Frey_curve_hardly_ramified}
  and thus reducible by theorem~\ref{hardly_ramified_reducible}.
\end{proof}

Our job of reducing FLT to theorems of the 1980s is hence reduced to proving
Theorem~\ref{hardly_ramified_reducible}.

\subsection{Hardly ramified mod $p$ representations are reducible}

In this section we will state three theorems, from which Theorem~\ref{hardly_ramified_reducible}
easily follows.

Firstly, we claim that
an irreducible hardly ramified mod $\ell$ representation lifts to an $\ell$-adic representation.

\begin{theorem}
  \lean{GaloisRepresentation.IsHardlyRamified.lifts}
  \label{hardly_ramified_lifts}
  \uses{hardly_ramified}
  \leanok
  If $\ell\geq3$ is prime and $\overline{\rho}:\GQ\to\GL_2(\Z/\ell\Z)$
  is hardly ramified and irreducible, then there exists a finite extension~$K$ of $\Q_\ell$
  with integer ring~$\calO$ and maximal ideal $\mathfrak{m}$
  and a hardly ramified representation
  $\rho:\GQ\to\GL_2(\calO)$ whose reduction modulo~$\m$ is isomorphic to $\rho$.
\end{theorem}
\begin{proof}
  Omitted for now {\bf TODO}
\end{proof}

Next we claim that a hardly ramified $\ell$-adic representation ``spreads out'' to a compatible
family of hardly ramified $q$-adic representations for all odd primes $q$ (note that we have
not made a definition of a hardly ramified 2-adic representation).

\begin{theorem}
  \lean{GaloisRepresentation.IsHardlyRamified.mem_isCompatible}
  \label{hardly_ramified_spreads_out}
  \uses{hardly_ramified}
  \leanok
  If $\ell\geq3$ is prime, $K$ is a finite extension of $\Q_\ell$
  with integers $\calO$ and if $\rho:\GQ\to\GL_2(\calO)$ is a hardly ramified representation
  whose reduction is irreducible,
  then there exists a number field $M$ and, for each finite place $\mu$ of $M$
  of characteristic prime to 2, with completion $M_\mu$ having integer ring $R_\mu$,
  a hardly ramified semisimple representation $\rho_\mu:\GQ\to\GL_2(R_\mu)$ (by which we
  mean the generic fibre is semisimple), with the following properties:
  \begin{itemize}
    \item There is some $\lambda\mid\ell$ of $M$ such that $\rho_\lambda\cong\rho$,
      the isomorphism happening over some appropriate local field containing a copy
      of $M_\lambda$ and a copy of~$K$;
    \item If $\mu_1$ and $\mu_2$ are two finite places of $M$ with odd residue characteristics $m_1$
      and $m_2$, and if $p\nmid 2m_1m_2$ is prime, then $\rho_{\mu_1}$ and $\rho_{\mu_2}$
      are both unramified at~$p$ and the characteristic polynomials $\rho_{\mu_1}(\Frob_p)$
      and $\rho_{\mu_2}(\Frob_p)$ lie in $M[X]$ and are equal.
  \end{itemize}
\end{theorem}
\begin{proof}
  Omitted for now {\bf TODO}
\end{proof}

In particular, we can ``move'' from an irreducible hardly ramified mod $\ell$ representation
to a hardly ramified 3-adic representation, and hence to a hardly ramified mod 3 representation.

However, we can essentially completely classify the hardly ramified mod 3 Galois representations:

\begin{theorem}
  \lean{GaloisRepresentation.IsHardlyRamified.mod_three}
  \label{hardly_ramified_mod3_reducible}
  \uses{hardly_ramified}
  \leanok
  Suppose $k$ is a finite field of characteristic 3, and suppose
  $\overline{rho}:\GQ\to\GL_2(k)$ is hardly ramified. Then $\overline{\rho}$ is an extension
  of the cyclotomic character by the trivial representation.
\end{theorem}
\begin{proof}
  Omitted for now. {\bf TODO}
\end{proof}

And we can use this to essentially completely classify the hardly ramified 3-adic Galois
representations:

\begin{theorem}
  \lean{GaloisRepresentation.IsHardlyRamified.three_adic}
  \label{hardly_ramified_3adic_reducible}
  \uses{hardly_ramified, hardly_ramified_mod3_reducible}
  \leanok
  Suppose $L/\Q_3$ is a finite extension, with integer ring $\calO_L$, and suppose
  $\rho_3:\GQ\to\GL_2(\calO_L)$ is hardly ramified. Then (considered as a representation
  to $\GL_2(L)$) $\rho_3^{ss}=1\oplus\chi_3$
  where $1$ is the trivial character and $\chi_3$ is the 3-adic cyclotomic character.
\end{theorem}
\begin{proof}
  Omitted for now {\bf TODO}
\end{proof}

Theorem~\ref{hardly_ramified_reducible} (if $\ell\geq 3$ is a prime and
$\overline{\rho}:\GQ\to\GL_2(\Z/\ell\Z)$ is hardly ramified,
then $\overline{\rho}$ is reducible) is an easy consequence of these theorems,
as we now show.

\begin{proof}
  \proves{hardly_ramified_reducible}
  \uses{hardly_ramified_lifts,
    hardly_ramified_spreads_out,hardly_ramified_3adic_reducible}
  Assume for a contradiction that $\overline{\rho}$ is irreducible. By theorem~\ref{hardly_ramified_lifts},
  $\overline{\rho}$ lifts to a hardly ramified $\ell$-adic reprepresentation $\rho$. By
  theorem~\ref{hardly_ramified_spreads_out}, $\rho$ is part of a compatible family of
  $q$-adic Galois representations. By theorem~\ref{hardly_ramified_3adic_reducible},
  any 3-adic member $\rho_3$ of this family has semisimplification $1\oplus\chi_3$ and in particular
  for $p\nmid 6$ we have that the characteristic polynomial of $\rho_3(\Frob_p)=(X-p)(X-1).$
  By compatibility of the family we deduce that for $p\nmid 6\ell$ the characteristic
  polynomial of $\rho(\Frob_p)$ is $(X-p)(X-1)$, and thus the characteristic polynomial
  of $\overline{\rho}(\Frob_p)$ is $(X-p)(X-1)$. By the Cebotarev density theorem,
  $\overline{\rho}$ and $1\oplus\chi$ have the same characteristic polynomials everywhere
  (here $\chi$ is the mod $\ell$ cyclotomic character). Thus by the Brauer-Nesbitt theorem,
  $\overline{\rho}$ is reducible, the contradiction we seek.
\end{proof}

What remains then (modulo several results which were known in the 1980s),
is to prove the three theorems~\ref{hardly_ramified_lifts},
\ref{hardly_ramified_spreads_out} and~\ref{hardly_ramified_3adic_reducible}.
By far the easiest is theorem~\ref{hardly_ramified_3adic_reducible}; this follows
from old estimates of Fontaine (ultimately relying on bounds for root discriminants due to
Odlyzko and Poitou), originally developed to prove that there was no
nontrivial abelian scheme over $\Z.$ The other two theorems are deeper, and both use
modern variants of Wiles' $R=T$ machinery.

We have not yet written any more LaTeX on how to proceed further; the rest of
this blueprint should be considered as more unfocussed thoughts.


%%FILE blueprint/src/chapter/ch04overview.tex
\chapter{An overview of the proof}\label{ch_overview}

So far we have seen that, modulo Mazur's theorem (and various other things which will still take some work to formalise but which are much easier), Fermat's Last Theorem can be reduced
to the statement that there is no prime $\ell\geq 5$ and hardly-ramified
irreducible 2-dimensional Galois representation $\rho:\GQ\to\GL_2(\Z/\ell\Z)$.

In this chapter we give an overview of our strategy for proving this, and collect
various results which we will need along the way. Note that we no longer need to assume that $\rho$ comes from the $\ell$-torsion in an elliptic curve.

\section{Potential modularity.}

We will only speak about modularity for 2-dimensional representations of the
absolute Galois group of a totally real field $F$ \emph{of even degree over} $\Q$.
What we will mean by ``modular'' is ``associated to an automorphic representation of the
units of the totally definite quaternion algebra over $F$ ramified at no finite places
and at all infinite places''. Such a quaternion algebra only exists if $F$ has even degree,
hence the restriction in our definition of modularity.
We can furthermore even demand that the infinity type of the form is trivial (the analogue
in this setting of classical weight 2 forms), as these are the only forms we shall need for FLT.

Assume we have a hardly-ramified representation~$\rho$ as above. Let $K$
be the number field corresponding to the kernel of~$\rho$. Our first claim
is that there is some totally real field $F$ of even degree, Galois over $\Q$,
unramified at $\ell$, and disjoint from $K$, such that $\rho|G_F$ is modular.
The proof of this is very long, and uses a host of machinery. For example:
\begin{itemize}
    \item Moret--Bailly's result~\cite{moret-bailly} on points on curves with prescribed
    local behaviour;
    \item several nontrivial results in global class field theory;
    \item the Jacquet--Langlands correspondence;
    \item The assertion that irreducible 2-dimensional mod $p$ representations induced from a character are modular (this follows from converse theorems);
    \item A modularity lifting theorem.
\end{itemize}

Almost everything here dates back to the 1980s or before.
The exception is the modularity lifting theorem, which we now state explicitly.

\section{A modularity lifting theorem}

Suppose $\ell\geq5$ is a prime, that $F$ is a totally real field of even degree in which $\ell$
is unramified, and that~$S$ is a finite set of finite places of~$F$ not dividing~$\ell$. Write
$G_F$ for the absolute Galois group of~$F$.

If $v\in S$ then let $F_v$ denote the completion of~$F$ at~$v$, fix an inclusion $\overline{F}\to\overline{F_v}$,
let $\calO_v$ denote the integers of $F_v$ and $k(v)$ the residue field. Let $I_v\subset G_F$ denote the inertia
subgroup at~$v$. Local class field theory (or a more elementary approach) gives a map $I_v\to\calO_{F_v}^\times$
and hence a map $I_v\to k(v)^\times$. Let $J_v$ denote the kernel of this map.

Let $R$ be a complete local Noetherian $\Z_\ell$-algebra with finite residue field of characteristic $\ell$.
We will be interested in representations $\rho:G_F\to\GL_2(R)$ with the following four properties.
\begin{itemize}
    \item $\det(\rho)$ is the cyclotomic character;
    \item $\rho$ is unramified outside $S\cup\{\ell\}$;
    \item If $v\in S$ then $\rho(g)$ has trace equal to~2 for all $g\in J_v$;
    \item If $v\mid\ell$ is a place of $F$ then $\rho$ is flat at~$v$.
\end{itemize}

In the last bullet point, ``flat'' means ``projective limit of representations arising from
finite flat group schemes''. Let us use the lousy temporary notation ``$S$-good'' to denote
representations with these four properties.

Say $k$ is a finite extension of $\Z/\ell\Z$ and $\rhobar:G_F\to\GL_2(k)$ is continuous,
absolutely irreducible when restricted to $F(\zeta_\ell)$, and $S$-good. One can check
that the functor representing $S$-good lifts of $\rhobar$ is representable.

\begin{theorem}
    \label{modularity_lifting_theorem}
    \uses{Skinner_Wiles_CFT_trick,local_galois_coh_dim_two,local_galois_coh_top_degree,
      local_galois_coh_poincare,local_galois_coh_euler_poincare,IsDedekindDomain.FiniteAdeleRing.baseChangeAlgEquiv,
      automorphic_representation_local_decomposition,cuspidal_automorphic_representation,
      Galois_representation_from_automorphic_representation_on_GL_2_form, moret-bailly,
      local_galois_coh_finite}
      \notready
        If $\rhobar$ is modular of level $\Gamma_1(S)$ and $\rho:G_F\to\GL_2(\calO)$ is
        an $S$-good lift of $\rhobar$ to $\calO$, the integers of a finite extension of $\Q_\ell$,
        then $\rho$ is also modular of level $\Gamma_1(S)$.
\end{theorem}

Right now we are very far from even stating this theorem in Lean.

I am not entirely sure where to find a proof of this in the literature, although it has certainly been known to the experts for some time. Theorem 3.3 of~\cite{taylor-mero-cont} comes close, although it assumes that $\ell$ is totally split in $F$ rather than just unramified. Another near-reference is Theorem~5.2 of~\cite{toby-modularity}, although this assumes
the slightly stronger assumption that the image of $\rho$ contains $\SL_2(\Z/p\Z)$ (however it is well-known to the experts that this can be weakened to give the result we need). One reference for the proof is \href{https://math.berkeley.edu/~fengt/249A_2018.pdf}{Richard Taylor's 2018 Stanford course}.

\begin{proof} (Sketch)

The proof is a two-stage procedure and has a nontrivial analytic input. First one uses the Skinner--Wiles trick to reduce to the ``minimal case'', and this needs cyclic base change for $\GL(2)$ and also a characterisation
of the image of the base change construction; this seems to need a multiplicity one result, which (because of our
definition of ``modular'') will need Jacquet--Langlands as well.

In the minimal case, the argument is the usual Taylor--Wiles trick, using refinements due to Kisin and others.
\end{proof}

Given this modularity lifting theorem, the strategy to show potential modularity of $\rho$ is to use Moret--Bailly to find an appropriate totally real field $F$, an auxiliary prime $p$, and an auxiliary elliptic curve over $F$ whose mod $\ell$ Galois representation is $\rho$ and whose
mod $p$ Galois representation is induced from a character. By converse theorems (for example)
the mod $p$ Galois representation is associated to an automorphic representation of
$\GL_2/F$ and hence by Jacquet--Langlands it is modular. Now we use the
modularity lifting theorem to deduce the modularity of the curve over $F$ and hence
the modularity of the $\ell$-torsion.

\section{Compatible families, and reduction at 3}

We now use Khare--Wintenberger to lift $\rho$ to a potentially modular $\ell$-adic
Galois representation of conductor 2, and put it into an $\ell$-adic family using
the Brauer's theorem trick in \cite{blggt}. Finally we look at the 3-adic specialisation
of this family. Reducing mod 3 we get a representation which is flat at 3 and tame at 2,
so must be reducible because
of the techniques introduced in Fontaine's paper on abelian varieties over $\Z$ (an irreducible
representation would cut out a number field whose discriminant violates the Odlyzko bounds).
One can now go on to deduce that the 3-adic representation must be reducible, which
contradicts the irreducibility of $\rho$.

We apologise for the sketchiness of what is here, however at the time of writing it is so far from what we are even able to \emph{state} in Lean that there seems to be little point right now in fleshing out the argument further. As this document grows, we will add a far more detailed discussion of what is going on here. Note in particular that stating the modularity lifting theorem in Lean is the first target.


%%FILE blueprint/src/chapter/ch05automorphicformexample.tex
\chapter{An example of an automorphic form}

\section{Introduction}

The key ingredient in Wiles' proof of Fermat's Last Theorem is a \emph{modularity lifting theorem},
sometimes called an $R=T$ theorem. For Wiles, the $R$ came from elliptic curves, the $T$ came from
classical modular forms, and the fact that they're equal is basically the Shimura--Taniyama--Weil
conjecture, now known as the Breuil--Conrad--Diamond--Taylor modularity theorem: any elliptic
curve over the rationals is modular.

At the heart of the proof we shall formalise is also an $R=T$ theorem, however the $T$ which we
shall use will be associated not to classical modular forms, but to spaces of more general
automorphic forms called quaternionic modular forms. Those of you who know something about
classical modular forms might well know that the groups $\SL_2(\R)$ and $\SL_2(\Z)$ are
intimately involved; these are the norm 1 units in the matrix rings $M_2(\R)$ and $M_2(\Z)$.
In the theory of quaternionic modular forms, the analogous groups are the norm 1 units
in rings such as Hamilton's quaternions $\R\oplus\R i\oplus\R j\oplus\R k$, and subrings
such as $\Z\oplus\Z i\oplus\Z j\oplus\Z k$.

One of the main goals of the FLT project at the time of writing this sentence, is formalising
the \emph{statement} of the modularity lifting theorem which we shall use. So we are going
to need to develop the theory of quaternionic modular forms, which is rather different
to the theory of classical modular forms (for example, in the cases we need, the definition
is completely algebraic; there are no holomorphic functions in sight, and the analogue
of the upper half plane in the quaternionic theory is a finite set of points).

We could just launch into the general theory over totally real fields, which will be the generality
which we'll need. But when I was a PhD student, I learnt about
these objects by playing with explicit examples. So, whilst not logically necessary for the proof,
I thought it would be fun, and perhaps also instructional, to compute a concrete example of a space
of quaternionic modular forms. The process of constructing the example might even inform what kind
of machinery we should be developing in general. Let's begin by discussing the quaternion algebra
we shall use.

\section{A quaternion algebra}

Let's define $D$ to be the quaternion algebra $\Q\oplus\Q i\oplus\Q j\oplus\Q k$. As a vector
space, $D$ is 4-dimensional over $\Q$ with $[1,i,j,k]$ giving a basis. It has a (non-commutative)ring structure,
with multiplication satisfying the usual quaternion algebra relations $i^2=j^2=k^2=ijk=-1$.
You can think of~$D$ as an analogue of $2\times 2$
matrices with rational coefficients, hence its units $D^\times$ are an analogue of the
group $\GL_2(\Q)$.

We will also need an analogue of the group $\GL_2(\Z)$, which will come from an integral structure
on~$D$. We choose the Hurwitz order, namely the subring $\calO:=\Z\oplus\Z i\oplus\Z j\oplus\Z \omega$,
where $\omega=\frac{-1+(i+j+k)}{2}$, a cube root of unity, as $(i+j+k)^2=-3$. The simplest way
to understand $\calO$ is that it's quaternions $a+bi+cj+dk$ where either $a,b,c,d$ are all integers
or are all in $\frac{1}{2}+\Z$.

Note that $\calO$ is a maximal order and a Euclidean domain, which is why we prefer it over the
more obvious sublattice $\Z\oplus\Z i\oplus\Z j\oplus\Z k.$

In this chapter, we are going to compute a complex vector space which could be called something
like the ``weight 2 level 2 modular forms for $D^\times$''. The main result will be that this
space is 1-dimensional.

Note that mathlib has modular forms, but it doesn't have enough complex analysis to deduce
that the space of modular forms of a given weight and level is finite-dimensional. If all
the `sorry`s in this chapter are completed
before mathlib gets the necessary complex analysis, then the first nonzero space of modular forms
to be proved finite-dimensional in Lean will be a space of quaternionic modular forms.

We will use a modern "adelic" definition of our modular forms, so the first thing we need
to do is to talk about profinite completions.

\section{\texorpdfstring{$\Zhat$}{Zhat}}

Classically automorphic forms were defined as functions on symmetric spaces (like the upper half
plane) which transformed well under the action of certain discrete groups (for example $\SL_2(\Z)$).
However such definitions became combinatorially problematic when generalised to number fields
with nontrivial class group, because the classical theory needed a \emph{number} $p$ to define
the Hecke operator $T_p$, and in the case where $p$ was a non-principal prime ideal in a number
field, there was no appropriate number. One fix is to take disjoint unions of symmetric spaces
indexed by the ideal class group of the field in question, but it is easier to work adelically,
which is morally what we shall do. However we are able to avoid introducing the adeles
explicitly; we can work instead with the conceptually simpler object $\Zhat$,
the profinite completion of $\Z$. So what is $\Zhat$? We offer a low-level definition of this object.

Given an integer $z$, we can reduce it mod $N$ for every positive natural
number and get elements $z_N=\overline{z}\in\Z/N\Z$. These elements are not completely arbitrary
though -- they must satisfy some compatibility conditions. For example there can be no
positive integer $z$ such that $z_{10}=6$ and $z_2=1$, because $z_{10}=6$ tells us that
$z$ ends in a 6 when written in base 10, and in particular it's even, so $z_2$ must be~0.
The general rule: if $D\mid N$ then $z_D$ must be equal to image of $z_N$ under the natural
ring homomorphism from $\Z/N\Z$ to $\Z/D\Z$. We say that a collection of elements
$z_N\in\Z/N\Z$ is \emph{compatible} if it satisfies this rule.

\begin{definition}\label{ZHat}\lean{ZHat}\leanok The profinite completion $\Zhat$ of $\Z$ is the set of
    all compatible collections $c=(c_N)_N$ of elements of $\Z/N\Z$ indexed by $\N^+:=\{1,2,3,\ldots\}$.
    A collection is said to be \emph{compatible} if for all positive integers
    $D\mid N$, we have $c_N$ mod $D$ equals $c_D$.
\end{definition}

\begin{lemma}
    \label{ZHat.commRing}
    \lean{ZHat.commRing}
    \uses{ZHat}
    \leanok
    $\Zhat$ is a subring of $\prod_{N\geq1}(Z/N\Z)$ and in particular is a ring.
\end{lemma}
\begin{proof} \leanok Follow your nose.
\end{proof}

Examples of elements of $\Zhat$ are are given by integers, where we define $z_N$ to be $z$
mod $N$ for all $N$. This gives us a natural map from $\Z$ to $\Zhat$. In particular we
can talk about $0\in\Zhat$ and $1\in\Zhat$.

\begin{lemma}
    \label{ZHat.nontrivial}
    \lean{ZHat.nontrivial}
    \uses{ZHat.commRing}
    \leanok
    $0\not=1$ in $\Zhat$.
\end{lemma}
\begin{proof}
    \leanok
    Recall that you can evaluate an element of $\Zhat$ at a positive integer.
    Evaluating $0$ at 2 gives $0$, and evaluating $1$ at $2$ gives $1$, and these
    are distinct elements of $\Z/2\Z$, so $0\not=1$ in $\Zhat$.
\end{proof}

\begin{lemma}
    \label{ZHat.charZero}
    \lean{ZHat.charZero}
    \leanok
    The map from the naturals into $\Zhat$ sending $n$ to $n$ is injective.
\end{lemma}
\begin{proof}
    \leanok
    Generalise the above idea. Feel free to write up a LaTeX proof and PR it.
\end{proof}

Note that it follows easily that that the map from the integers to $\Zhat$ is injective.

But $\Zhat$
is \emph{much} larger than $\Z$; it has the same cardinality as the reals in fact.
Let's write down an explicit example of an element of $\Zhat$ which isn't obviously in $\Z$.

\begin{definition}
    \label{ZHat.e}
    \lean{ZHat.e}
    \uses{ZHat.commRing}
    \leanok
    The infinite sum $0!+1!+2!+3!+4!+5!+\cdots$ looks
    like it makes no sense at all; it is the sum of an infinite series of larger and larger
    positive numbers.
    However, the sum is \emph{finite} modulo $N$ for every positive integer $N$, because
    all the terms from $N!$ onwards are multiples of $N$ and thus are zero in $\Z/N\Z$.
    Thus it makes sense to define $e_N$ to be the value of the finite sum modulo $N$.
    Explicitly, $e_N=0!+1!+\cdots+(N-1)!$ modulo $N$.
\end{definition}

\begin{lemma}
    \label{ZHat.e_def}
    \lean{ZHat.e_def}
    \uses{ZHat.e}
    \leanok
    The collection $(e_N)_N$ is an element of $\Zhat$.
\end{lemma}
\begin{proof}
    \leanok
    This boils down to checking that $D!+(D+1)!+\cdots+(N-1)!$ is a multiple of~$D$.
\end{proof}

\begin{lemma}
    \label{ZHat.e_not_in_Int}
    \lean{ZHat.e_not_in_Int}
    \uses{ZHat.e}
    \leanok
    The element $(e_N)_N$ of $\Zhat$ is not in $\Z$.
\end{lemma}
\begin{proof}
    \leanok
    First imagine that $e=n$ with $n\in\Z$ and $0\leq n$. In this case, choose $j$
    such that $0!+1!+2!+\cdots+j!>n$ and check also that the sum is less
    than $(j+1)!$. Now set $N=(j+1)!$ and let's compare $e_N$
    and $n_N=n$. The trick is that $e_N$ must be $0!+1!+\cdots+j!$ mod $N$,
    because all the terms beyond this are multiples not just of $(j+1)$ but
    of $(j+1)!=N$. Thus mod $N$ we have $0\leq n<e_N<N$ so $n\not=e$.

    Now we deal with $n=-t<0$; choose $j$ large such that
    $(j+1)!-(0!+1!+\cdots+j!)>t$ (possible because the sum is at most $2\times j!$)
    and then set $N=(j+1)!$ and we have $0 < e_N<N-t<N$ so we cannot have $e_N=-t$ in $\Z/N\Z$,
    so again $e\not=n$.
\end{proof}

Let's prove some more basic facts about $\Zhat$.

\begin{lemma}
    \label{ZHat.torsionfree}
    \lean{ZHat.torsionfree}
    \uses{ZHat.commRing, ZHat.charZero}
    \leanok
    If $0<N$ is an integer then multiplication by $N$ is injective on $\Zhat$.
\end{lemma}
\begin{proof}
    \leanok
    Suppose that $(z_i)_i\in\Zhat$ and $Nz=0$. This means that $Nz_i=0\in\Z/i\Z$ for all $i$.
    Let us fix an arbitrary positive integer~$j$; we need to prove that $z_j=0\in\Z/j\Z$.
    Consider the element $z_{Nj}\in\Z/Nj\Z$. By assumption, we have $Nz_{Nj}=0$, meaning that
    if we lift $z_{Nj}$ to an integer, we have $Nj\mid Nz_{Nj}$, and thus $j\mid z_{Nj}$.
    Thus by the compatibility assumption on the $z_i$ we have that $z_j\in\Z/j\Z$ is the
    mod~$j$ reduction of $z_{Nj}$ and hence is zero.
\end{proof}

We will also need to understand exactly which elements of $\Zhat$
are multiples of~$N$.

\begin{lemma}
    \label{ZHat.multiples}
    \lean{ZHat.multiples}
    \uses{ZHat.commRing}
    \leanok
    The multiples of~$N$ in $\Zhat$ are precisely the compatible collections $(z_i)_i\in\Zhat$
    with $z_N=0$.
\end{lemma}
\begin{proof} \leanok
    Clearly $z_N=0$ is a necessary condition to be a multiple of~$N$. To see it is sufficient,
    take a general $(z_i)\in\Zhat$ such that $z_N=0$,
    and now define a new element $(y_j)_j$ of $\Zhat$
    by $y_j=z_{Nj}/N$. Just to clarify what this means: $z_{Nj}\in\Z/Nj\Z$ reduces mod~$N$
    to $z_N=0$ by the compatibility assumption, so it is in the subgroup $N\Z/Nj\Z$ of $\Z/Nj\Z$,
    which is isomorphic (via "division by $N$") to the group $\Z/j\Z$; this is how we construct
    $y_j$. It is easily checked that the $y_j$ are compatible and that $Ny=z$.
\end{proof}

\section{More advanced remarks on \texorpdfstring{$\Zhat$}{Zhat} versus \texorpdfstring{$\Q$}{Q}}

This section can be skipped on first reading.

People who have seen some more advanced algebra might recognise this construction of $\Zhat$
as being the profinite completion of the additive abelian group $\Z$, so it is a fundamental
object of mathematics in some sense. But usually, when building mathematics, after $\Z$ we
go to $\Q$, a multiplicative localisation of $\Z$, and only complete after that (to get $\R$).
The process of ``completing before localising'' gives us a far more arithmetic completion
of $\Z$.

Even though $\Q$ is a divisible abelian group and hence its profinite completion vanishes,
we can still attempt to "locally profinitely complete it" by defining $\Qhat:=\Q\otimes_{\Z}\Zhat$.
This object is more commonly known as the \emph{finite adeles} of $\Q$. More generally if $F$ is
any number field then $F\otimes_{\Z}\Zhat$ is the ring of finite adeles of $F$. To get to
the full ring of adeles of a number field~$F$ you need to take the product with the
ring of infinite adeles of $F$, which is $F\otimes_{\Q}\R$: some kind of universal
archimedean completion of $F$. I don't know a reference which develops the theory of adeles
in this way, so this is what we shall do here.

\section{\texorpdfstring{$\Qhat$}{Qhat} and tensor products.}

The definition of $\Qhat$ is easy if you know about tensor products
of additive abelian groups.

\begin{definition}
    \label{QHat}
    \lean{QHat}
    \uses{ZHat}
    \leanok
    % blue node
    The profinite completion $\Qhat$ of $\Q$ is the tensor product $\Q\otimes_{\Z}\Zhat$,
    or $\Qhat=\Q\otimes\Zhat$ for short.
\end{definition}

\section{A crash course in tensor products}

We've defined $\Qhat$ to be $\Q\otimes\Zhat$. Whatever does this mean? Well just to orient yourself,
if $A$ and $B$ are additive abelian groups, then $A\otimes B$ is also an abelian group.
And if $A$ and $B$ are commutative rings (as they are in our case), then $A\otimes B$ is also
a commutative ring.

Even if $A$ and $B$ are completely concrete commutative rings, their tensor product $A\otimes B$
might be incomprehensible. For example $\bbC\otimes\bbC$ is completely incomprehensible (note that we are
tensoring over the integers). It is not like the product of groups or the disjoint union of
two sets, where you have a completely explicit unambiguous formula for each element.

In this sense, the theory of tensor products is a bit like the theory of continuous functions.
Humanity started off studying concrete polynomial equations such as $x^2+1$ and then moved on to
concrete analytic functions such as $\log(x)$ and $\sin(x)$, but eventually the abstract concept
of a continuous function from the reals to the reals was born. There is no ``formula'' for a general
continuous function, and continuous functions such as $e^{-1/x^2}$ or $|x|$ have no power series.
Even if there \emph{were} a formula for a specific continuous function of interest, it is not clear
in general how to make sense of the claim that it's the ``best'' formula. In other words, there is
no "canonical form" for a general continuous function, and yet we prove things about them anyway.
We shall adopt the same attitude for elements of $A\otimes B$.

The first thing to know about the tensor product $A\otimes B$ of two abelian groups $A$ and $B$
is a ``constructor'' for the type. In other words, how can we make elements $A\otimes B$?
Well, it turns out that given elements $a\in A$ and $b\in B$, we can form the element
$a\otimes_t b\in A\otimes B$.

\begin{example} Recall that the sum of all the factorials is an element $e\in\Zhat$, and $22/7$ is certainly
a rational number, so we can make the element $\frac{22}{7}\otimes_te\in\Qhat$.
\end{example}

This example is in the Lean code.

Elements of the form $a\otimes_t b\in A\otimes B$ are known as \emph{pure tensors}. In the
literature, pure tensors
are often written $a\otimes b$, but we shall follow {\tt mathlib}'s convention in reserving
the $\otimes$ symbol for \emph{groups} like $A \otimes B$, and adorning it with a $t$ when
using it on \emph{elements} of the groups (or, as Lean calls them, \emph{terms}, which explains
the notation).

Addition of pure tensors obeys the ``distributivity'' rules
$a\otimes_t b_1+a\otimes_t b_2=a\otimes_t(b_1+b_2)$ and
$a_1\otimes_t b+a_2\otimes_t b=(a_1+a_2)\otimes_t b$, but there is no rule which simplifies
a general sum $a\otimes_t b + c\otimes_t d$ into a pure tensor. Indeed, in general it is \emph{not}
the case that every element of a tensor product $A\otimes B$ is of the form $a\otimes_t b$; there can be
tensors which aren't pure. However every element of $A\otimes B$ is a finite sum of pure tensors,
with the result that one can attempt to define additive maps from $A\otimes B$ by saying what they
do on pure tensors, and then extending linearly.

Another thing worth understanding is that just like how rational numbers can be written as
quotients of integers in several ways (for example $1/2=2/4=3/6=\cdots$), a general pure tensor
in $A\otimes B$ can be represented as $a\otimes_t b$ in many ways. For example, in $\Qhat$
we have $1\otimes_t 2=2\otimes_t 1$. A general rule for equality of pure tensors is that if
$a\in A$ and $b\in B$ and $z\in\Z$, then $za\otimes_tb=a\otimes_tzb$; integers can move over the
tensor symbol. But equality is hard: in general there may not be an algorithm to decide
whether two pure tensors $a\otimes_t b$ and $c\otimes_t d$ are equal in $A\otimes B$.

\begin{remark} A summary of the situation: if $A$ and $B$ are abelian groups, then every element
    of $A\otimes B$ can be written in the form $\sum_{i=1}^Na_i\otimes_tb_i$. It's just that this
    representation is highly nonunique, and furthermore given explicit elements $a_1,a_2\in A$
    and $b_1,b_2\in B$ it might be a hard problem to figure out if
    $a_1\otimes_t b_1=a_2\otimes_t b_2$.

    For example, it turns out that $(\Z/2\Z)\otimes(\Z/3\Z)=0$
    and so in this tensor product all the $a\otimes_t b$ are equal to each other and to $0\otimes 0$.
\end{remark}

Having said all of that, one nice property of $\Qhat$ is that every tensor \emph{is} pure;
let's prove this now.

\begin{lemma}
    \label{QHat.canonicalForm}
    \lean{QHat.canonicalForm}
    \uses{QHat,ZHat.commRing}
    \leanok
    Every element of $\Qhat:=\Q\otimes\Zhat$ can be written as $q\otimes_t z$ with $q\in\Q$ and $z\in\Zhat$.
    Furthermore one can even assume that $q=\frac{1}{N}$ for some positive integer $N$.
\end{lemma}
\begin{proof} \leanok
    A proof I would write on the board would look like the following. Take a general
    element of $\Qhat$; we know it can be expressed as a finite sum
    $\sum_i q_i\otimes_t z_i$ with $q_i\in\Q$ and $z_i\in\Zhat$. Now choose a large
    positive integer $N$, the lowest common multiple of all the denominators showing up in the
    $q_i$, and then rewrite $\sum_i q_i\otimes_t z_i$ as $\sum_i \frac{n_i}{N}\otimes z_i$ with
    $n_i\in\Z$. Now using the fundamental fact that $na\otimes_t b=a\otimes_t nb$ for $n\in\Z$,
    we can rewrite the sum as $\sum_i \frac{1}{N}\otimes_t n_i z_i$
    which is equal to the pure tensor $\frac{1}{N}\otimes(\sum_i n_i z_i)$.

    In Lean I would prove this using {\tt TensorProduct.induction\_on}, which quickly
    reduces us to the claim that the sum of two pure tensors is pure, which we can prove
    using the above technique whilst avoiding the general theory of finite sums.
\end{proof}

Be careful though: just because every element of $\Qhat$ can be written as $q\otimes z$, this
representation may not be unique. For example $2\otimes 1=1\otimes 2$. However, writing
$\frac{1}{N}\otimes_t z$ as $z/N$ does tempt us into the following definition.

\begin{definition}
    \label{QHat.IsCoprime}
    \lean{QHat.IsCoprime}
    \uses{ZHat.commRing}
    \leanok
    If $N\in\N^+$ and $z\in\Zhat$ then we say that $N$ and $z$ are \emph{coprime} if
    $z_N\in(\Z/N\Z)^\times$. We write $z/N$ as notation
    for the element $\frac{1}{N}\otimes_tz$.
\end{definition}

\begin{lemma}
    \label{QHat.lowestTerms}
    \lean{QHat.lowestTerms}
    \uses{QHat.IsCoprime}
    \leanok
    Every element of $\Qhat$ can be uniquely written as $z/N$ with $z\in\Zhat$, $N\in\N^+$,
    and with $N$ and $z$ coprime.
\end{lemma}
\begin{proof}
    Existence: by the previous lemma, an arbitrary element can be written as $z/N$; let $D$
    be the greatest common divisor of $N$ and $z_N$ (lifted to a natural). If $D=1$
    then the fraction is by definition in lowest terms. However if $1<D\mid N$ then $z_D$
    is the reduction of $z_N$ and is hence 0. By lemma~\ref{ZHat.multiples} we deduce that $z=Dy$
    is a multiple of~$D$, and hence $z/N=\frac{1}{N}\otimes_tDy=\frac{1}{E}\otimes y$, where
    $E=N/D$. Now if a natural divided both $y_E$ and $E$ then this natural would divide both $z_N/D$
    and $N/D$, contradicting the fact that $D$ is the greatest common divisor.

    Uniqueness: if $z/N=w/M$, we deduce $1\otimes_t Mz=1\otimes_t Nw$,
    and by injectivity of $\Zhat\to\Qhat$ we deduce that $Mz=Nw=y$.
    In particular, if $L$ is the lowest common multiple of $M$ and $N$ then $y_L$ is a multiple of both $M$ and $N$ and is
    hence zero, so $y=Lx$ is a multiple of~$L$ by~\ref{ZHat.multiples}, and we deduce
    from torsionfreeness that $z=(L/M)x$ and $w=(L/N)x$. If some prime divided $L/M$
    then it would have to divide~$N$ which means that $z$ is not in lowest terms;
    similarly if some prime divided $L/N$ then $w/M$ would not be in lowest terms.
    We deduce that $L=M=N$ and hence $z=w$ by torsionfreeness.
\end{proof}

If $A$ and $B$ are additive abelian groups then $A\otimes B$ is also an additive abelian group.
However if $A$ and $B$ are commutative rings, then $A\otimes B$ also inherits the structure
of a commutative ring, with $0=0\otimes_t 0$ and $1=1\otimes_t 1$. Pure tensors multiply in the
obvious way: the product of $a_1\otimes_t b_1$ and $a_2\otimes_t b_2$ is $a_1a_2\otimes_t b_1b_2.$
There are ring homomorphisms $A\to A\otimes B$ and $B\to A\otimes B$ sending $a$ to $a\otimes_t 1$
and $b$ to $1\otimes_t b$. In general such maps are not injective, but in the case of
$\Qhat=\Q\otimes\Zhat$ both maps from $\Q$ and $\Zhat$ are inclusions.

\begin{lemma}
    \label{QHat.injective_rat}
    \lean{QHat.injective_rat}
    \uses{QHat,ZHat.commRing}
    \leanok
    The ring homomorphism $\Q\to\Qhat$ sending $q$ to $q\otimes_t 1$
    is injective.
\end{lemma}
\begin{proof} \leanok
    We have seen that the map from $\Z$ to $\Zhat$ is
    injective. Now $\Q$ is a flat $\Z$-module, because it's
    torsion-free, so tensoring up we deduce that the map
    from $\Q=\Q\otimes\Z$ to $\Qhat=\Q\otimes\Zhat$ is also injective.
    There is no doubt a more elementary proof of this fact.
\end{proof}

\begin{lemma}
    \label{QHat.injective_zHat}
    \lean{QHat.injective_zHat}
    \uses{QHat, ZHat.torsionfree}
    \leanok
    The ring homomorphism $\Zhat\to\Qhat$ sending
    $z$ to $1\otimes_t z$ is injective.
\end{lemma}
\begin{proof} \leanok
    The map from $\Z$ to $\Q$ is injective, and we have seen
    that $\Zhat$ is a torsion-free and thus flat $\Z$-module,
    so the map from $\Zhat$ to $\Qhat$ is also injective.
\end{proof}

We can thus identify $\Q=\Q\otimes\Z$ and $\Zhat=\Z\otimes\Zhat$ with subrings of $\Qhat=\Q\otimes\Zhat$.
Note that, being commutative rings, $\Q$ and $\Zhat$ both contain a copy of $\Z$ as a subring, and
the corresponding copies of $\Z$ in $\Qhat$ are equal; this is because $1\otimes a=a\otimes 1$
for all $a\in\Z$.

\section{Additive structure of \texorpdfstring{$\Qhat$}{Qhat}.}

Here we forget the ring structure on everything, and analyse $\Qhat$ as an additive
abelian group, and in particular how the subgroups $\Z$, $\Q$ and $\Zhat$ sit within it.

The two results we prove in this section are that $\Q\cap\Zhat=\Z$ and
that $\Q+\Zhat=\Qhat$. Using lattice-theoretic notation we could write
these results as $\Q\sqcap\Zhat=\Z$ and $\Q\sqcup\Zhat=\Qhat$.

\begin{lemma}
    \label{QHat.rat_meet_zHat}
    \lean{QHat.rat_meet_zHat}
    \uses{QHat,ZHat.commRing}
    \leanok
    The intersection of $\Q$ and $\Zhat$ in $\Qhat$ is $\Z$.
\end{lemma}
\begin{proof}
    \uses{QHat.lowestTerms}
    \leanok
    Clearly $\Z\subseteq\Q\cap\Zhat$. Now suppose that $x\in\Q\cap\Zhat$.
    Because $x$ is rational we can write it as $\frac{A}{B}\otimes_t1$ for some
    fraction $A/B$ in lowest terms, and hence $x=A/B$ where now we regard $A\in\Zhat$
    and note that $A/B$ is still in lowest terms. However $x\in\Zhat$ implies that
    $x=x/1$ is in lowest terms, so we deduce that $B=1$ and thus $x=A\in\Z$.
\end{proof}

\begin{lemma}
    \label{QHat.rat_join_zHat}
    \lean{QHat.rat_join_zHat}
    \uses{QHat,ZHat.commRing}
    \leanok
    The sum of $\Q$ and $\Zhat$ in $\Qhat$ is $\Qhat$.
    More precisely, every element of $\Qhat$ can be written as $q+z$ with $q\in\Q$ and $z\in\Zhat$,
    or more precisely as $q\otimes_t 1+1\otimes_t z$.
\end{lemma}
\begin{proof}
    \uses{QHat.canonicalForm}
    \leanok
    Write $x\in\Qhat$ as $x=z/N$ in lowest terms. Lift $z_N$ to an integer $t$ and observe
    that $(z-t)_N=0$, hence $z-t=Ny$ for some $y\in\Zhat$. Now $x=t/N+y\in\Q+\Zhat$.
\end{proof}

\section{Multiplicative structure of the units of \texorpdfstring{$\Qhat$}{Qhat}.}

We now forget the additive structure on the commutative ring $\Qhat$ and consider
the multiplicative structure of its group of units $\Qhat^\times$ (which I couldn't
get into the section title). We have the obvious
subgroups $\Q^\times$, $\Z^\times$ and $\Zhat^\times$.

\begin{lemma}
    \label{Qhat.unitsrat_meet_unitszHat}
    \lean{QHat.unitsrat_meet_unitszHat}
    \uses{QHat,ZHat.commRing}
    \leanok
    The intersection of $\Q^\times$ and $\Zhat^\times$ in $\Qhat^\times$ is $\Z^\times$.
\end{lemma}
\begin{proof}
    \uses{QHat.lowestTerms,QHat.rat_meet_zHat}
    \leanok
    Clearly the intersection is contained within $\Q\cap\Zhat=\Z$. If $n\in\Z$ is in $\Zhat^\times$
    then $n\not=0$ and its inverse $1/n=\pm1/|n|$ is in lowest terms but also in $\Zhat$,
    and hence $|n|=1$ by uniqueness of lowest term representation.
\end{proof}

\begin{lemma}
    \label{QHat.unitsrat_join_unitszHat}
    \lean{QHat.unitsrat_join_unitszHat}
    \uses{QHat,ZHat.commRing}
    \leanok
    The product of $\Q^\times$ and $\Zhat^\times$ in $\Qhat^\times$ is all of $\Qhat^\times$.
    More precisely, every element of $\Qhat^\times$ can be written as $qz$ with $q\in\Q^\times$
    and $z\in\Zhat^\times$.
\end{lemma}
Note that by the previous lemma, this representation will be unique up to sign.
\begin{proof}
    \uses{ZHat.multiples,ZHat.eq_zero_of_mul_eq_zero,QHat.canonicalForm}
    \leanok
  We already know that a general element of $\Qhat^\times$ can be written as $x/N$ with $N$
  positive, so this reduces us to proving that a general element $x\in\Zhat$ which is invertible
  in $\Qhat^\times$ can be written as $qz$ with $q\in\Q^\times$ and $z\in\Zhat^\times$.

  We know $1/x$ can be written in lowest terms as $y/M$, and multiplying up we deduce
  that $xy=M$, and hence $x$ divides a positive integer. If $i:\Z\to\Zhat$ denotes
  the inclusion, then we've just seen that the preimage of the principal
  ideal $(x)$, namely, $J:=i^{-1}(x\Zhat)$ is nonzero, as it contains $M$.
  Let $g\in J$ be the smallest positive integer; it's well-known that $J=(g)$.

  I claim that it suffices to show that $x\Zhat=g\Zhat$. Because knowing $g=yx$ and
  $x=gz$ for some $y,z\in\Zhat$ tells us that $g(1-yz)=0$, and we know that multiplication by $g$ is injective,
  hence $yz=1$, so $z$ is a unit and we have written $x=gz$ with $g\in\Q^\times$ and $z\in\Zhat^\times$.

  It remains to prove the claim. By definition $g\in J\subseteq x\Zhat$ so this is one
  inclusion. For the other, it suffices to prove that $x_g=0$. However if $0<x_g<g$
  lifts $x_g$ to the naturals then I claim that $x_g\in J$, for $x_g-x$ is a multiple
  of~$g$ and hence of~$x$, and this contradicts minimality of~$g$.
\end{proof}

We are nearly ready to embark upon the multiplicative adelic theory for quaternion algebras.
However before we do this, we need to develop the theory of the Hurwitz quaternions a bit
more formally.

\section{The Hurwitz quaternions}

\begin{definition}
    \label{Hurwitz}
    \lean{Hurwitz}
    \leanok
        The Hurwitz quaternions are the set $\calO := \Z\oplus\Z \omega\oplus\Z i\oplus
    \Z i\omega$ (as an abstract abelian group or as a subgroup of the usual quaternions).
    Here $\omega=\frac{-1+(i+j+k)}{2}$ and note that $(i+j+k)^2=-3$.
    We have $\overline{\omega}=\omega^2=-(\omega+1)$. A general quaternion
    $a+bi+cj+dk$ is a Hurwitz quaternion if either $a,b,c,d\in\Z$
    or $a,b,c,d\in\Z+\frac{1}{2}$.
\end{definition}

\begin{lemma}
    \label{Hurwitz.ring}
    \lean{Hurwitz.ring}
    \uses{Hurwitz}
    \leanok
    The Hurwitz quaternions form a ring.
\end{lemma}
\begin{proof}
    \leanok
    Follow your nose.
\end{proof}


This ring is isomorphic to $\Z^4$ as an additive group, and
$\calO\otimes_{\Z}\R=\R\oplus \R i\oplus \R j\oplus\R \omega$
is the usual Hamilton quaternions.

\begin{definition}
    \label{Hurwitz.starRing}
    \lean{Hurwitz.starRing}
    \uses{Hurwitz.ring}
    \leanok
    There's a conjugation map (which we'll call "star") from the Hurwitz quaternions to themselves, sending
    integers to themselves and purely imaginary elements like $2\omega+1$ to minus themselves.
    It satisfies $(x^*)^*=x$, $(xy)^*=y^*x^*$ and $(x+y)^*=x^*+y^*$. In particular, the Hurwitz
    quaternions are a "star ring" in the sense of mathlib.
\end{definition}

\begin{definition}
    \label{Hurwitz.norm}
    \lean{Hurwitz.norm}
    \uses{Hurwitz}
    \leanok
    The Hurwitz quaternions come equipped with an integer-valued norm, which is
    $a^2+b^2+c^2+d^2$ on $a+bi+cj+dk$ but needs to be modified a bit to deal with $\omega$.
\end{definition}

\begin{lemma}
    \label{Hurwitz.norm_eq_mul_conj}
    \lean{Hurwitz.norm_eq_mul_conj}
    \leanok
    We have $N(x)=x\overline{x}$.
\end{lemma}
\begin{proof}
    \uses{Hurwitz.norm, Hurwitz.starRing}
    \leanok
    Easy calculation.
\end{proof}

\begin{lemma}
    \label{Hurwitz.norm_zero}
    \lean{Hurwitz.norm_zero}
    \uses{Hurwitz.norm}
    \leanok
    The norm of $0$ is $0$.
\end{lemma}
\begin{proof}
    \leanok
    A calculation.
\end{proof}

\begin{lemma}
    \label{Hurwitz.norm_one}
    \lean{Hurwitz.norm_one}
    \uses{Hurwitz.norm}
    \leanok
    The norm of $1$ is $1$.
\end{lemma}
\begin{proof}
    \leanok
    A calculation.
\end{proof}

\begin{lemma}
    \label{Hurwitz.norm_mul}
    \lean{Hurwitz.norm_mul}
    \uses{Hurwitz.norm}
    \leanok
    The norm of a product is the product of the norms.
\end{lemma}
\begin{proof}
    \leanok
    A calculation.
\end{proof}

\begin{lemma}
    \label{Hurwitz.norm_nonneg}
    \lean{Hurwitz.norm_nonneg}
    \uses{Hurwitz.norm}
    \leanok
    The norm of an element is nonnegative.
\end{lemma}
\begin{proof} \leanok
    It's a sum of rational squares.
\end{proof}

\begin{lemma}
    \label{Hurwitz.norm_eq_zero}
    \lean{Hurwitz.norm_eq_zero}
    \uses{Hurwitz.norm}
    \leanok
    The norm of an element is zero if and only if the element is zero.
\end{lemma}
\begin{proof}
    \leanok
    It's a sum of rational squares.
\end{proof}

\begin{lemma}
    \label{Hurwitz.exists_near}
    \lean{Hurwitz.exists_near}
    \uses{Hurwitz.norm}
    \leanok
    Given a ``usual'' quaternion $a=x+yi+zj+wk$ with $x,y,z,w\in\R$,
    there exists a Hurwitz quaternion $q$ such that $N(a-q)<1$.
\end{lemma}
\begin{proof}
    \leanok
  If $[r]$ denotes the nearest integer to the real number $r$, then $|r-[r]|\leq \frac{1}{2}$.
  Hence if $q=[x]+[y]i+[z]j+[w]k$ then $N(a-q)=|x-[x]|^2+\cdots
  \leq \frac{1}{4}+\frac{1}{4}+\frac{1}{4}+\frac{1}{4}\leq 1$, with strict inequality unless
  $|x-[x]|=|y-[y]|=|z-[z]|=|w-[w]|=\frac{1}{2}$, in which case $a\in\mathcal{O}$ because $a-\omega$
  has integer coordinates.
\end{proof}

\begin{lemma}
    \label{Hurwitz.quot_rem}
    \lean{Hurwitz.quot_rem}
    \uses{Hurwitz.norm, Hurwitz.exists_near}
    \leanok
    Given two Hurwitz quaternions $a$ and $b$ with $b$ nonzero, there exists
    $q$ and $r$ such that $a=qb+r$ and $N(r)<N(b)$.
\end{lemma}
\begin{proof}
    \leanok
  Let $q$ be the Hurwitz quaternion obtained by applying Lemma~\ref{Hurwitz.exists_near}
  to $a/b := ab^{-1}$; then $N(a/b-q)<1$ and now everything follows after multiplying up.
\end{proof}

\begin{corollary}
    \label{Hurwitz.left_ideal_princ}
    \lean{Hurwitz.left_ideal_princ}
    \uses{Hurwitz.norm}
    \leanok
    All left ideals of $\calO$ are principal.
\end{corollary}
\begin{proof}
    \uses{Hurwitz.quot_rem}
    \leanok
    If the ideal is 0, use 0. Otherwise, choose a nonzero element of smallest norm.
\end{proof}

\begin{remark}
    All right ideals are principal too, because there's
    another version of Euclid saying $a=bq+r$.
\end{remark}

\section{Profinite completion of the Hurwitz quaternions}

We define $\calOhat$ to be $\calO\otimes\Zhat$, so it's elements $a+bi+cj+d\omega$
with $a,b,c,d\in\Zhat$. The basic thing we need is this:

\begin{theorem}
    \label{Hurwitz.surjective_pnat_quotient}
    \uses{Hurwitz}
    If $N$ is a positive natural then the obvious map $\calO\to\calOhat/N\calOhat$ is surjective.
\end{theorem}
\begin{proof}
    This is just four copies of the surjection $\Z\to\Zhat/N\Zhat$.
    Note that this latter map is surjective because $\Z\to\Z/N\Z$ is surjective,
    hence given $z\in\Zhat$ you can subtract an integer $w$ such
    that $(z-w)_N=0$, so $z-w$ is a multiple of~$N$.
\end{proof}

We define $D:=\Q\otimes\calO=\Q\oplus\Q i\oplus\Q j\oplus\Q\omega=\Q\oplus\Q i\oplus\Q j\oplus\Q k$.
Finally, we define $\widehat{D}:=D\otimes\Zhat$. Just as with $\Qhat$ we have

\begin{lemma}
    \label{HurwitzRatHat.canonicalForm}
    \lean{HurwitzRatHat.canonicalForm}
    \uses{QHat, ZHat, Hurwitz}
    \leanok
    Every element of $\widehat{D}$ can be written as $z/N$ with $z\in\calOhat$ and $N\in\N^+$.
\end{lemma}
\begin{proof}
    Same as the proof for $\Qhat$.
\end{proof}

It is not hard to check that $\widehat{D}$ contains $\widehat{\calO}$
and $D$ as subrings, and that as additive abelian groups we have
$\widehat{\calO}\cap D=\calO$ and $\widehat{\calO}+D=\widehat{D}$.
This is because $\calO$ is just four copies of $\Z$ and we've proved
the analogous result for $\Z$.

However the multiplicative structure is more interesting, especially
as $D$ is not commutative. For a general quaternion algebra it is \emph{not}
true that $(\widehat{D})^\times=D^\times(\widehat{\calO})^\times$, because
there are "class group obstructions". The double coset space is some
kind of non-commutative analogue of a class group. However for our
particular choice of $D$ and $\calO$ the result is true.

\begin{theorem}
    \label{HurwitzRatHat.completed_units}
    \lean{HurwitzRatHat.completed_units}
    \uses{HurwitzRatHat.canonicalForm}
    \leanok
    The group of units of $\widehat{D}$ is $D^\times\calOhat^\times$.
    More precisely, every element of $\widehat{D}^\times$ can be written as a product $\delta u$
    with $\delta\in D^\times$ and $u\in\calOhat^\times$.
\end{theorem}
\begin{proof}
    Given an element $x$ of $\widehat{D}^\times$, we can use lemma~\ref{HurwitzRatHat.canonicalForm}
    to write it as $z/N$ with $N$ a positive integer
    and $z\in\widehat{\calO}$. Note that $N$ is central and in $D^\times$. Similarly, we can
    write $x^{-1}$ as $y/M$ with $M$ a positive integer and $y\in\widehat{\calO}$. Then
    $1=xx^{-1}=zy/NM$ and so $zy=NM=MN$, and $1=x^{-1}x=yz/MN$ so $yz=MN$ too. In particular $y$
    both left and right divides a positive integer.

    Now consider the left ideal $\widehat{\calO}y$ generated by $y$. We've just seen that
    this ideal has nontrivial intersection with $\calO$, because it contains $MN>0$. Hence
    its intersection with $\calO$ is a nonzero left ideal of $\calO$, which is
    hence principal by corollary~\ref{Hurwitz.left_ideal_princ}. Write it as $\calO\alpha$
    with $0\not=\alpha\in\calO$.

    It suffices to show that $\calOhat\alpha=\calOhat y$. For this would imply that
    $u\alpha=y$ and $vy=\alpha$ for some $u,v\in\calOhat$ and thus $(vu-1)\alpha=0$
    and $(uv-1)y=0$, and both $\alpha$ and $y$ are left divisors of positive integers
    (the norm of $\alpha$, and $MN$ respectively), so now using the fact that $\calOhat$
    is $\Z$-torsion-free (is the tensor product of torsion-free abelian groups torsion-free? That
    would be a cheap way of doing it. Otherwise use $\calO=\Z^4$) we deduce that $u$ and $v$
    are units, and thus $x^{-1}=\frac{1}{M}u\alpha$ so $x=(M\alpha^{-1})v\in D^\times\calOhat^\times$.

    What remains is this. We have $y\in\calOhat$ which left and right divides some positive integer.
    We've defined $0\not=\alpha\in\calO$ such that $\calO\alpha$ is the pullback of the
    abelian group $\calOhat y$ along the map $\calO\to\calOhat$. We need to show that when we push
    this ideal $\calO\alpha$ forwards to $\calOhat$ we get $\calOhat y$ again. The fact that
    $\calOhat\alpha\subseteq\calOhat y$ is easy, because $\alpha\in\calOhat y$ by definition.
    So it remains to show that $y\in\calOhat\alpha$.

    Let's define $T$ to be a positive integer which is both a left and right multiple of
    both $y$ and $\alpha$ (for example $T=MN\alpha\overline{\alpha}$ will do). Now note that we
    have an isomorphism $\calO/T\calO=\calOhat/T\calOhat$, so we can choose
    some $\beta\in\calO$ such that $\beta-y\in T\calOhat$ is a multiple of $T$. Next note that
    $\beta\in y+\calOhat T\subset \calOhat y$ is in $\calOhat y\cap\calO=\calO\alpha$, meaning
    $\beta=\gamma\alpha$ for some $\gamma\in\calO$.
    Hence $y\in\beta+\calOhat T\subseteq\calOhat\alpha$.
\end{proof}


%%FILE blueprint/src/chapter/ch06automorphicrepresentations.tex
\chapter{Stating the modularity lifting theorems}

I think that a nice and accessible goal (which will maybe take a month or two) would be to
\emph{state} the modularity lifting theorems which we'll be formalising. There are in fact two;
one (the "minimal case") is proved using an extension of the original Taylor--Wiles techniques, and
the other is deduced from it using various more modern tricks which were developed later. This
chapter (currently work in progress) will contain a detailed discussion of all the things
involved in the statement of the theorem.

\section{Automorphic forms and analysis}

Modular forms were historically the first nontrivial examples of automorphic forms, but by the
1950s or so it was realised that they were special cases of a very general notion of an automorphic
form, as were Dirichlet characters! Modular forms are holomorphic automorphic forms for the group
$\GL_2/\Q$, and Dirichlet characters are automorphic forms for the group $\GL_1/\Q$. It's possible
to make sense of the notion of an automorphic form for the group $G/k$. Here $k$ is
a ``global field'' -- that is, a field which is either a finite extension of $\Q$ (a number field)
or a finite extension of $(\Z/p\Z)(T)$ (a function field), and $G$ is a connected reductive
group variety over $k$.

The reason that the definition of a modular form involves some analysis (they are holomorphic
functions) is that if you quotient out the group $\GL_2(\R)$ by its centre and the maximal compact
subgroup $O_2(\R)$, you get something which can be naturally identified with the upper half plane, a
symmetric space with lots of interesting differential operators associated to it (for example a
Casimir operator). However if you do the same thing with $\GL_1(\R)$ then you get a one point set,
which is why a Dirichlet character is just a combinatorial object; it's a group homomorphism
$(\Z/N\Z)^\times\to\bbC^\times$ where $N$ is some positive integer. It turns out that there are many
other connected reductive groups where the associated symmetric space is 0-dimensional, and in these
cases the definition of an automorphic form is again combinatorial. An example would be the group
variety associated to the units of a totally definite quaternion algebra over a totally real field.
In this case, the analogue of $\GL_2(\R)$ would be the units $\bbH^\times$ in the Hamilton
quaternions, a maximal compact subgroup would be the quaternions of norm 1 (homeomorphic to the
3-sphere $S^3$) and quotienting out $\bbH^\times$ by its centre $\R^\times$ and $S^3$ again just
gives you 1 point.

Before we talk about quaternion algebras, let's talk about central simple algebras.

\section{Central simple algebras}

Convention: in this section, fields are commutative, but algebras over a field may not be.

Recall that a \emph{central simple algebra} over a field $K$ is a
nonzero $K$-algebra $D$ such that $K$ is the
centre of $D$ and that $D$ has no nontrivial two-sided ideals.

Another way of saying that $D$ has no nontrivial two-sided ideals: every surjective ring
homomorphism $D\twoheadrightarrow A$ to any ring~$A$ is either an isomorphism, or the zero map
to the zero ring. Note that this latter condition
has nothing to do with~$K$.

\begin{lemma}
    \label{MatrixRing.isCentralSimple}
    %\lean{MatrixRing.isCentralSimple}
    %\uses{IsCentralSimple}
    \discussion{47}
    \leanok
    If $n\geq1$ then the $n\times n$ matrices $M_n(K)$ are a central simple algebra over~$K$.
\end{lemma}
\begin{proof}
We prove more generally that matrices with coefficients in~$K$ and indexed by an arbitrary nonempty
finite type are a central simple algebra over~$K$.

They are clearly an algebra over $K$, with $K$ embedded via scalar matrices as usual
(the injectivity of the map from~$K$ comes from nonemptiness of the finite index type).
The centre clearly contains $K$; to show that it
equals~$K$, we argue as follows. Let $e(i,j)$ be the matrix with a 1 in the $i$th row and $j$th
column, and zeros everywhere else. An element $Z=(Z_{s,t})_{s,t}$ of the centre commutes with
all matrices $e(i,j)$ for $i\not=j$ and these equations immediately imply that $Z_{i,j}=0$ if
$i\not=j$ and that $Z_{i,i}=Z_{j,j}$.

It suffices to prove that any nonzero two-sided ideal~$I$ is all of $M_n(K)$. So say $0\not=M\in I$
and let's fix $(i,j)$ such that $M_{i,j}\not=0$. One easily checks that
$M_{i,j} \mathrm{id} = \sum_{k}e(k,i)\times M\times e(j,k)\in I$ (where $\mathrm{id} \in M_n(K)$
is the identity matrix). Therefore, $\mathrm{id} \in I$, so $I = M_n(K)$.

The definition also requires that the ring be non-zero, but this follows from the index type being
nonempty.
\end{proof}

\begin{lemma}
    \label{IsCentralSimple.baseChange} % no Lean yet because Lean didn't seem to know L \otimes_K D was a ring
    %\uses{IsCentralSimple}
    If $D$ is a central simple algebra over~$K$ and $L/K$ is a field extension, then $L\otimes_KD$
    is a central simple algebra over~$L$.
\end{lemma}
\begin{proof}
    This is not too hard: it's lemma b of section 12.4 in Peirce's "Associative algebras".
\end{proof}


Next: define trace and norm.


%\begin{lemma}
%    \label{IsQuaternionAlgebra.exists_quaternionAlgebra_iso}
%    \lean{IsQuaternionAlgebra.exists_quaternionAlgebra_iso}
%  Any such quaternion algebra has a basis $1,i,j,k$ with $i^2=a$, $j^2=b$ and $ij=-ji=k$.
%\end{lemma}
%\begin{proof}
%  Choose $d\in D\backslash K$. Then
%\end{proof}


%%FILE blueprint/src/chapter/ch07exampleGLn.tex
\chapter{Automorphic forms and the Langlands Conjectures}

This chapter came from discussions between Patrick, Mario and myself, all currently
visiting the Hausdorff Research Institute for
Mathematics in Bonn. The ultimate goal is to formally state some version of the global Langlands
reciprocity conjectures for $\GL_n$ over $\Q$.

\section[Definition of an automorphic form]{Definition of an automorphic form for $\GL_n$ over $\Q$.}

The global Langlands reciprocity conjectures relate automorphic forms to Galois representations.
The statements for a general connected reductive group involve the construction of the Langlands
dual group, and we do not have quite enough Lie algebra theory to push this definition through
in general. However if we restrict the special case of the group $\GL_n/\Q$, the dual group
is just $\GL_n(\bbC)$ and several other technical obstructions are also removed. In this
section we will explain the definition of an automorphic form for the group $\GL_n/\Q$, following
the exposition by Borel and Jacquet in Corvallis.

\section{The finite adeles of the rationals.}

Mathlib already has the definition of the finite adeles $\A_{\Q}^f$ of the rationals as a
commutative $\Q$-algebra, and the proof that it's a topological ring.

\section[The adelic general linear group]{The group $\GL_n(\A_{\Q})$.}

The adeles $\A_{\Q}$ of $\Q$ are the product $\A_{\Q}^f\times\R$, with the product topology.
They are a topological ring. Hence $\GL_n(\A_{\Q})=\GL_n(\A_{\Q}^f)\times\GL_n(\R)$ is a
topological group, where we are being a bit liberal with our use of the equality symbol.

\section{Smooth functions}

\begin{definition}
  \label{AutomorphicForm.GLn.IsSmooth}
  \lean{AutomorphicForm.GLn.IsSmooth}
  A function $f:\GL_n(\A_{\Q}^f)\times\GL_n(\R)\to\bbC$ is \emph{smooth}
  if it has the following three properties.

  \begin{enumerate}
    \item $f$ is continuous.
    \item For all $x\in\GL_n(\A_{\Q}^f)$, the function $y\mapsto f(x,y)$ is smooth.
    \item For all $y\in\GL_n(\R)$, the function $x\mapsto f(x,y)$ is locally constant.
  \end{enumerate}
\end{definition}

Current state of this definition: I've half-formalised it; I don't know how to say the the
function is smooth on the infinite part, because I have never used the manifold library before
and I have no idea what my model with corners is supposed to be.

\section{Slowly-increasing functions}

Automorphic representations satisfy a growth condition which we may as well factor out
into a separate definition.

We define the following temporary ``size'' function $s:\GL_n(\R)\to\R$ by
$s(M)=trace(MM^T+M^{-1}M^{-T})$ where $M^{-T}$ denotes inverse-transpose. Note that
$s(M)$ is always positive, and is large if $M$ has a very large or very small
(in absolute value) eigenvalue.

\begin{definition}
  \label{AutomorphicForm.GLn.IsSlowlyIncreasing}
  \lean{AutomorphicForm.GLn.IsSlowlyIncreasing}
  \leanok
  We say that a function $f:\GL_n(\R)\to\bbC$ is \emph{slowly-increasing}
  if there's some real constant $C$ and positive integer $n$ such that $|f(M)|\leq Cs(M)^n$
  for all $M\in\GL_n(\R)$.
\end{definition}

Note: the book says $n$ is positive, but $\{M|s(M)\leq 1\}$ is compact so I don't
think it makes any difference.

\section{Weights at infinity}

\begin{definition}
  \label{AutomorphicForm.GLn.Weight}
  \lean{AutomorphicForm.GLn.Weight}

The \emph{weight} of an automorphic form for $\GL_n/\Q$ can be thought of as a finite-dimensional
continuous complex representation $\rho$ of a maximal compact subgroup of $\GL_n(\R)$,
and it's convenient to choose one (they're all conjugate) so we choose $O_n(\R)$.
\end{definition}

The Lean definition is incomplete right now -- I don't demand irreducibility
(I wasn't sure whether I was doing this the right way; if I used category theory
then I might have struggled to say that the representation was continuous).


\section{The action of the universal enveloping algebra.}

\begin{definition}
  \label{instLieAlgebraAction}
There is a natural action of the real Lie algebra of $\GL_n(\R)$ on the complex vector space of
smooth complex-valued functions on $\GL_n(\R)$.
\end{definition}

\begin{definition}
  \label{instComplexLieAlgebraAction}
  \uses{instLieAlgebraAction}

  This extends to is a natural complex Lie algebra action of the complexification of
  the real Lie algebra, on the smooth complex functions on $\GL_n(\R)$.

\end{definition}

\begin{definition}
  \label{instUniversalEnvelopingAlgebraAction}
  \uses{instComplexLieAlgebraAction}
By functoriality, we get an action of the universal enveloping algebra of this
complexified Lie algebra on the smooth complex functions.

\end{definition}

\begin{definition}
\label{instCentreAction}
\uses{instUniversalEnvelopingAlgebraAction}
Thus the \emph{centre} $\Z_n$ of this universal enveloping algebra also acts on the smooth
complex functions.

\end{definition}

\begin{remark} The centre we just defined is a commutative ring which contains a copy of $\bbC$.
  Note that Harish-Chandra, or possibly this was known earlier, showed that it is a polynomial
  ring in $n$ variables over the complexes. We shall not need this.
\end{remark}

\section{Automorphic forms}

From here on there is no more Lean right now, only LaTeX.

\begin{definition}
  \label{AutomorphicForm.GLn.AutomorphicFormForGLnOverQ}
  \lean{AutomorphicForm.GLn.AutomorphicFormForGLnOverQ}
  \uses{AutomorphicForm.GLn.IsSmooth, AutomorphicForm.GLn.IsSlowlyIncreasing,
    AutomorphicForm.GLn.Weight, instCentreAction}
  A smooth function $f:\GL_n(\A_{\Q}^f)\times\GL_n(\R)\to\bbC$ is
  an $O_n(\R)$-\emph{automorphic form} on $\GL_n(\A_{\Q})$ if it satisfies the following
  five conditions.
  \begin{enumerate}
    \item (periodicity) For all $g\in\GL_n(\Q)$, we have $f(gx,gy)=f(x,y)$.
    \item (has a finite level) There exists a compact open subgroup $U\subseteq\GL_n(\A_{\Q}^f)$
      such that $f(xu,y)=f(x,y)$ for all $u\in U$, $x\in \GL_n(\A_{\Q}^f)$ and $y\in\GL_n(\R)$.
    \item (weight $\rho$) There exists a continuous finite-dimensional irreducible complex
    representation $\rho$ of $O_n(\R)$ such that for every $(x,y)\in\GL_n(\A_{\Q})$, the
    set of functions $k\mapsto f(x,yk)$ span a finite-dimensional complex vector space isomorphic
    as $O_n(\R)$-representation to a direct sum of copies of $\rho$.
    \item (has an infinite level) There is an ideal $I$ of the centre $Z_n$ described in the
    previous section, which has finite complex codimension, and which annihiliates the
    function $y \mapsto f(x,y)$ for all $x\in \GL_n(\A_{\Q}^f)$. Note that this is a very fancy
    way of saying ``the function satisfies some natural differential equations''. In the
    case of modular forms, the differential equations are the Cauchy-Riemann equations, which
    is why modular forms are holomorphic.
    \item (growth condition) For every $x\in\GL_n(\A_{\Q}^f)$, the function $y\mapsto f(x,y)$
  on $\GL_n(\R)$ is slowly-increasing.
  \end{enumerate}

\end{definition}

Automorphic forms of a fixed weight $\rho$ form a complex vector space, and if we also
fix the finite level $U$ and the infinite level $I$ then we get a subspace which is
finite-dimensional; this is a theorem of Harish-Chandra. There is also the concept
of a cusp form, meaning an automorphic form for which furthermore some adelic integrals
vanish.

\section{Hecke operators}

\begin{lemma} The group $\GL_n(\A_{\Q}^f)$ acts (on the left) on the space of automorphic forms
  for $\GL_n(\A_{\Q})$ by the formula $(g\cdot f)(x,y)=f(xg,y)$.
\end{lemma}
\begin{proof}
  This is obvious. Note that the conjugate of a compact open subgroup is still
  compact and open.
\end{proof}

A formal development of the theory of Hecke operators looks like the following.

Let $U$ be a fixed compact open subgroup of $\GL_n(\A_{\Q}^f)$, and let's also fix
a weight $\rho$, and let $M_\rho(n)$ denote the complex vector space of automorphic
forms for $\GL_n/\Q$ of weight $\rho$. The level $U$ forms $M_\rho(n,U)$ are just the $U$-invariants
of this space. If $g\in\GL_n(\A_{\Q}^f)$, then I
claim that the double coset space $UgU$ can be written as a \emph{finite} disjoint union
of single cosets $g_iU$; one way of saying this is that the double coset space is certainly
a disjoint union of left cosets, but the double coset space is compact and the left cosets
are open.

Define the Hecke operator $T_g:M_\rho(n,U)\to M_\rho(n,U)$ by
$T_g(f)=\sum g_i\cdot f$.

\begin{lemma} This function is well-defined, i.e., it sends a $U$-invariant form to
  a $U$-invariant form which is independent of the choice of $g_i$.
\end{lemma}
\begin{proof} Easy group theory.
\end{proof}


%%FILE blueprint/src/chapter/chtopbestiary.tex
\chapter{Appendix: A collection of results which are needed in the proof.}\label{ch_bestiary}

In this (temporary, unorganised) appendix we list a whole host of definitions and theorems which were known to humanity by the end of the 1980s and which we shall need. These definitions and theorems will find their way into more relevant sections of the blueprint once I have written more details. Note that some of these things are straightforward; others are probably multi-year research projects. The purpose of this chapter right now is simply to give the community (and possibly AIs) some kind of idea of the task we face. Note also that many of the \emph{definitions} here are yet to be formalised in Lean, and this needs to be done before we can start talking about formalising theorems.

\section{Results from class field theory}

We start with the local case. In fact we restrict to the $p$-adic case, but only for simplicity of exposition because it's all we'll need (and, to be frank, because I'm not 100 percent of what is true in the function field case).

Let $K$ be a finite extension of $\Q_p$. We write $\widehat{\Z}$ for the profinite completion of $\Z$; it is isomorphic to $\prod_p\Z_p$ where $\Z_p$ is the $p$-adic integers and the product is over all primes.

\begin{theorem}\label{maximal_unramified_extension_of_p-adic_field}\notready The maximal unramified extension $K^{un}$ in a given algebraic closure of $K$
    is Galois over $K$ with Galois group ``canonically'' isomorphic to $\widehat{\Z}$ in two ways; one of these two isomorphisms identifies $1\in\widehat{\Z}$ with an arithmetic Frobenius (the endomorphism inducing $x\mapsto x^q$ on the residue field of $K^{un}$, where $q$ is the size of the residue field of $K$). The other identifies 1 with geometric Frobenius (defined to be the inverse of arithmetic Frobenius).
\end{theorem}

It is impossible to say which of the two canonical isomorphisms is ``the most canonical''; people working in different areas make different choices in order to locally minimise the number of minus signs in their results.

As a result, the absolute Galois group of $K$ surjects onto $\widehat{\Z}$; its kernel is said to be the \emph{inertia subgroup} of this Galois group. Now pull back this surjection along the continuous map from $\Z$ (with its discrete topology) to $\widehat{\Z}$, in the category of topological groups. We end up with a group containing the inertia group as an open normal subgroup, and with quotient isomorphic to the integers.

\begin{definition}\label{local_Weil_group}\uses{maximal_unramified_extension_of_p-adic_field}\notready The topological group described above is called the \emph{Weil group} of $K$.
\end{definition}

The following theorem is nontrivial.

\begin{theorem}\label{local_class_field_theory}\uses{local_Weil_group}\notready If $K$ is a finite extension of $\Q_p$ then there are two ``canonical'' isomorphisms of topological abelian groups, between $K^\times$ and the abelianisation of the Weil group of $K$.
\end{theorem}
\begin{proof} This is the main theorem of local class field theory; see for example the relevant articles in~\cite{cf} or many other places.
\end{proof}

Note that Mar\'ia In\'es de Frutos Fern\'andez and Filippo Nuccio are working on a formalisation of the proof of this using Lubin--Tate formal groups.

Now let $M$ be an abelian group (with the discrete topology) equipped with a continuous
action of $G_K$, the Galois group $\GK$ where we fix an algebraic closure $\Kbar$ of $K$. Note that
if one doesn't want to choose an algebraic closure of $K$ one can instead think of $M$ as being an
etale sheaf of abelian groups on $\Spec(K)$.

Continuous group cohomology $H^i(G_K,M)$ in this setting can be defined using continuous cocycles and continuous coboundaries, or just as a colimit of usual group cohomology over the finite quotients of this absolute Galois group (or as etale cohomology, if you prefer). Here are some of the facts we will need about cohomology in this situation. A nice summary is that cohomology of a local Galois group behaves like the cohomology of a compact connected 2-manifold. All the theorems below will need extensive planning.

\begin{theorem}
    \label{local_galois_coh_finite}
    \notready
    If $M$ is finite then the cohomology groups $H^i(G_K,M)$ all finite.
\end{theorem}
\begin{proof}
    This is Proposition~14 in section~5.2 of~\cite{serre-galcoh}.
\end{proof}

\begin{theorem} ["the dimension is 2"]\label{local_galois_coh_dim_two}\notready
  If $M$ is torsion then $H^i(G_K,M)=0$ if $i>2$.
\end{theorem}
\begin{proof} This follows from Proposition~15 in~section 5.3 of~\cite{serre-galcoh}.
\end{proof}

\begin{theorem} ["top degree"]
    \label{local_galois_coh_top_degree}
    \notready
    $H^2(G_K,\mu_n)$ is ``canonically'' isomorphic to $\Z/n\Z$.
\end{theorem}
\begin{proof} This is also included in Lemma 2 of section 5.2 of \cite{serre-galcoh} (Serre just writes that the groups are equal; he clearly is not a Lean user. I can see no explanation in his book of this use of the equality symbol. When the statement of this ``theorem'' is formalised in Lean it may well actually be a definition, giving the map).
\end{proof}

\begin{theorem}\notready There is a ``canonical'' isomorphism $H^2(K,\mu_\infty)=\Q/\Z$.
\end{theorem}
\begin{proof}\notready
This is in Theorem II.5.2 in~\cite{serre-galcoh}.
\end{proof}

\begin{theorem} ["Poincar\'e duality"]\label{local_galois_coh_poincare}\notready
    If $\mu=\bigcup_{n\geq1}\mu_n$ and $M':=\Hom(M,\mu)$ is the dual of $M$ then for
    $0\leq i\leq 2$ the cup product pairing $H^i(G_K,M)\times H^{2-i}(G_K,M')\to H^2(G_K,\mu)=\Q/\Z$ is perfect.
\end{theorem}
\begin{proof}\notready This is Theorem 2 in section 5.2 in \cite{serre-galcoh}. Note again
    the dubious (as far as Lean is concerned) use of the equality symbol.
\end{proof}

\begin{theorem} ["Euler-Poincar\'e characteristic"]\label{local_galois_coh_euler_poincare}\notready
    If $h^i(M)$ denotes the order of $H^i(G_K,M)$ then $h^0(M)-h^1(M)+h^2(M)=0$.
\end{theorem}

We now move onto the global case. If $N$ is a number field, that is, a finite extension of $\Q$,
then let $\A_N^f:= N\otimes_{\Z}\widehat{\Z}$
denote the finite adeles of $N$ and let $N_\infty := N\otimes_{\Q}\R$ denote the product of the completions of $N$ at the infinite places, so $\A_N:=\A_N^f\times N_\infty$ is the ring of adeles of $N$.

\begin{theorem}\label{global_class_field_theory}\uses{local_class_field_theory}\notready If $N$ is a finite extension of $\Q$ then there are two ``canonical'' isomorphisms of topological groups between the profinite abelian groups $\pi_0(\A_N^\times/N^\times)$ and $\GN^{\ab}$; one sends local uniformisers to arithmetic Frobenii and the other to geometric Frobenii; each of the global isomorphisms is compatible with the local isomorphisms above.
\end{theorem}
\begin{proof}\notready This is the main theorem of global class field theory; see for example Tate's article in~\cite{cf}.
\end{proof}

We need the following consequence:

\begin{theorem}\label{Skinner_Wiles_CFT_trick}\uses{global_class_field_theory}\notready Let $S$ be a finite set of places of a number field $K$ . For each $v \in S$
let $L_v/K_v$ be a finite Galois extension. Then there is a finite solvable Galois extension
$L/K$ such that if $w$ is a place of $L$ dividing $v \in S$, then $L_w/K_v$ is isomorphic to $L_v/K_v$ as $K_v$-algebra. Moreover, if $K^{\avoid} /K$ is
any finite extension then we can choose $L$ to be linearly disjoint from $K^{\avoid}$.
\end{theorem}

We also need Poitou-Tate duality; I'll refrain from writing it down for now, because we don't even have Galois cohomology in Lean yet (although we are very close to it).

\section{Structures on the points of an affine variety.}

All rings and algebras in this section are commutative with a 1, and all morphisms send 1 to 1.

Let $X=\Spec(A)$ be an affine scheme of finite type over a field $K$. For example $X$ could be an affine algebraic variety; in fact we shall only be interested in smooth affine varieties in the applications, but the initial definition and theorem are fine for all finite type schemes.

If $R$ is any $K$-algebra then one can talk about the $R$-points $X(R)$ of $X$, which in this case
naturally bijects with the $K$-algebra homomorphisms from $A$ to $R$.

\begin{definition}\label{topology_on_affine_variety_points} If $X$ is an affine scheme of finite
    type over $K$, and if $R$ is a $K$-algebra which is also a topological ring, then we define a topology on the $R$-points $X(R)$ of $K$ by embedding the $K$-algebra homomorphisms from $A$ to $R$ into the set-theoretic maps from $A$ to $R$ with its product topology, and giving it the subspace topology.
\end{definition}

\begin{theorem}\label{topology_on_affine_variety_computation}\uses{topology_on_affine_variety_points}
    If $X$ is as above and $X\to\mathbb{A}^n_K$ is a closed immersion, then the induced map from $X(R)$ with its topology as above to $R^n$ is an embedding of topological spaces (that is, a homeomorphism onto its image).
\end{theorem}
\begin{proof} See \href{https://math.stanford.edu/~conrad/papers/adelictop.pdf}{Conrad's notes}.
\end{proof}

We now specialise to the smooth case. I want to make the following conjectural ``definition'':

\begin{definition}\label{manifold_on_algebraic_variety_points}\notready Let $K$ be a field equipped with an isomorphism to the reals, complexes, or a finite extension of the $p$-adic numbers. Let $X$ be a smooth affine algebraic variety over $K$. Then the points $X(K)$ naturally inherit the structure of a manifold over $K$.
\end{definition}

\begin{remark} Probably this is fine for a broader class of fields $K$.
\end{remark}

\begin{theorem}\label{manifold_on_algebraic_variety_computation}\uses{manifold_on_algebraic_variety_points}\notready
    If $X$ is as in the previous definition and $X\to\mathbb{A}^n_K$ is a closed immersion, then the induced map from $X(K)$ with its manifold structure to $K^n$ is an embedding of manifolds.
\end{theorem}
\begin{proof} I'm assuming this is standard, if true.
\end{proof}
\begin{corollary}\label{lie_group_from_algebraic_group}\uses{manifold_on_algebraic_variety_computation}\notready
    If $G$ is an affine algebraic group of finite type over $K=\R$ or $\bbC$ then $G(K)$ is naturally a real or complex Lie group.
\end{corollary}
\begin{remark}

    The corollary, for sure, is true! And it's all we need. I have not yet made any serious effort to find a reference for the definition or independence, although there seem to be some ideas \href{https://leanprover.zulipchat.com/#narrow/stream/116395-maths/topic/top.20space.20.2F.20manifold.20structure.20on.20points.20of.20alg.20varieties/near/431812525}{here}. As a toy example, one can embed $\GL_n(\R)$ into either $\R^{n^2+1}$ via $M\mapsto (M,\det(M)^{-1})$ or into $\R^{2n^2}$ via $M\mapsto (M,M^{-1})$ and the claim is that the two induced manifold structures are the same.
\end{remark}
\section{Algebraic groups.}

The concept of an affine algebraic group over a field $K$ can be implemented in Lean as a commutative Hopf algebra over $K$, as a group object in the category of affine schemes over $K$, as a representable group functor on the category of affine schemes over $K$, or as a representable group functor on the category of schemes over $K$ which is represented by an affine scheme. All of these are the same to mathematicians but different to Lean and some thought should go into which of these should be the actual definition, and which should be proved to be the same thing as the definition.

\begin{definition}\label{connected_reductive_group}\notready An affine algebraic group~$G$ of finite type over a field~$k$ is said to be \emph{connected} if it is connected as a scheme, and \emph{reductive} if $G_{\overline{k}}$ has no nontrivial smooth connected unipotent normal $k$-subgroup.
\end{definition}

\section{Automorphic forms and representations}

This section needs a lot of work; I am just attempting to write down some approximation to the (well-known) definitions but in great generality (far greater than we need). Some definitions below are short on details; indeed there may even be errors or imprecisions right now (because we are working in more generality than I am used to). It will be a very interesting project to get these details down. One reference (which leaves a lot of exercises) is Borel-Casselman in \cite{corvallis1}. Even \emph{stating} these definitions will be a big challenge in Lean; indeed one of the motivations of the project is that it forces us to write down all the below properly.

Let $G$ be a connected reductive group over a number field $N$.  We note
that $G(\A_N^f)$ is a (locally profinite) topological space and $G(N_\infty)$ is a real Lie group;
their product is $G(\A_N)$. If $g\in G(\A_N)$, write $g_f\in G(\A_N^f)$ for the finite part and $g_\infty\in G(N_\infty)$ for its infinite part.

For some reason, in the literature people seem to fix a choice of maximal compact subgroup $U_\infty$ of $G(N_\infty)$. I believe that all such subgroups are conjugate, and probably this gives some route between the different definitions coming from the different choices.

Example: if $G=\GL_2$ and $N=\Q$ then $N_\infty=\R$ and $G(N_\infty)$ is just $\GL_2(\R)$ with its usual Lie group structure and we can take $U_\infty$ to be $O_2(\R)$; $G(\A_N^f)$ is the restricted product of $\GL_2(\Q_p)$ over $\GL_2(\Z_p)$, for all primes $p$.

By assumption, $G(N_\infty)$ admits a finite-dimensional (algebraic) representation $\rho$ with finite kernel. Consider $\rho$ as taking values in $GL_N(\bbC)=\Aut_{\bbC}(V)$. Fix a hermitian sesquilinear form on $V$ which
is $U_\infty$ invariant, and now define a norm $||g||_\rho$ on $G(N_\infty)$ by $$||g||_\rho=(\tr \rho(g)^*\rho(g))^{1/2},$$ where the asterisk denotes adjoint with respect to the sesquilinear form. According to the article by Borel--Jacquet in \cite{corvallis1} (p189), if $\rho'$ is another such choice then there exists a positive real $C$ and a positive integer $n$ such that $||g||_{\rho'}\leq C||g||_\rho^n$ for all $g\in G(N_\infty)$.

\begin{definition}\label{slowly_increasing}\notready A function $f : G(N_\infty)\to\bbC$ is \emph{slowly-increasing} if there exists some $C>0$
    and $n\geq1$ such that $|f(x)\leq C||x||_\rho^n$.
\end{definition}

\begin{theorem}\label{slowly_increasing_well_defined}\uses{slowly_increasing}\notready This is independent of the choice of $\rho$ as above.
\end{theorem}
\begin{proof} Follows from the above.
\end{proof}

We can now give the definition of an automorphic form. For FLT we only need the definition for $G$ being either an abelian algebraic group, or an inner form of $GL(2)$, but we have chosen to work in full generality here.

\begin{definition}\label{automorphic_form}\uses{slowly_increasing_well_defined,connected_reductive_group, lie_group_from_algebraic_group, topology_on_affine_variety_computation}\notready An \emph{automorphic form} is a function $\phi:G(\A_N)\to\bbC$ satisfying the following conditions:
    \begin{itemize}
        \item $\phi$ is locally constant on $G(\A_N^f)$ and $C^\infty$ on $G(N_\infty)$. In other words, for every $g_\infty$, $\phi(-,g_\infty)$ is locally constant, and for every $g_f$, $\phi(g_f,-)$ is smooth.
        \item $\phi$ is left-invariant under $G(N)$;
        \item $\phi$ is right-$U_\infty$-finite (that is, the space spanned by $x\mapsto \phi(xu)$ as $u$ varies over $U_\infty$ is finite-dimensional);
        \item $\phi$ is right $K_f$-finite, where $K_f$ is one (or equivalently all) compact open subgroups of $G(\A_N^f)$;
        \item $\phi$ is $\mathcal{z}$-finite, where $\mathcal{z}$ is the centre of the universal enveloping algebra of the Lie algebra of $G(N_\infty)$, acting via differential operators. Equivalently $\phi$ is annihiliated by a finite index ideal of this centre, so morally $\phi$ satisfies lots of differential equations of a certain type;
        \item For all $g_f$, the function $g_\infty\mapsto \phi(g_f g\infty)$ is slowly-increasing in the sense above.
    \end{itemize}
\end{definition}

Automorphic forms form a typically infinite-dimensional vector space.

\begin{definition}\label{cuspidal_automorphic_form}\uses{automorphic_form}\notready An automorphic form is \emph{cuspidal} (or ``a cusp form'') if it furthermore satisfies $\int_{U(N)\backslash U(\A_N)}\phi(ux)du=0$, where $P$ runs through all the proper parabolic subgroups of $G$ defined over $N$ and $U$ is the unipotent radical of $P$, and the integral is with respect to the measure coming from Haar measure.
\end{definition}

The cuspidal automorphic forms form a complex subspace of the space of automorphic forms.

\begin{definition}\label{automorphic_form_actions}\uses{automorphic_form}\notready The group $G(\A_N)$ acts on itself on the right, and this induces a left action of its
    subgroup $G(\A_N^f)\times U_\infty$ on the spaces of automorphic forms and cusp forms. The Lie algebra $\mathfrak{g}$ of $G(N_\infty)$ also acts, via differential operators. Furthermore the actions of $\mathfrak{g}$ and $U_\infty$ are compatible in the sense that the differential of the $U_\infty$ action is the action of its Lie algebra considered as a subalgebra of $\mathfrak{g}$. We say that the spaces are $(G(\A_N^f)\times U_\infty,\mathfrak{g})$-modules.
\end{definition}

\begin{theorem}\label{cuspidal_automorphic_form_decomposition}\uses{cuspidal_automorphic_form, automorphic_form_actions} The cusp forms decompose as a (typically infinite) direct sum of irreducible $(G(\A_N^f)\times U_\infty,\mathfrak{g})$-modules.
\end{theorem}
\begin{definition}\label{cuspidal_automorphic_representation}\uses{cuspidal_automorphic_form_decomposition}\notready A cuspidal automorphic representation is an irreducible $(G(\A_N^f)\times U_\infty,\mathfrak{g})$-module isomorphic to an irreducible summand of the space of cusp forms.
\end{definition}

For non-cuspidal representations, they do not decompose as a direct sum; there is a continuous spectrum which decomposes as a direct integral. We may not ever need these. As a result the definition of an automorphic representation has to be slightly modified in the non-cuspidal case.

\begin{definition}\label{automorphic_representation}\uses{automorphic_form_actions}\notready An automorphic representation is an irreducible $(G(\A_N^f)\times U_\infty,\mathfrak{g})$-module isomorphic to an irreducible subquotient of the space of automorphic forms.
\end{definition}

Admissibility is a finiteness condition on an irreducible representation of $(G(\A_N^f)\times U_\infty,\mathfrak{g})$; automorphic representations are admissible, and this seems to boil down to theorems of Godement and Harish-Chandra in the general case.


\begin{theorem}\label{automorphic_representation_local_decomposition}\uses{automorphic_representation}\notready An irreducible admissible $(G(\A_N^f)\times U_\infty,\mathfrak{g})$-module is a restricted tensor product of irreducible representations $\pi_v$ of $G(N_v)$ as $v$ runs through the finite places of $N$, tensored with a tensor product of irreducible $(\mathfrak{g}_v,U_{\infty,v})$-modules as $v$ runs through the infinite places of $N$. The representations $\pi_v$ are unramified for all but finitely many $v$.
\end{theorem}
\begin{proof} See Flath's article in~\cite{corvallis2}.
\end{proof}

As mentioned above, we only need all of this for abelian algebraic groups and for inner forms of $\GL_2$ over totally real fields, where everything can be made more concrete (and in particular where I can write down concrete definitions, although this still needs to be done). In particular, we don't strictly speaking need all of the above, we could just cheat and deal with $\GL_2(\R)$ and $\bbH^\times$ separately.

The theorems I need are: Jacquet-Langlands for inner forms of $\GL_2$ over totally real fields,
and multiplicity 1 for these inner forms. We also need cyclic base change plus classification of image, all for totally definite quaternion algebras, and we need
automorphic induction from $\GL_1(K)$ to $\GL_2(F)$ when $K/F$ is a degree 2 totally imaginary extension. There seems to be little point formalising the statements of the theorems if we cannot yet even formalise the definition of an automorphic representation properly.

\section{Galois representations}

Ivan Farabella has formalised the definition of a compatible family of Galois
representations, modulo the existence of Frobenius elements, which has been
established by Jou Glasheen.

\begin{definition}\label{compatible_family}\discussion{23} Let $N$ be a number field. A \emph{compatible family of $d$-dimensional Galois representations over $N$} is a finite set of finite places $S$ of $N$,
a number field $E$, a monic degree $d$ polynomial $F_{\p}(X)\in E[X]$ for each finite place $\p$ of $K$ not in $S$ and, for each prime number $\ell$ and field embedding $\phi : E\to\Qlbar$ (or essentially equivalently for each finite place of $E$), a continuous homomorphism $\rho:\GK\to\GL_2(\Qlbar)$ unramified outside $S$ and the primes of $K$ above $\ell$, such that $\rho(\Frob_\p)$ has characteristic polynomial $P_\pi(X)$ if $\pi$ lies above a prime number $p\not=\ell$ with $p\not\in S$.
\end{definition}

The big theorem, which again we are far from even \emph{stating} right now, is

\begin{theorem}\label{Galois_representation_from_automorphic_representation_on_GL_2_form}\uses{automorphic_representation,Shimura_varieties,compatible_family}\notready Given an automorphic representation $\pi$ for an inner form of $\GL_2$ over a totally real field and with reflex field~$E$, such that $\pi$ is weight 2 discrete series at every infinite place, there exists a compatible family of 2-dimensional Galois representations associated to $\pi$, with $S$ being the places at which $\pi$ is ramified, and $F_{\p}(X)$ being the monic polynomial with roots the two Satake parameters for $\pi$ at $\p$.
\end{theorem}

\section{Algebraic geometry}

We have already mentioned Mazur's Theorem on torsion subgroups of elliptic curves (theorem~\ref{mazur}).
The proof of this is the main theorem of~\cite{mazur-torsion}, 150 pages of subtle arithmetic
geometry involving the bad reduction of modular curves, exotic cohomology theories (etale and more),
and the consequences of this for the Neron models of their Jacobians. After a beautiful
introductory chapter containing a history and examples, the convention is established that
throughout the paper, $N$ will denote a prime number which is at least 5. And then the first
sentence of chapter 1 of the paper proper is ``Consider quasi-finite separated commutative group
schemes of finite presentation over the base $S:=\Spec{\Z}$ which are finite flat group schemes
over $S':=\Spec(Z[1/N]).$''. At the time of writing (May 2024), Lean's algebraic geometry cannot
get us through \emph{the first sentence of Mazur's proof}, which occupies pages 43 to 172 of
the paper (not including the appendix or references, that's just the proof). Anyone interested
in formalising Mazur's paper should make a formalisation of its first sentence their first
milestone.

Talking of modular curves, we also need the existence of Shimura curves and surfaces over totally real fields~$F$ (of degree greater than~2, so always compact). The curves are "modeles \'etranges" in the sense of Deligne, so we also need moduli spaces of unitary Shimura varieties over CM extensions. We need to decompose the first and second etale cohomology groups of these varieties into Galois representations, by understanding them
in terms of automorphic representations.

\begin{definition}\label{Shimura_varieties}\notready We need the definition of (the canonical model over $F$ of) the Shimura curve attached to an inner form of $\GL_2$ with precisely one split infinite place, and the same for the Shimura surface associated to an inner form split at two infinite places (and ramified elsewhere, so it's compact).
\end{definition}

We also need Moret-Bailly's theorem from~\cite{moret-bailly}:

\begin{theorem}\label{moret-bailly}\notready Let $K^{\avoid}/K$ be a Galois extension of number fields. Suppose also
that $S$ is a finite set of places of $K$. For $v\in S$ let $L_v/K_v$ be a finite Galois extension.
Suppose also that $T /K$ is a smooth, geometrically connected curve and that for each
$v\in S$ we are given a nonempty, $\Gal(L_v/K_v)$-invariant, open subset $\Omega_v\subseteq (L_v)$.
Then there is a finite Galois extension $L/K$ and a point $P ∈ T (L)$ such that
\begin{itemize}
\item $L/K$ is Galois and linearly disjoint from $K^{\avoid}$ over $K$;
\item if $v\in S$ and $w$ is a prime of $L$ above $v$ then $L_w /K_v$ is isomorphic to $L_v/K_v$;
\item and $P \in\Omega_v\subseteq T (L_v) \cong (L_w)$ via one such $K_v$-algebra morphism
(this makes sense as $\Omega_v$ is $\Gal(L_v/K v)$-invariant).
\end{itemize}
\end{theorem}

Note that we do not even have the definition of a curve over a field in Lean.

\section{Algebra}

We need the classification of finite subgroups of $\PGL_2(\overline{\F}_p)$. The answer is that they are all cyclic, dihedral, $A_4$, $S_4$, $A_5$, or isomorphic to $\PSL_2(k)$ or $\PGL_2(k)$ for some finite field of characteristic~$p$. This should at least be easy to state!


%%FILE blueprint/src/chapter/global_langlands.tex
This is not imported right now.

\chapter{The Global Langlands conjectures}

\section{Overview of the chapter}

In this section we discuss the problem of attempting to state the motivic global Langlands conjectures for connected reductive groups over number fields. More precisely, the goal is to formally state the Buzzard--Gee version of the conjecture, which applies to algebraic automorphic representations. The main difficulties here are in writing down precisely what is meant by phrases such as ``automorphic representation'', ``connected reductive group'', ``local-global compatibility'', ``de Rham Galois representation'', and so on; these words hide a large amount of technical material.

\section{Statement of the conjecture}

Let $R$ be a commutative base ring; it will often be a field but we shall develop the theory in more generality when there is no extra effort needed to do.

\begin{definition}\label{affine_group_scheme_over_affine_base}\lean{???}

  An \emph{affine group scheme over $R$} is a group object in the category of affine schemes over $R$.

  \begin{definition}\label{Hopf_algebra}\lean{TODO}% we have Hopf algebras in mathlib

    %%%%%%%%%%%%%%%%%%%

    TODO: Connected and reductive


%%FILE blueprint/src/content.tex
% Add chapters of the blueprint here

\input{chapter/ch01introduction}
\input{chapter/ch02reductions}
\input{chapter/ch03freyreduction}
%\input{chapter/ch03freyold} -- replaced by the chapter above
\input{chapter/ch04overview}
\input{chapter/ch05automorphicformexample}
\input{chapter/ch06automorphicrepresentations}
\input{chapter/ch07exampleGLn}
\input{chapter/FrobeniusProject}
\input{chapter/AdeleMiniproject}
\input{chapter/HaarCharacterProject}
\input{chapter/FujisakiProject}
\input{chapter/QuaternionAlgebraProject}
\input{chapter/HeckeOperatorProject.tex}
\input{chapter/chtopbestiary}
\input{chapter/biblio}

